% !TeX root = ../../Book3.tex
% 纲要标题：## 第5章　完整群、平行结构与黎曼分解
% 纲要位置：几何学计划纲要.md，第 2260 行。
% 正文状态：初稿；尚未正式审核。

% 纲要内容开始
% 
% > 本章目标与证明范围：从平行移动得到完整群，研究其如何决定平行张量和局部分解。de Rham 的整体版本保留单连通及完备性；特殊完整群分类单列引用而不倒推基础结果。
% 
% 纲要内容结束

\chapter{完整群、平行结构与黎曼分解}
\label{g3:ch:05}
\BookChapterNumber{g3:ch:05}

\begin{chapterabstract}
	沿闭路平行运输的返回变换组成完整群。其连续部分由沿路径搬回基点的曲率算子决定，整体分支则还记录环路的拓扑。本章建立完整群的浸入 Lie 群结构，证明 Ambrose--Singer 定理，再把不变张量和不变子空间分别转化为平行张量和正交平行分布。局部乘积在单连通与完备条件下延拓为整体 de Rham 分解。随后从不变度量、点反射与对称 Lie 代数构造齐性及对称空间，计算球面、双曲空间和特殊线性群商的曲率、秩与限制根。特殊完整群及对称空间的分类给出明确的适用条件和外部来源。
\end{chapterabstract}

\begin{learninggoals}
	\begin{itemize}
		\item 区分整体与受限完整群、浸入子群与闭子群，并用小矩形及路径变分恢复完整代数。
		\item 由完整群不变量构造平行张量和分布，证明局部乘积并追踪它的整体延拓条件。
		\item 使用 de Rham 分解识别 Euclid 因子、不可约因子及商空间中的离散扭曲。
		\item 从对称分解计算模型曲率、最大平坦与秩，核验紧型和非紧型对偶。
		\item 通过平行形式理解特殊完整群，并计算 \(\operatorname{SL}(n,\mathbb R)/\operatorname{SO}(n)\) 的限制根及重数，辨明分类表中的局部与整体信息。
	\end{itemize}
\end{learninggoals}

一个向量沿闭路回到起点以后，长度没有改变，方向却可能已经改变。球面上的小环路产生与曲率相联系的转动；平坦流形则可能在绕过一个不可缩环路后翻转某些方向。例如，在 Euclid 平面上用平移和滑移反射构造平坦 Klein 瓶\footnote{Felix Christian Klein（1849--1925），德国数学家，以变换群组织几何的 Erlangen 纲领及函数论研究著称。}，局部的联络系数与曲率都可以为零，但穿过滑移反射所对应环路的运输仍会改变一个坐标方向。要说明哪些向量、张量或方向分解能够在整个流形上保持平行，就必须同时记录这两种返回现象。

\ProseParagraphBreak

记录全部返回变换，得到基点切空间上的完整群。把小环路连续缩小，会把群的问题变成 Lie 代数的问题；曲率恰好给出面积阶的首个变化。一个点的曲率值还不够，因为环路可以先走到别处，再绕小圈，最后返回。因此需要把各点的曲率沿路径搬回同一个切空间。这将导出完整代数的生成定理，也解释为什么平行张量可以从一个切空间中的不变量恢复。

\ProseParagraphBreak

不变子空间进一步给出平行分布。Levi-Civita 联络的无挠性使这些分布可积，正交性则把局部度量分成乘积。由局部坐标盒走到整体乘积，需要防止路径延拓在有限长度处中断，并消除沿不同路径继续时的单值障碍；完备性和单连通性正是在这两处发挥作用。以后研究曲率比较和几何方程时，这一分解会把问题归结到不可约因子。

\ProseParagraphBreak

本章以第一章的 Levi-Civita 联络、第二章的完备性及 Hopf--Rinow 定理、第三章的曲率恒等式为直接基础。第一卷第13.4节的 Frobenius 定理、第16.3节的 Lie 子代数积分、第17章的齐性商和第18章的主丛联络在使用处恢复所需结论；涉及整体提升时就近说明覆盖的构造。主干的完整群与分解论证在本章展开。对称空间的标准齐性结构以及穷尽分类分别按 Loos 第一卷第II章、第IV章和第二卷第VII章引用\cite[第II章，第1--2节；第IV章，第1节]{Loos1969SymmetricI}\cite[第VII章，第2--3节]{Loos1969SymmetricII}；特殊完整群的名单及例外群模型依据 Joyce 第2节\cite[第2.1--2.5节，第111--115页]{Joyce2006ExceptionalHolonomy}。

% !TeX root = ../../Book3.tex
% 纲要标题：### 5.1 完整群的 Lie 群结构与曲率
% 纲要位置：几何学计划纲要.md，第 2264 行。
% 正文状态：初稿；尚未正式审核。

% 纲要内容开始
% 
% * 平行运输的路径群
% * 受限完整群和整体完整群
% * 局部平行截面及浸入 Lie 群结构
% * 小矩形运输的二阶展开
% * 沿路径搬运的曲率算子
% * Ambrose–Singer 定理
% 
% 纲要内容结束

\section{完整群的 Lie 群结构与曲率}
\label{g3:sec:0501}

沿闭曲线运输一个向量，回到起点时，它一般已经转向。把所有可能的返回变换放在一起，就得到完整群。曲率则测量一个很小的闭曲线所造成的首个非零变化。本节要把这两个尺度联系起来：先证明返回变换确实组成一个有限维 Lie 群，再证明它的 Lie 代数由沿路径搬回基点的曲率算子张成。

\ProseParagraphBreak

证明适用于比切丛稍广的对象。以下设 \(M\) 为连通、有限维、Hausdorff、第二可数且无边界的光滑流形，\(E\to M\) 为实秩 \(m\geq1\) 的光滑向量丛，带正定纤维内积和度量联络 \(\nabla^E\)。其曲率记为
\[
	R^E(X,Y)=\nabla^E_X\nabla^E_Y-\nabla^E_Y\nabla^E_X-\nabla^E_{[X,Y]}.
\]
它是纤维上的斜自伴算子。取 \(E=\mathrm TM\)、\(\nabla^E=\nabla\) 时，就得到黎曼度量的完整群；此时 \(R^E=R\)，符号与第三章一致。以下不需要底流形完备。

\subsection{平行运输的路径群}
\label{g3:subsec:0501:01}

一条路径总指紧区间上的分片光滑路径。若 \(\gamma\) 从 \(p\) 到 \(q\)，记其平行运输为
\[
	P_\gamma:E_p\longrightarrow E_q.
\]
第一卷第15.2节的“平行移动”定理及运算规律已经证明：\(P_\gamma\) 存在、唯一，对具有共同有限分段的光滑路径族光滑依赖；若 \(\gamma*\eta\) 表示先走 \(\gamma\)、再走 \(\eta\)，而 \(\bar\gamma\) 表示反向路径，则
\begin{equation}\label{g3:eq:holonomy-transport-laws}
	P_{\gamma*\eta}=P_\eta P_\gamma,
	\qquad P_{\bar\gamma}=P_\gamma^{-1}.
\end{equation}
这些结论来自局部线性方程 \(u'=-A(\dot\gamma)u\) 的存在唯一性和参数依赖，并允许插入常值路段。度量相容性又给出
\[
	\frac{\mathrm d}{\mathrm dt}\langle P_{\gamma|_{[0,t]}}v,P_{\gamma|_{[0,t]}}w\rangle_E=0,
\]
故每个 \(P_\gamma\) 都是线性等距。

\ProseParagraphBreak

固定基点 \(p\)。基点环路的拼接在参数化层面并不严格结合，但把保向重参数化、插入常值路段和插入或删除一段“走出后立即原路返回”视为等价后，就得到一个群：取乘法次序为先右因子、后左因子，常值环路为单位元，反向环路为逆。这里等价关系也对前后拼接封闭，所以群运算良定义。我们把这个商称为这里使用的路径群；平行运输由\eqref{g3:eq:holonomy-transport-laws}给出它到 \(\mathrm O(E_p)\) 的群同态。它保留路径的几何信息，比只保留同伦类的基本群更细。

\begin{definition}\label{g3:def:holonomy-group}
	在上述度量向量丛上，基点 \(p\) 处的\term[wanzhengqun]{完整群}[holonomy group]为
	\begin{equation}\label{g3:eq:holonomy-group}
		\operatorname{Hol}_p(E,\nabla^E)
		=\{P_\gamma:\gamma\text{ 是以 }p\text{ 为基点的分片光滑环路}\}
		\subseteq\mathrm O(E_p).
	\end{equation}
	简记为 \(\operatorname{Hol}_p\)。在 \(E=\mathrm TM\) 的情形，也写作 \(\operatorname{Hol}_p(M,g)\)。
\end{definition}
\symidx{\(\operatorname{Hol}_p\)}

第一卷主丛语言中的“环路平行移动群”与这里的完整群描述同一返回现象。在选定 \(E_p\) 的一个正交标架后，上式成为 \(\mathrm O(m)\) 的子群；换标架只把它共轭。若 \(\alpha:p\to q\)，则绕行 \(\bar\alpha*\gamma*\alpha\) 给出
\begin{equation}\label{g3:eq:holonomy-basepoint-change}
	\operatorname{Hol}_q=P_\alpha\operatorname{Hol}_pP_\alpha^{-1}.
\end{equation}
这说明连通流形上完整群的共轭类型与基点无关，但具体算子仍作用于指定的纤维。

\subsection{受限完整群和整体完整群}
\label{g3:subsec:0501:02}

\begin{definition}\label{g3:def:restricted-holonomy}
	\term[shouxianwanzhengqun]{受限完整群}[restricted holonomy group]定义为
	\[
		\operatorname{Hol}_p^0
		=\{P_\gamma:\gamma\text{ 在保持基点的连续同伦下可缩为常值环路}\}.
	\]
	这里 \(\gamma\) 本身仍要求分片光滑。相较之下，允许全部基点环路所得的 \(\operatorname{Hol}_p\) 也称整体完整群。
\end{definition}
\symidx{\(\operatorname{Hol}_p^0\)}

上标 \(0\) 暂时只是记号；在建立适当的 Lie 群拓扑后，它才会被证明恰好是单位元连通分支。不能在此先把它解释为 \(\mathrm O(E_p)\) 的子空间拓扑下的连通分支。

\begin{proposition}\label{g3:prop:restricted-holonomy-normal}
	对连通光滑流形上的度量向量丛，\(\operatorname{Hol}_p^0\) 是 \(\operatorname{Hol}_p\) 的正规子群，且存在满同态
	\[
		\pi_1(M,p)\longrightarrow\operatorname{Hol}_p/\operatorname{Hol}_p^0.
	\]
	特别地，\(M\) 单连通时，两种完整群相同。
\end{proposition}
\begin{proof}
	可缩环路的拼接与反向仍可缩；把可缩环路前后接上互为反向的固定路径，也仍可缩。因此\eqref{g3:eq:holonomy-transport-laws}证明子群性和正规性。若 \(\gamma,\eta\) 定端点同伦，则 \(\gamma*\bar\eta\) 可缩，所以 \(P_\eta^{-1}P_\gamma\in\operatorname{Hol}_p^0\)。因而 \(P_\gamma\) 的陪集只依赖 \([\gamma]\)。若基本群乘法采用先左后右的路径次序，就取 \([\gamma]\mapsto P_\gamma^{-1}\operatorname{Hol}_p^0\)，它由\eqref{g3:eq:holonomy-transport-laws}成为同态。

	每条连续环路都可分成有限段，每段落在一个坐标球中，再以坐标线段代替；原段与新段在该球中定端点同伦。这给出同伦类的分片光滑代表，也说明以上映射定义在整个基本群上。满射性来自完整群的定义。
\end{proof}

命题说的是运输在商群中的类具有同伦不变性。一般曲率不为零时，同伦环路的实际运输可以不同。例如，一个小可缩环路本来就可能转动向量。

\begin{example}\label{g3:exm:dense-flat-holonomy}
	在 \(S^1=\mathbb R/(2\pi\mathbb Z)\) 上取平凡的秩二正交丛，令
	\[
		J=\begin{pmatrix}0&-1\\1&0\end{pmatrix},\qquad
		\nabla^E=\mathrm d+\alpha J\,\mathrm d\theta,
	\]
	其中 \(\mathrm d\theta\) 是由覆盖实线上的角坐标微分下降的一形式。沿绕数为 \(k\) 的环路，平行方程直接给出
	\[
		P_\gamma=\exp(-2\pi k\alpha J).
	\]
	因此 \(\operatorname{Hol}_p^0=\{I\}\)，而整体完整群由一次绕行的旋转生成。若 \(\alpha\) 无理，这个循环子群在 \(\mathrm{SO}(2)\) 中稠密。其稠密性可由抽屉原理看出：\(0,\alpha,\ldots,N\alpha\) 的小数部分中有两项距离小于 \(1/N\)，故子群含有任意小的非零旋转；这类旋转的有限次幂能以任意精度逼近圆上的每一点。后面构造的 Lie 群拓扑使这个完整群成为可数离散群。
\end{example}

\subsection{局部平行截面及浸入 Lie 群结构}
\label{g3:subsec:0501:03}

在一个以 \(p\) 为坐标原点的星形坐标域内，把 \(E_p\) 的基沿坐标径向线运输，就得到光滑局部标架。这个标架沿所选径向线平行，却未必沿其他方向平行。若要求一个标架的全部截面在所有方向都平行，条件就变为曲率为零：必要性由曲率定义立刻得到，充分性是\cref{g3:lem:flat-connection-local-frame}。完整群研究的正是沿不同路径构造这些局部截面时，它们之间可能出现的转动。

\ProseParagraphBreak

下面把“改变环路所产生的方向”变成有限维空间。称 \(s\mapsto\gamma_s\) 为光滑环路族，是指参数在一个开区间内变化，\(\gamma_s(0)=\gamma_s(1)=p\)，而 \([0,1]\) 有一个不依赖 \(s\) 的有限分段，使 \((s,t)\mapsto\gamma_s(t)\) 在每一段上光滑并延至端点。允许的族也包括有限次拼接、反向及固定路径的前后连接。第一卷的平行移动定理保证
\[
	a(s)=P_{\gamma_s}
\]
是 \(\mathrm O(E_p)\) 中的光滑曲线。对多参数族采用相同定义。

\begin{lemma}\label{g3:lem:contractible-loop-generators}
	设 \(M\) 为连通光滑流形。任一可缩分片光滑基点环路的运输，可以写成有限个因子的乘积，每个因子来自一个形如 \(\alpha*\beta*\bar\alpha\) 的环路：\(\alpha\) 从 \(p\) 到某个 \(q\)，\(\beta\) 是以 \(q\) 为基点且完全位于一个凸坐标球中的环路。每个这样的因子都可以通过一条光滑环路族的运输从 \(I\) 连到它。
\end{lemma}
\begin{proof}
	把可缩同伦写在一个闭正方形上，使一边是所给环路，其余三边恒为 \(p\)。由紧性，取足够细的有限矩形网格，使每个小格的像落在一个凸坐标球中，并把边界环路的分段点纳入网格。保留外边界路径；对每条内边，选择连接其两个端点的分片光滑路径，且使它留在所有相邻小格对应坐标球的交中。这样的选择存在：原来的连续边已经落在这个开交集内，用有限个小坐标球逐段代替即可。

	现在每个小格的边界都是某一凸坐标球内的分片光滑环路。逐格扫过正方形，把已扫区域的边界路径换成加入一格后的边界路径；每一步的差是用一条网格路径从 \(p\) 引到该格顶点后，绕其边界再返回的环路。内部公共边以相反方向出现，运输逐对消去。有限次操作后，外边界的运输便写成这些带引路的小格环路运输的乘积。

	设 \(x\) 是其中一个坐标球的坐标，\(\beta(0)=\beta(1)=q\)。用
	\[
		x(\beta_s(t))=x(q)+s\bigl(x(\beta(t))-x(q)\bigr),\qquad0\leq s\leq1,
	\]
	收缩 \(\beta\)。凸性保证其像留在坐标球中，分片光滑性及共同分段也保持。由于这族路径的像是球内的紧集，参数还可稍微延伸到 \([0,1]\) 的两侧。固定引路 \(\alpha\) 后，所得环路族在 \(s=0\) 的运输为 \(I\)，在 \(s=1\) 为所需因子。
\end{proof}

\begin{theorem}[完整群的 Lie 群结构]\label{g3:thm:holonomy-lie-group}
	设 \(M\) 为连通、有限维、Hausdorff、第二可数且无边界的光滑流形，\((E,\langle\ ,\ \rangle_E,\nabla^E)\) 为有限秩度量向量丛。则 \(\operatorname{Hol}_p^0\) 有自然的连通浸入 Lie 子群结构；\(\operatorname{Hol}_p\) 有使其各个 \(\operatorname{Hol}_p^0\) 陪集为开集的 Lie 群结构，并且
	\[
		\bigl(\operatorname{Hol}_p\bigr)_{\mathrm{id}}=\operatorname{Hol}_p^0,
		\qquad
		\operatorname{Lie}(\operatorname{Hol}_p)=\operatorname{Lie}(\operatorname{Hol}_p^0)
		\subseteq\mathfrak{so}(E_p).
	\]
	包含映射到 \(\mathrm O(E_p)\) 是光滑浸入。光滑环路族的运输对这里的 Lie 群结构光滑。
\end{theorem}
\begin{proof}
	记 \(H=\operatorname{Hol}_p\)、\(H_0=\operatorname{Hol}_p^0\)，并在 \(\mathfrak{so}(E_p)\) 中定义
	\[
		\mathfrak k=\operatorname{span}_{\mathbb R}
		\{a'(0):a(s)=P_{\gamma_s},\ a(0)=I\}.
	\]
	有限个这样的导数之和仍可由一个环路族实现：对相应运输曲线作乘积，并把各参数乘以所需实系数，乘积在零点的导数就是该线性组合。因此 \(\mathfrak k\) 中每个向量都有这种实现。若 \(a(0)=I\) 而 \(\gamma_0\) 本身不一定可缩，可把族改成 \(\gamma_s*\bar\gamma_0\)；其运输仍为 \(a(s)\)，并且每个环路可缩。这说明上述实现曲线均可取值于 \(H_0\)。

	对任意 \(h\in H\)，用一个固定环路及其反向共轭整族路径，得到 \(h\mathfrak k h^{-1}=\mathfrak k\)。若 \(A,B\in\mathfrak k\)，取 \(a(0)=I\)、\(a'(0)=A\)，则 \(a(s)Ba(s)^{-1}\) 始终在 \(\mathfrak k\) 中。求导得
	\[
		[A,B]=\left.\frac{\mathrm d}{\mathrm ds}\right|_{s=0}a(s)Ba(s)^{-1}\in\mathfrak k.
	\]
	这里用到的闭性仅是有限维线性子空间 \(\mathfrak k\) 的闭性。于是 \(\mathfrak k\) 是 Lie 子代数。

	第一卷第16.3节“Lie 子代数的积分”定理给出唯一的连通浸入 Lie 子群 \(K\subseteq\mathrm O(E_p)\)，其 Lie 代数为 \(\mathfrak k\)。具体地，\(K\) 是左不变分布
	\[
		D_g=\{gA:A\in\mathfrak k\}
	\]
	经过 \(I\) 的最大积分叶，带叶自身的拓扑。我们还要用该节“取值于叶的光滑映射”引理：一个连通参数域上的环境光滑映射，若微分处处取值于 \(D\)，其像就在一片叶内，并且对叶的流形结构光滑。这个事实在叶状坐标 \((u,v)\) 中很直接：\(\mathrm d(v\circ f)=0\) 使横向坐标局部恒定，故 \(f\) 局部取值于一张开叶片，叶片坐标就是 \(u\circ f\)。

	对任意光滑环路族的运输 \(a(s)\)，固定 \(s_0\)，曲线 \(a(s_0)^{-1}a(s_0+t)\) 仍是一个环路族的运输且在 \(t=0\) 等于 \(I\)。所以
	\begin{equation}\label{g3:eq:holonomy-family-logarithmic-derivative}
		a(s)^{-1}a'(s)\in\mathfrak k.
	\end{equation}
	叶值映射引理说明：若 \(a(0)=I\)，则整条曲线在 \(K\) 中，且作为 \(K\) 中的曲线光滑。多参数族对每个参数方向都满足同一条件，也得到相同结论。结合\cref{g3:lem:contractible-loop-generators}，每个可缩环路的运输是这种曲线终值的有限乘积，故 \(H_0\subseteq K\)。

	反过来，选 \(\mathfrak k\) 的基 \(A_1,\ldots,A_d\)，用取值于 \(H_0\) 的环路族运输 \(a_i(s)\) 实现 \(a_i(0)=I\)、\(a_i'(0)=A_i\)。刚才的论证已保证它们对 \(K\) 的拓扑光滑。映射
	\[
		(t^1,\ldots,t^d)\longmapsto a_1(t^1)\cdots a_d(t^d)
	\]
	在原点的微分是到 \(\mathrm T_IK=\mathfrak k\) 的同构。逆函数定理使其像包含 \(K\) 中单位元的一个开邻域，而整个像在 \(H_0\) 内。因此 \(H_0\) 是 \(K\) 的开子群；它的其他陪集也开，连通性迫使 \(H_0=K\)。\(d=0\) 时同样得到 \(H_0=\{I\}\)。

	给 \(H\) 的每个左陪集 \(hK\) 赋予由左平移得到的 \(K\) 的流形结构，并把这些陪集取不交并。共轭映射 \(g\mapsto hgh^{-1}\) 保持 \(D\)，所以由叶值映射引理，它在 \(K\) 上是光滑自同构。在两个陪集上，群乘法可写成
	\[
		(ak)(b\ell)=ab\,(b^{-1}kb)\ell,
		\qquad (ak)^{-1}=a^{-1}(ak^{-1}a^{-1}),
	\]
	故乘法与取逆都光滑；包含映射在各陪集上也是浸入。这份拓扑 Hausdorff，因为各个陪集是 Hausdorff 开流形，互异陪集本来就不相交。

	还须验证第二可数性。连通第二可数流形的基本群可数，下面恢复所需理由。取可数个凸坐标球覆盖，在每球中选一点，在每两个球的非空交的每个连通分支中再选一点及通向两球所选点的路径。开集的连通分支至多可数，因为它们两两不交且每个都包含可数拓扑基中的一个元素；它们又由局部路径连通性而路径连通。任一环路可作有限分段，每段落在一个球中。利用交分支中的路径移动分段点，再利用每球的单连通性，把环路定端点同伦到有限条已选路径的拼接。加入基点所在球中的固定连接路径后，同伦类便由可数个字母的有限字表示，故至多可数。由\cref{g3:prop:restricted-holonomy-normal}，\(H/K\) 也至多可数。每个 \(K\) 陪集已有可数基，取它们的并就得到 \(H\) 的可数基。

	最后，\(K\) 连通而各陪集开且闭，所以 \(K\) 恰为单位元连通分支。任意光滑环路族局部可由固定运输左移到 \(K\)，前面的叶值映射论证也证明其内在光滑性。
\end{proof}

这个证明区分了两个问题：先在正交群中寻找允许的切方向，再证明对应积分叶中的每个点都由实际环路产生。后一步使用环路族给出的局部坐标，因此所得是完整群本身，而不是其闭包。

\begin{example}\label{g3:exm:nonclosed-restricted-holonomy}
	受限完整群也未必在环境正交群中闭。取 \(M=\mathbb R^2\)、\(E=M\times\mathbb R^4\)，在全局正交标架中令
	\[
		B=\begin{pmatrix}J&0\\0&\alpha J\end{pmatrix},\qquad
		\nabla^E=\mathrm d+Bx\,\mathrm dy,
		\qquad \alpha\notin\mathbb Q,
	\]
	其中 \(J\) 如\cref{g3:exm:dense-flat-holonomy}。所有联络矩阵彼此交换，故
	\[
		P_\gamma=\exp\left(-B\oint_\gamma x\,\mathrm dy\right),\qquad
		R^E=B\,\mathrm dx\wedge\mathrm dy.
	\]
	以原点为顶点的矩形及其反向绕行可使积分取任意实数，所以
	\[
		\operatorname{Hol}_0^0=\operatorname{Hol}_0=\{\exp(tB):t\in\mathbb R\}.
	\]
	环路族的积分 \(\oint x\,\mathrm dy\) 光滑依赖参数，且矩形面积可任意变化，因此上一定理构造的速度空间恰为 \(\mathbb RB\)。该参数化单射：\(t\in2\pi\mathbb Z\) 且 \(\alpha t\in2\pi\mathbb Z\) 只能同时在 \(t=0\) 成立。它的像在 \(\mathrm{SO}(2)\times\mathrm{SO}(2)\) 中稠密；固定第一角后取 \(t=\theta+2\pi k\)，第二角由无理旋转稠密性逼近任意给定角。因此完整群自身是一维 Lie 群，拓扑同 \(\mathbb R\)，闭包却是二维环面。这里给出的是度量向量丛联络的例子。
\end{example}

\subsection{小矩形运输的二阶展开}
\label{g3:subsec:0501:04}

曲率项的符号取决于绕行方向。我们固定先沿第一坐标正向、再沿第二坐标正向，随后沿两条反向边返回。与只令两条边同时缩短相比，把两边长度分别记为 \(s,t\) 能直接得到后面需要的一个参数导数。

\begin{proposition}\label{g3:prop:small-rectangle-holonomy}
	在度量向量丛的一个局部正交标架域内，取以 \(q\) 为原点的坐标。设 \(\dim M\geq2\)，令 \(L(s,t)\in\mathrm O(E_q)\) 为依次经过
	\[
		(0,0),\ (s,0),\ (s,t),\ (0,t),\ (0,0)
	\]
	的矩形边界运输，其余坐标为零。则 \(L\) 在 \((0,0)\) 附近光滑，且
	\begin{align}
		L(s,t)&=I-stR^E_q(\partial_1,\partial_2)
		+st\,O(|s|+|t|),\label{g3:eq:small-rectangle-holonomy}\\
		\partial_sL(0,t)&=-tR^E_q(\partial_1,\partial_2)+O(t^2).
		\label{g3:eq:small-rectangle-first-variation}
	\end{align}
	余项用纤维算子范数衡量；\(s,t\) 都允许正、负或零值。
\end{proposition}
\begin{proof}
	把四条边分别用固定的单位时间区间参数化，顶点对 \(s,t\) 光滑依赖，故平行方程的参数依赖给出 \(L\) 光滑。退化矩形只是往返路段，所以 \(L(s,0)=L(0,t)=I\)。逐分量两次应用微积分基本定理，得到
	\[
		L(s,t)-I=stC(s,t),\qquad
		C(s,t)=\int_0^1\!\int_0^1
		\partial_s\partial_tL(us,vt)\,\mathrm du\,\mathrm dv.
	\]
	因此 \(C\) 光滑，只需计算 \(C(0,0)\)。

	设联络矩阵为 \(A=\sum A_i\,\mathrm dx^i\)。把标架右乘正交矩阵
	\[
		k(x)=\exp\left(-\sum_i x^iA_i(q)\right)
	\]
	后，变换公式 \(\widetilde A=k^{-1}Ak+k^{-1}\mathrm dk\) 给出 \(\widetilde A(q)=0\)。新旧标架在 \(q\) 相同，故 \(L\) 不变。以下省略波浪号。若 \(\delta=|s|+|t|\)，沿每条边的联络积分为 \(O(\delta^2)\)；运输积分方程的二次迭代项为 \(O(\delta^4)\)。把 \(A\) 在原点作一次 Taylor 展开，四条边各自运输的二次项依次为
	\[
		-\frac{s^2}{2}\partial_1A_1,\qquad
		-st\partial_1A_2-\frac{t^2}{2}\partial_2A_2,\qquad
		\frac{s^2}{2}\partial_1A_1+st\partial_2A_1,\qquad
		\frac{t^2}{2}\partial_2A_2.
	\]
	这里各偏导均在 \(q\) 取值。乘起四个运输矩阵时，二次项直接相加，其余项为 \(O(\delta^3)\)。故混合二阶项为
	\[
		-st(\partial_1A_2-\partial_2A_1)(q)
		=-stR^E_q(\partial_1,\partial_2),
	\]
	最后用到 \(R^E=\mathrm dA+A\wedge A\) 及 \(A(q)=0\)。于是 \(C(0,0)=-R^E_q(\partial_1,\partial_2)\)。\(C\) 光滑便给出 \(C(s,t)=C(0,0)+O(|s|+|t|)\)，从而得到第一式；对恒等式 \(L-I=stC\) 求 \(s\) 导数并令 \(s=0\)，又得第二式。
\end{proof}

这也恢复了第一卷第18.3节小方环公式在向量丛中的形式。它表示曲率是单位面积所产生的无穷小旋转，负号来自既定的绕行方向和平行方程。若要把曲率证明为完整群 Lie 代数中的元素，\eqref{g3:eq:small-rectangle-first-variation}比单独的 \(L(\varepsilon,\varepsilon)\) 更方便：固定一条边后，另一条边的变化已经是一条在单位元出发的光滑群曲线。

\subsection{沿路径搬运的曲率算子}
\label{g3:subsec:0501:05}

不同点的曲率作用在不同纤维上，不能直接相加。若 \(\gamma:p\to q\)，就用平行运输把 \(q\) 处的曲率变回 \(E_p\) 的算子：
\begin{equation}\label{g3:eq:transported-curvature}
	\mathcal R_\gamma(X,Y)=P_\gamma^{-1}R^E_q(X,Y)P_\gamma,
	\qquad X,Y\in\mathrm T_qM.
\end{equation}
下面的交换图说明了输入和输出怎样比较：
\[
	\begin{tikzcd}[column sep=large]
		E_p \arrow[r,"P_\gamma"] \arrow[d,"{\mathcal R_\gamma(X,Y)}"']
		& E_q \arrow[d,"{R^E_q(X,Y)}"]\\
		E_p \arrow[r,"P_\gamma"'] & E_q.
	\end{tikzcd}
\]
路径 \(\gamma\) 的选择是数据的一部分。同一个终点，用不同路径搬回的算子可能不同。定义有限维线性空间
\begin{equation}\label{g3:eq:transported-curvature-span}
	\mathfrak r_p=\operatorname{span}_{\mathbb R}
	\{P_\gamma^{-1}R^E_q(X,Y)P_\gamma:
	\gamma:p\to q,\ X,Y\in\mathrm T_qM\}
	\subseteq\mathfrak{so}(E_p).
\end{equation}
在基点前接一个环路，可知 \(\mathfrak r_p\) 在 \(\operatorname{Hol}_p\) 的共轭作用下不变。此时还没有假定它对括号封闭。

\ProseParagraphBreak

曲率沿整条路径的贡献由运输变分公式表达。第一卷第18.3节在主丛上证明过这一公式；这里把固定端点的向量丛版本完整写出，便于直接看出所有被积算子都属于 \(\mathfrak r_p\)。

\begin{proposition}\label{g3:prop:holonomy-variation}
	设 \(\Gamma(s,t)\) 为上述具有共同有限分段的光滑基点环路族，\(\Gamma(s,0)=\Gamma(s,1)=p\)。记
	\[
		T=\partial_t\Gamma,\qquad V=\partial_s\Gamma,\qquad
		U(s,t)=P_{\Gamma(s,\cdot)|_{[0,t]}},\qquad a(s)=U(s,1).
	\]
	则
	\begin{equation}\label{g3:eq:holonomy-variation}
		a(s)^{-1}a'(s)
		=\int_0^1 U(s,t)^{-1}R^E_{\Gamma(s,t)}(T,V)U(s,t)\,\mathrm dt.
	\end{equation}
	积分按各光滑段相加，取值于固定的 \(\operatorname{End}(E_p)\)。特别地，左边属于 \(\mathfrak r_p\)。
\end{proposition}
\begin{proof}
	固定 \(v\in E_p\)，令 \(Z(s,t)=U(s,t)v\)，视为拉回丛 \(\Gamma^*E\) 的截面。记其沿参数方向的协变导数为 \(D_s,D_t\)。平行性给出 \(D_tZ=0\)。拉回联络的曲率恒等式为
	\[
		D_tD_sZ-D_sD_tZ=R^E(T,V)Z.
	\]
	这个式子不要求 \(\Gamma\) 为浸入：在任一局部标架内写 \(D_t=\partial_t+A(T)\)、\(D_s=\partial_s+A(V)\)，交换两算子后正是 \((\mathrm dA+A\wedge A)(T,V)\)。因此
	\[
		\frac{\mathrm d}{\mathrm dt}\bigl(U(s,t)^{-1}D_sZ(s,t)\bigr)
		=U(s,t)^{-1}R^E(T,V)U(s,t)v.
	\]
	起点固定且 \(Z(s,0)=v\)，故初值项为零。终点也固定，故 \(D_sZ(s,1)=a'(s)v\)。积分便得到所需公式作用于 \(v\) 的结果。

	若有分段点，则 \(Z(s,t)\) 从两侧的取值相同，并且作为 \(s\) 的光滑截面相同，所以 \(D_sZ\) 的两侧端值也相同。逐段积分时内部端点项抵消，公式仍成立。每个被积算子都是沿路径初段搬回的曲率；\(\mathfrak r_p\) 是有限维线性子空间，积分的 Riemann 和及其极限均在其中。
\end{proof}

输入次序在这里是 \((T,V)\)：先沿路径的方向，再沿改变路径的方向。将它们互换会改变符号。公式还说明，有限环路的运输通常由一族彼此不交换的曲率算子共同决定，不能把它直接写成未作运输比较的曲率面积积分的指数。

\subsection{Ambrose–Singer 定理}
\label{g3:subsec:0501:06}

Ambrose\footnote{Warren Ambrose（1914--1995），美国数学家，研究微分几何、偏微分方程和概率论。}与 Singer\footnote{Isadore Manuel Singer（1924--2021），美国数学家，研究微分几何和全局分析，并参与建立椭圆算子的指标理论。}的\term[ambrose singer dingli]{Ambrose--Singer 定理}[Ambrose--Singer holonomy theorem]把前面的两个方向合在一起：小矩形产生曲率方向，而一切环路族的变化都由搬运后的曲率给出。

\begin{theorem}[Ambrose--Singer]\label{g3:thm:ambrose-singer}
	设 \(M\) 为连通、有限维、Hausdorff、第二可数且无边界的光滑流形，\((E,\langle\ ,\ \rangle_E,\nabla^E)\) 为有限秩度量向量丛。给完整群赋予\cref{g3:thm:holonomy-lie-group}的浸入 Lie 群结构，则
	\begin{equation}\label{g3:eq:ambrose-singer}
		\mathfrak{hol}_p
		:=\operatorname{Lie}(\operatorname{Hol}_p^0)
		=\operatorname{span}_{\mathbb R}
		\{P_\gamma^{-1}R^E_q(X,Y)P_\gamma:
		\gamma:p\to q,\ X,Y\in\mathrm T_qM\}.
	\end{equation}
	特别地，右侧的线性张成已经是 Lie 子代数。
\end{theorem}
\symidx{\(\mathfrak{hol}_p\)}
\begin{proof}
	沿用完整群结构定理中的 \(\mathfrak k\)，已知它等于 \(\mathfrak{hol}_p\)。先取其一个生成方向 \(a'(0)\)，其中 \(a(s)\) 来自光滑环路族且 \(a(0)=I\)。由\cref{g3:prop:holonomy-variation}，
	\[
		a'(0)=\int_0^1U(0,t)^{-1}R^E(T,V)U(0,t)\,\mathrm dt\in\mathfrak r_p.
	\]
	所以 \(\mathfrak k\subseteq\mathfrak r_p\)。

	反过来，固定 \(\gamma:p\to q\)。先取 \(q\) 处两个坐标基向量 \(\partial_i,\partial_j\)。在对应坐标平面内作\cref{g3:prop:small-rectangle-holonomy}的矩形，记其运输为 \(L_q(s,t)\)。沿 \(\gamma\) 出发、绕矩形、再沿 \(\bar\gamma\) 返回的运输为
	\[
		b_t(s)=P_\gamma^{-1}L_q(s,t)P_\gamma.
	\]
	对每个固定的小 \(t\)，它是满足 \(b_t(0)=I\) 的光滑环路族运输，故 \(b_t'(0)\in\mathfrak k\)。当 \(t\ne0\) 时，\eqref{g3:eq:small-rectangle-first-variation}给出
	\[
		\frac{b_t'(0)}t
		=-P_\gamma^{-1}R^E_q(\partial_i,\partial_j)P_\gamma+O(t).
	\]
	让 \(t\to0\)，利用有限维线性子空间 \(\mathfrak k\) 的闭性，得到右边主项属于 \(\mathfrak k\)。任意 \(X,Y\) 的情形由双线性推出，所以 \(\mathfrak r_p\subseteq\mathfrak k\)。当 \(\dim M\leq1\) 时曲率二形式恒为零，第一步已给出 \(\mathfrak k=0\)，结论同样成立。
\end{proof}

证明中的极限始终在有限维向量空间 \(\mathfrak{so}(E_p)\) 内进行，没有用到完整群在正交群中闭。定理中的全部路径也有实质作用：只看 \(R^E_p\)，一般不足以看见其他位置所产生的旋转。下一节将把这个 Lie 代数对张量的作用转成平行结构。

\begin{corollary}\label{g3:cor:flat-restricted-holonomy}
	连通光滑流形上的度量联络平坦，当且仅当一个、从而每个基点的受限完整群为 \(\{I\}\)。若底流形还单连通，则这又等价于整体完整群为 \(\{I\}\)。
\end{corollary}
\begin{proof}
	平坦时\eqref{g3:eq:ambrose-singer}给出 \(\mathfrak{hol}_p=0\)，而受限完整群连通，所以它只有单位元。反过来，若受限完整群平凡，则所有搬回的曲率均为零；任取到 \(q\) 的路径并用运输可逆性，得到 \(R^E_q=0\)。单连通情形再用\cref{g3:prop:restricted-holonomy-normal}。
\end{proof}

\begin{corollary}\label{g3:cor:parallel-curvature-holonomy}
	设 \((M,g)\) 为连通光滑黎曼流形。若 \(\nabla R=0\)，则
	\[
		\mathfrak{hol}_p(M,g)
		=\operatorname{span}_{\mathbb R}\{R_p(u,v):u,v\in\mathrm T_pM\}.
	\]
\end{corollary}
\begin{proof}
	沿任意 \(\gamma:p\to q\) 平行运输 \(u,v,w\)。由 \(\nabla R=0\)，\(R(P_\gamma u,P_\gamma v)P_\gamma w\) 也沿路径平行，故
	\[
		P_\gamma^{-1}R_q(P_\gamma u,P_\gamma v)P_\gamma=R_p(u,v).
	\]
	因为 \(P_\gamma\) 满射，\eqref{g3:eq:ambrose-singer}中的每个生成元都化为基点曲率算子；反向包含取常值路径即可。
\end{proof}

% 纲要条目排为编号小节标题，并进入目录与书签；说明段仍保留于纲要注释中。

% !TeX root = ../../Book3.tex
% 纲要标题：### 5.2 平行张量与分布
% 纲要位置：几何学计划纲要.md，第 2273 行。
% 正文状态：初稿；尚未正式审核。

% 纲要内容开始
% 
% * 不变向量与平行截面
% * 不变子空间和正交分布
% * 平行分布的可积性
% * 局部乘积的构造
% 
% 纲要内容结束

\section{平行张量与分布}
\label{g3:sec:0502}

完整群把所有闭路的平行移动集中在一个切空间中。若一份张量在闭路运输后保持不变，就可以沿任意道路把它送到其他点，而不必记住所选道路。对一个子空间作同样的运输，则得到随点变化的分布。Levi-Civita 联络的无挠性和度量相容性会把这些线性不变量进一步转化为全测地叶与局部黎曼乘积。

\subsection{不变向量与平行截面}
\label{g3:subsec:0502:01}

本节的黎曼流形均光滑、有限维且无边界。对分片光滑道路 \(\gamma\)，记其切向平行移动为 \(P_\gamma\)，沿用式\eqref{g3:eq:holonomy-transport-laws}的拼接与逆路约定。在
\[
	\mathcal T^a_bM=(\mathrm{T}M)^{\otimes a}
		\otimes(\mathrm{T}^*M)^{\otimes b}
\]
上使用诱导联络；其平行移动在向量槽上使用 \(P_\gamma\)，在协向量槽上使用 \((P_\gamma^{-1})^*\)。这些操作与张量积、配对和收缩相容。满足 \(\nabla S=0\) 的张量场称为\term[pingxingzhangliang]{平行张量}[parallel tensor]。沿道路的张量场平行，当且仅当它的每个值都是初值经过上述运输所得。

\begin{theorem}\label{g3:thm:holonomy-parallel-tensors}
	设 \((M,g)\) 为连通光滑黎曼流形，\(p\in M\)，\(a,b\) 为非负整数，\(\nabla\) 为 Levi-Civita 联络在张量丛上的诱导联络，\(\mathcal T^a_bM=(\mathrm TM)^{\otimes a}\otimes(\mathrm T^*M)^{\otimes b}\)。在张量空间 \(\mathcal T^a_b(\mathrm{T}_pM)\) 上使用整体完整群 \(\operatorname{Hol}_p(M)\) 的自然作用。取值映射
	\[
		\{S\in\Gamma(\mathcal T^a_bM):\nabla S=0\}
		\longrightarrow
		\mathcal T^a_b(\mathrm{T}_pM)^{\operatorname{Hol}_p(M)},
		\qquad S\longmapsto S_p
	\]
	为线性同构。特别地，整体平行向量场恰与整体完整群固定的切向量对应。
\end{theorem}
\begin{proof}
	若 \(\nabla S=0\)，沿任一道路的限制都是平行场，所以
	\(S_q=P_\gamma S_p\)。当 \(\gamma\) 是基于 \(p\) 的闭路时，终值又是 \(S_p\)，故 \(S_p\) 为完整群不变量。连通流形上任意两点可由分片光滑道路连接，同一等式还说明 \(S_p\) 决定全部 \(S\)，因此取值映射单射。

	反过来，设 \(s\) 为该群固定的一个张量。对任意 \(q\)，选从 \(p\) 到 \(q\) 的道路 \(\gamma\)，定义
	\[
		S_q=P_\gamma s.
	\]
	若另选 \(\beta\)，则先走 \(\gamma\) 再走 \(\beta\) 的反路得到一个基于 \(p\) 的闭路，因而
	\(P_\beta^{-1}P_\gamma s=s\)。所以定义与道路无关。

	还要检验光滑性。固定 \(q_0\) 及一条从 \(p\) 到 \(q_0\) 的道路，在 \(q_0\) 的法坐标小球内，用径向测地线连接 \(q_0\) 与变化的点 \(q\)。平行移动满足线性常微分方程，其系数及初值光滑依赖 \(q\)；上一节的平行移动构造因而给出 \(S_q\) 对 \(q\) 的光滑依赖。最后，对任意从 \(q_0\) 出发的光滑曲线，把原来连接 \(p,q_0\) 的道路与该曲线初段拼接，由道路无关性得到
	\(S_{\gamma(t)}=P_{\gamma|_{[0,t]}}S_{q_0}\)。它沿每条曲线平行，故 \(\nabla S=0\)。整个构造线性，且是取值映射的逆。
\end{proof}

这里必须使用所有闭路产生的整体完整群。仅由可缩闭路固定的张量，可能在绕一个不可缩闭路后改变；它首先给出万有覆盖上的平行张量，能否下降还要检查这些额外闭路。

\ProseParagraphBreak

例如 \(g\) 本身平行，故完整群保持内积。对连通的定向黎曼流形，正体积形式也平行：在正交单位标架中，联络矩阵斜对称，其迹为零，所以对顶次外积的联络作用为零。这等价于整体完整群包含在 \(\operatorname{SO}(\mathrm{T}_pM)\) 内。反过来，若完整群保持某个非零顶次交替形式，上述定理给出一个处处非零的平行顶次形式，从而确定 \(M\) 的定向。

\begin{proposition}\label{g3:prop:parallel-tensor-constant-data}
	设 \((M,g)\) 为连通光滑黎曼流形，平行性均相对于 Levi-Civita 联络及其张量诱导联络。平行张量的范数为常数；平行自同态场 \(A:\mathrm{T}M\to\mathrm{T}M\) 在不同点的值由平行移动共轭，所以其特征多项式与各特征空间的维数不随点变化。
\end{proposition}
\begin{proof}
	诱导联络与张量内积相容，故
	\(X|S|^2=2\langle\nabla_XS,S\rangle=0\)。若 \(A\) 平行，沿从 \(p\) 到 \(q\) 的道路有
	\[
		A_qP_\gamma=P_\gamma A_p,
		\qquad A_q=P_\gamma A_pP_\gamma^{-1}.
	\]
	共轭保持特征多项式及各特征空间维数，得到结论。
\end{proof}

\subsection{不变子空间和正交分布}
\label{g3:subsec:0502:02}

一个完整群不变子空间不要求其中每个向量都固定。运输后保留下来的是整块方向空间，而不是预先选定的基。这对应于平行分布。

\begin{definition}\label{g3:def:parallel-distribution}
	设 \((M,g)\) 为光滑黎曼流形。记 \(\nabla\) 为其 Levi-Civita 联络。常秩光滑子丛 \(E\subset\mathrm{T}M\) 若满足
	\[
		\nabla_XY\in\Gamma(E)
		\qquad
		\bigl(X\in\Gamma(\mathrm{T}M),\ Y\in\Gamma(E)\bigr),
	\]
	则称为\term[pingxingfenbu]{平行分布}[parallel distribution]。求导方向 \(X\) 可以是任意切向方向。
\end{definition}

\begin{proposition}\label{g3:prop:parallel-invariant-distributions}
	设 \((M,g)\) 为连通光滑黎曼流形，\(\nabla\) 为其 Levi-Civita 联络，\(P_\gamma\) 为沿分片光滑道路的平行移动，固定 \(p\in M\)。整体完整群不变子空间 \(V\subset\mathrm{T}_pM\) 与 \(M\) 上的平行分布一一对应，映射由
	\[
		E_q=P_\gamma V,\qquad \gamma:p\longrightarrow q
	\]
	给出。平行分布的正交补也平行；对 \(E\) 的正交投影 \(\Pi_E\) 满足 \(\nabla\Pi_E=0\)。\symidx{\(\Pi_E\)：到平行分布 \(E\) 的正交投影}
\end{proposition}
\begin{proof}
	若 \(V\) 为完整群不变子空间，对两条同端点道路 \(\gamma,\beta\)，
	\(P_\beta^{-1}P_\gamma V=V\)，所以 \(E_q\) 与道路无关。固定一点附近的径向道路，运输一个 \(V\) 的基，得到光滑且处处线性无关的局部标架，故 \(E\) 是秩 \(\dim V\) 的光滑子丛。

	这个子丛被所有平行移动保持。为把它改写成联络条件，沿一条曲线 \(\gamma\) 将 \(E_{\gamma(t)}\) 运输回起点。若 \(Y(t)\) 始终取值于 \(E_{\gamma(t)}\)，则 \(P_{\gamma|_{[0,t]}}^{-1}Y(t)\) 始终属于同一个线性子空间 \(E_{\gamma(0)}\)，其导数也在该空间内。该导数运输回 \(\gamma(t)\) 正是 \(\mathrm D_tY\)，所以 \(\mathrm D_tY\in E\)。对局部场及任意曲线应用此式，得到平行分布条件。

	反过来，若 \(E\) 平行，在局部标架中沿曲线求平行场时，以 \(E\) 的基作为前几项，求导不会从这些基产生正交补分量。因此在 \(E\) 内解同一线性方程，就得到初值在 \(E\) 内的平行场；由唯一性，它也是整个切丛中的平行场。沿有限段坐标域拼接，平行移动始终保持 \(E\)。特别地，闭路保持 \(E_p\)，而道路运输公式又恢复整个 \(E\)。

	对 \(Z\in\Gamma(E^\perp)\)、\(Y\in\Gamma(E)\)，有
	\[
		g(\nabla_XZ,Y)=Xg(Z,Y)-g(Z,\nabla_XY)=0.
	\]
	所以 \(E^\perp\) 平行。将任意场写成 \(Y=Y_E+Y_{E^\perp}\)，两个分量的导数仍在各自子丛，故
	\(\nabla_X(\Pi_EY)=\Pi_E(\nabla_XY)\)，即 \(\nabla\Pi_E=0\)。
\end{proof}

例如一个自伴平行自同态 \(A\) 的不同特征空间正交。其特征值在连通流形上固定，且由运输共轭，各特征空间组成平行子丛。也可以直接把对应投影写成
\[
	\Pi_\lambda
	=\prod_{\mu\ne\lambda}\frac{A-\mu I}{\lambda-\mu},
\]
其中乘积遍历其他不同特征值；这是光滑平行的自同态，像正是所求子丛。

\subsection{平行分布的可积性}
\label{g3:subsec:0502:03}

分布只指定每点允许的方向。要把这些方向拼成子流形，需要验证沿这些方向行走所得的交换子仍不离开分布。无挠性在这里把协变导数与 Lie 括号联系起来。

\begin{proposition}\label{g3:prop:parallel-distribution-leaves}
	设 \(E\) 为光滑黎曼流形 \((M,g)\) 上的平行分布。则 \(E\) 可积，具有唯一的最大连通积分叶，各叶以诱导度量为全测地浸入子流形。若 \(M\) 完备，则每条叶对其自身的诱导距离完备。
\end{proposition}
\begin{proof}
	若 \(X,Y\) 取值于 \(E\)，则
	\[
		[X,Y]=\nabla_XY-\nabla_YX\in\Gamma(E),
	\]
	所以 \(E\) 对合。第一卷第13.4节的\term[frobeniusdingli]{Frobenius 定理}[Frobenius theorem]说明，光滑常秩对合分布在每点附近具有叶状坐标；同节的最大连通积分叶定理将这些叶片拼成唯一的叶，并赋予其 Hausdorff、第二可数的光滑流形结构。\footnote{Ferdinand Georg Frobenius（1849--1917），德国数学家，研究群的表示与特征标以及微分方程。}
	% 跨卷依赖：g1:thm:frobenius、g1:thm:maximal-connected-leaf，已核读陈述与证明。
	这里 \(E\) 的光滑性、常秩与刚验证的对合性正好满足这些假设。叶采用这个叶拓扑；它未必是 \(M\) 的嵌入子流形。

	在一条叶上，\(\nabla_XY\) 对切向场 \(X,Y\) 仍在 \(E\)，故叶的第二基本形式为零。由\cref{g3:prop:totally-geodesic-characterization}，叶内测地线也是环境测地线。

	设 \(M\) 完备。在任一叶点取叶切向初值，环境测地线 \(\gamma\) 全时存在。它的速度沿自身平行，初速度在 \(E\)，故全部速度都在 \(E\)。在每个叶状坐标盒内，横向坐标的导数为零，所以曲线的每个小段都落在一张叶片中。沿时间区间拼接，它始终属于起点所在的最大叶，并且作为到该叶的映射光滑。于是任意叶测地线初值都全时可解，\cref{g3:thm:hopf-rinow}给出叶的诱导度量完备性。
\end{proof}

正交补同时平行，因而同时可积。这比“某一分布的叶全测地”更强：平行性还要求沿横向求导时保留该分布，这正是下一步消去两个因子之间度量变化的条件。

\subsection{局部乘积的构造}
\label{g3:subsec:0502:04}

先把两族叶片同时变成坐标切片，再检查度量是否真正分离。正交性将消去混合项；平行性则使每个对角块只依赖自己的坐标。

\begin{theorem}\label{g3:thm:parallel-local-product}
	设 \((M,g)\) 为光滑黎曼流形，且
	\(\mathrm{T}M=E\oplus F\) 为两个正交平行分布的直和。则每点附近存在乘积坐标
	\((x^1,\ldots,x^r,y^1,\ldots,y^s)\)，其中 \(r=\operatorname{rank}E\)、\(s=\operatorname{rank}F\)，使
	\[
		E=\operatorname{span}\{\partial_{x^i}\},\qquad
		F=\operatorname{span}\{\partial_{y^\alpha}\},
	\]
	并且
	\begin{equation}\label{g3:eq:parallel-local-product-metric}
		g=g_{ij}(x)\,\mathrm dx^i\mathrm dx^j
			+h_{\alpha\beta}(y)\,\mathrm dy^\alpha\mathrm dy^\beta.
	\end{equation}
	因此该邻域与两张局部积分叶片的黎曼乘积等距。
\end{theorem}
\begin{proof}
	由两分布的可积性，在共同的小邻域上存在独立函数
	\(y^1,\ldots,y^s\)，使 \(E=\ker\mathrm dy\)，也存在
	\(x^1,\ldots,x^r\)，使 \(F=\ker\mathrm dx\)。联合映射
	\((x,y)\) 的微分若杀掉一个向量，该向量同时属于 \(E,F\)，只能为零；维数相等，微分因而可逆。逆函数定理给出联合坐标，将像缩小为乘积开盒，便得到同时适应两分布的坐标。

	正交性给出 \(g(\partial_{x^i},\partial_{y^\alpha})=0\)。平行性又说明
	\[
		\nabla_{\partial_{x^i}}\partial_{y^\alpha}\in F,\qquad
		\nabla_{\partial_{y^\alpha}}\partial_{x^i}\in E.
	\]
	坐标向量场交换，无挠性使这两个导数相等；它们属于正交交集，故均为零。因此
	\[
		\begin{aligned}
			\partial_{y^\alpha}g(\partial_{x^i},\partial_{x^j})
			&=g(\nabla_{\partial_{y^\alpha}}\partial_{x^i},\partial_{x^j})
				+g(\partial_{x^i},\nabla_{\partial_{y^\alpha}}\partial_{x^j})
			=0.
		\end{aligned}
	\]
	同理 \(\partial_{x^i}g(\partial_{y^\alpha},\partial_{y^\beta})=0\)。各块遂只依赖自己的一组坐标，得到式\eqref{g3:eq:parallel-local-product-metric}。固定 \(y\) 或固定 \(x\) 后取得的正是两个因子的诱导度量。
\end{proof}

\begin{corollary}\label{g3:cor:parallel-mixed-curvature}
	设 \((M,g)\) 为光滑黎曼流形，\(\mathrm TM=E\oplus F\) 是相对于 Levi-Civita 联络的正交平行分解，\(R\) 为该联络的曲率。若 \(X\in E\)、\(Y\in F\)，则 \(R(X,Y)=0\)。此外，\(R(E,E)\) 在 \(F\) 上的作用为零，\(R(F,F)\) 在 \(E\) 上的作用为零。
\end{corollary}
\begin{proof}
	局部乘积坐标中的混合协变导数为零，\(E\) 块的联络系数只依赖 \(x\)，\(F\) 块的联络系数只依赖 \(y\)。代入曲率坐标公式，混合求导方向的全部系数为零，同一因子的曲率也不作用到另一因子。张量性使结论对任意所列向量成立。
\end{proof}

在局部乘积盒中，一条 \(E\) 叶片道路可搬到相邻的 \(E\) 叶片上，保持 \(x\) 坐标而只改变固定的 \(y\) 值；此搬运是局部等距。同样可以搬运 \(F\) 道路。两个操作在共同定义的参数矩形上交换，且相对的边有相同长度。这是把局部乘积推广到整体的基本构造。

\begin{example}\label{g3:exm:orthogonal-foliations-not-parallel}
	在平面去掉原点后使用极坐标，度量为
	\[
		g=\mathrm dr^2+r^2\mathrm d\theta^2.
	\]
	径向线与圆周给出两族正交可积的一维分布，但圆周分布不平行，因为
	\(\nabla_{\partial_\theta}\partial_\theta=-r\partial_r\)。
	相应对角块 \(r^2\mathrm d\theta^2\) 依赖另一组坐标，故这两族指定分布不能产生黎曼乘积。可积与正交本身不足以替代平行性。
\end{example}
% 纲要条目排为编号小节标题，并进入目录与书签；说明段仍保留于纲要注释中。

% !TeX root = ../../Book3.tex
% 纲要标题：### 5.3 de Rham 分解
% 纲要位置：几何学计划纲要.md，第 2280 行。
% 正文状态：初稿；尚未正式审核。

% 纲要内容开始
% 
% * 局部分解的存在
% * 单连通完备情形的整体乘积
% * 因子的不可约性与唯一性
% * 一般商空间中的扭曲现象
% 
% 纲要内容结束

\section{de Rham 分解}
\label{g3:sec:0503}

上一节已经把完整群的不变子空间变成正交平行分布，并把这样的分布变成局部乘积。本节要解决两个整体问题：局部乘积能否沿道路继续铺开，所得因子是否由度量唯一决定。\term[derhamfenjie]{de Rham 分解}[de Rham decomposition]回答这两个问题；整体定理中的完备性与单连通性各有明确作用。\footnote{Georges de Rham（1903--1990），瑞士数学家，研究微分拓扑、微分形式与代数拓扑。}

\subsection{局部分解的存在}
\label{g3:subsec:0503:01}

设在连通黎曼流形上已经得到有限个彼此正交的平行分布
\[
	\mathrm{T}M=E_0\oplus E_1\oplus\cdots\oplus E_s.
\]
将\cref{g3:thm:parallel-local-product}先用于 \(E_0\) 与其正交补，再在第二个因子中继续应用，便在每点附近得到
\begin{equation}\label{g3:eq:multiple-local-product}
	(U,g|_U)\cong
	(U_0,g_0)\times\cdots\times(U_s,g_s).
\end{equation}
每个 \(E_i\) 对应一个因子的切空间，\(g_i\) 只依赖该因子的坐标。特别地，给定完整群的一个正交不变分解，\cref{g3:prop:parallel-invariant-distributions}便给出这种局部分解。

这里的乘积坐标还有一项重要的相容性。若两套坐标 \((x,y)\)、\((x',y')\) 都适应 \(E\oplus F\)，则在交集的一个充分小的连通乘积盒中，
\[
	x'=\varphi(x),\qquad y'=\psi(y).
\]
事实上，\(\mathrm dx'\) 在 \(F\) 上为零，故 \(x'\) 与 \(y\) 无关；另一式同理。因而“保持 \(x\) 而改变 \(y\)”这一局部搬运不依赖所用坐标。它在不同盒中接续时保持长度，也保持 \(E\)、\(F\) 两个方向。

局部结论只要求相应平行分布存在。下面把这些乘积盒沿有限长道路拼接；这一步将用到完备性，不能仅由每点附近的式\eqref{g3:eq:multiple-local-product}推出。

\subsection{单连通完备情形的整体乘积}
\label{g3:subsec:0503:02}

先证明一个把局部等距映射识别为覆盖的工具。这里映射须已在整个源流形上定义。

\begin{lemma}\label{g3:lem:complete-local-isometry-covering}
	设 \(\Phi:(N,h)\to(M,g)\) 是连通黎曼流形之间的局部等距映射，且 \(N\) 完备，则 \(\Phi\) 满射并且是覆盖映射。若 \(M\) 单连通，则 \(\Phi\) 是整体等距同构。
\end{lemma}
\begin{proof}
	局部等距保持 Levi-Civita 联络，因而保持带仿射参数的测地线。源流形完备，由 Hopf--Rinow 定理，其测地线对所有实参数存在。因此，\(M\) 中从 \(\Phi(z)\) 出发的任意有限测地线段，都可由初速度经 \((\mathrm d\Phi_z)^{-1}\) 提回后提升到 \(N\) 中。目标中的任意两点可以用有限条短测地线段连接，逐段提升便知 \(\Phi\) 满射。

	固定 \(q\in M\)，选 \(\rho>0\)，使 \(\exp_q\) 在切球 \(B(0,\rho)\) 上为微分同胚，记其像为 \(U\)。对每个 \(z\in\Phi^{-1}(q)\)，定义
	\[
		s_z(\exp_q v)=\exp_z\bigl((\mathrm d\Phi_z)^{-1}v\bigr),
		\qquad |v|<\rho.
	\]
	源流形测地完备使右端有定义；测地线唯一性给出
	\(\Phi\circ s_z=\operatorname{id}_U\)。因此 \(\mathrm ds_z\) 可逆，\(s_z\) 是到开集 \(V_z=s_z(U)\) 的微分同胚，逆映射为 \(\Phi|_{V_z}\)。

	若 \(V_z\) 与 \(V_{z'}\) 相交，在交点上提升从其像返回 \(q\) 的径向测地线。提升的初速度由 \(\mathrm d\Phi\) 唯一决定，测地线唯一性迫使终点相同，即 \(z=z'\)。反过来，任意 \(w\in\Phi^{-1}(U)\) 都能提升这条反向径向线，得到某个 \(z\in\Phi^{-1}(q)\)，于是 \(w\in V_z\)。故
	\[
		\Phi^{-1}(U)=\coprod_{z\in\Phi^{-1}(q)}V_z
	\]
	是均匀覆盖邻域。

	若 \(M\) 单连通，连通覆盖只有一层：从源中两点到同一目标点所成的投影闭路可以收缩，覆盖的道路与同伦提升唯一性使这两点相同。因此 \(\Phi\) 为微分同胚，局部等距性随即成为整体等距性。
	% 跨卷依赖：第一卷第5.3节 g1:thm:path-lifting、g1:thm:homotopy-lifting、g1:cor:lifted-path-homotopy，已核读。
\end{proof}

接下来构造上一引理所需的全域映射。设
\(\mathrm{T}M=E\oplus F\) 是正交平行分解。所谓由一条 \(E\) 道路 \(\alpha\) 和一条 \(F\) 道路 \(\beta\) 张成的\term[pingxingjuxing]{平行矩形}[parallel rectangle]，是一个分片光滑映射
\[
	Q:[0,1]^2\longrightarrow M,\qquad
	Q(s,0)=\alpha(s),\quad Q(0,t)=\beta(t),
\]
它在每个充分小的乘积坐标盒中具有形式
\[
	Q(s,t)=(x(s),y(t)).
\]
这一定义允许道路停顿、折返以及矩形像的自交。“矩形”指参数域及两个方向之间的乘积关系，并不要求其像是一个嵌入曲面。

\begin{figure}[htbp]
	\centering
	% !TeX root = ../Book3.tex
% 参数域示意：图中方格不表示 Q 的像在流形中嵌入。
\begin{tikzpicture}[x=1cm,y=1cm]
	\fill[gray!18] (1.2,0.9) rectangle (2.4,1.8);
	\draw[gray!55,thin] (0,0) grid[xstep=1.2,ystep=0.9] (6,3.6);
	\draw[thick] (0,0) rectangle (6,3.6);
	\draw[->,thick] (0,0) -- (6,0)
		node[midway,below=7pt] {\(\alpha(s)=Q(s,0)\)};
	\draw[->,thick] (0,0) -- (0,3.6)
		node[midway,left=7pt] {\(\beta(t)=Q(0,t)\)};
	\draw[->] (1.8,0.55) -- (1.8,1.25);
	\draw[->] (0.85,1.35) -- (1.55,1.35);
	\fill (0,0) circle (1.7pt);
	\fill (6,0) circle (1.7pt);
	\fill (0,3.6) circle (1.7pt);
	\fill (6,3.6) circle (1.7pt);
	\node[below left=4pt] at (0,0) {\(p\)};
	\node[below right=4pt] at (6,0) {\(\alpha(1)\)};
	\node[above left=4pt] at (0,3.6) {\(\beta(1)\)};
	\node[above right=4pt] at (6,3.6) {\(Q(1,1)\)};
	\node[above=5pt] at (3,3.6) {\([0,1]^2\) 中的有限细分};
\end{tikzpicture}

	\caption[平行矩形的逐格延拓]{平行矩形的逐格延拓。图示为参数域；每个小格由下边和左边确定，阴影格标出一次延拓。矩形像不要求在流形中嵌入。}
	\label{g3:fig:parallel-product-rectangle}
\end{figure}


\begin{lemma}\label{g3:lem:parallel-rectangle-extension}
	设 \(M\) 连通且完备，\(E,F\) 为正交互补平行分布。任给从同一点 \(p\) 出发的分片光滑 \(E\) 道路 \(\alpha\) 和 \(F\) 道路 \(\beta\)，存在唯一的平行矩形 \(Q_{\alpha,\beta}\)。其相对边的对应子段长度相等。矩形对共同有限分段的光滑道路族光滑依赖，且终点 \(Q_{\alpha,\beta}(1,1)\) 仅依赖两条道路在各自积分叶内的固定端点同伦类。
\end{lemma}
\begin{proof}
	在一个乘积盒中，已知下边 \((x(s),y_0)\) 和左边 \((x_0,y(t))\)，唯一的填法是 \((x(s),y(t))\)。由于度量的两个块各自独立，相对边对应子段等长。上一小节的换坐标公式说明，这些填法在交叠处相容。

	令 \(L=\operatorname{Len}(\alpha)+\operatorname{Len}(\beta)\)。Hopf--Rinow 定理使闭球
	\[
		K=\overline B(p,L+1)
	\]
	紧。用有限个较小乘积盒覆盖 \(K\)，并使每个较小盒的闭包包含在另一个乘积盒内部。可以选同一个 \(\delta>0\)，使从较小盒中任一点出发的两条长度小于 \(\delta\) 的 \(E\)、\(F\) 道路，在相应大盒内有上述坐标乘积。说明这一点时，在小盒和大盒之间再放入一个闭包紧含于大盒的中间盒，先用小盒闭包到中间盒边界的正距离保证短道路留在中间盒，再利用中间盒闭包上的度量下界控制坐标变化，使两个坐标分量合成后仍位于大盒内。有限覆盖允许取所有这些正数的最小值。

	将 \(\alpha,\beta\) 各自分成有限段，使每段长度小于 \(\delta\)，并包含原有的光滑分段点。先填左下角的小矩形，再逐列、逐行填入其余矩形。每次使用已经得到的下边与左边；局部等距搬运保证新边仍具有对应原道路段的长度。任一已构造的点，都可沿这两类边从 \(p\) 到达，总长度不超过 \(L\)。它始终位于 \(K\)，故同一个 \(\delta\) 在全过程有效，有限格点构造遂覆盖整个参数正方形。这也说明完备性在这里提供的是紧球内统一的坐标控制。

	若存在两个填法，取参数域的共同细分，使比较落在乘积盒内。一个小格的下边与左边确定整个格子，逐格归纳得二者相同。加细原分段不会改变结果。对一个具有共同有限分段的光滑道路族，可在每个参数值附近使用共同的有限格子和乘积盒；各步都是光滑坐标运算，因而所得矩形光滑依赖于参数。

	最后让 \(\alpha_u\) 在同一 \(E\) 叶内作固定端点的分片光滑同伦，暂时固定 \(\beta\)。在 \(Q_u=Q_{\alpha_u,\beta}\) 中，变分向量 \(\partial_uQ_u\) 始终属于 \(E\)：它在左边为零，而各格的 \(F\) 坐标不随 \(s\) 改变；该结论沿格子接续。沿竖直方向，这个 \(E\) 变分向量平行，因为在乘积坐标中其 \(E\) 分量不随 \(t\) 改变，混合协变导数也为零。右下角的变分向量为
	\(\partial_u\alpha_u(1)=0\)，故整个右边上的变分向量为零，特别是
	\[
		\partial_uQ_{\alpha_u,\beta}(1,1)=0.
	\]
	交换两个方向，可知改变 \(\beta\) 的固定端点叶内同伦也不改变终点。

	这里使用分片光滑同伦不限制通常的道路同伦类。对流形中的一个连续道路同伦，可把参数正方形细分得足够小，使每个格子的像落在坐标球中；先在边上以坐标线段替换，再在各格内用坐标线性收缩填充，边界上保持给定的分片光滑道路。如此得到的有限分片光滑同伦仍固定端点。此操作在积分叶自身的流形拓扑中进行，因而以上结论适用于通常的叶内同伦类。
\end{proof}

\begin{theorem}\label{g3:thm:global-parallel-product}
	设 \((M,g)\) 连通、单连通且完备，并有正交平行分解
	\(\mathrm{T}M=E\oplus F\)。固定 \(p\in M\)，以 \(L,N\) 分别表示通过 \(p\) 的最大连通 \(E\) 叶与 \(F\) 叶，赋予诱导度量。则 \(L,N\) 完备、单连通并且嵌入 \(M\)，存在等距同构
	\[
		\Phi:(L\times N,g_L\oplus g_N)\longrightarrow(M,g),
	\]
	使 \(\Phi(x,p)=x\)、\(\Phi(p,y)=y\)，且两个因子的切空间分别映到 \(E,F\)。
\end{theorem}
\begin{proof}
	\cref{g3:prop:parallel-distribution-leaves}已给出叶的流形结构、全测地性与完备性。此时不预设叶嵌入。取基于 \(p\) 的万有覆盖
	\[
		\pi_L:\widetilde L\longrightarrow L,\qquad
		\pi_N:\widetilde N\longrightarrow N.
	\]
	可使用第一卷第5.4节的道路类构造：\(\widetilde L\) 的点是从 \(p\) 出发的道路的固定端点同伦类，邻域由在终点附近追加短道路得到，\(\widetilde N\) 同理。流形局部单连通，故此构造适用；覆盖坐标给出光滑结构。取可数坐标球及其交集分支后，有限坐标球链给出可数覆盖图册，因而这些覆盖仍是第二可数流形。
	% 跨卷依赖：g1:thm:path-class-universal-cover，已核读陈述与证明。

	在覆盖空间上拉回叶的度量。它们仍完备：叶中的测地线对所有时间存在，每个有限参数区间上的道路提升相互相容，得到覆盖空间中对所有时间存在的测地线；局部等距保证提升仍满足测地线方程。再用 Hopf--Rinow 定理即可。由乘积联络公式\eqref{g3:eq:product-connection}，乘积中的测地线恰为两个因子测地线的配对，故
	\(\widetilde L\times\widetilde N\) 也完备。

	对道路类定义
	\[
		\widetilde\Phi([\alpha],[\beta])
		=Q_{\alpha,\beta}(1,1).
	\]
	\cref{g3:lem:parallel-rectangle-extension}同时保证它在整个乘积上有定义，并与道路代表无关。还需核实它是局部等距。在 \([\alpha]\) 附近用向终点追加一段短 \(E\) 道路来给覆盖坐标，在 \([\beta]\) 附近同样追加短 \(F\) 道路。原矩形沿这两段的延拓，在终点附近恰为乘积盒中的坐标乘积。只改变第一个端点时，所得变化是沿原 \(F\) 道路的逐盒等距搬运；只改变第二个端点时，是沿原 \(E\) 道路的逐盒等距搬运。因此两个方向的微分分别为到 \(E,F\) 的线性等距，并且像正交。矩形对参数的光滑依赖又保证 \(\widetilde\Phi\) 光滑，故它是局部等距。

	现在源流形连通且完备，\cref{g3:lem:complete-local-isometry-covering}说明 \(\widetilde\Phi\) 是到 \(M\) 的覆盖；目标单连通，故它是一对一的整体等距同构。在坐标轴上，
	\[
		\widetilde\Phi([\alpha],[\text{常路}])=\pi_L([\alpha]),
		\qquad
		\widetilde\Phi([\text{常路}],[\beta])=\pi_N([\beta]),
	\]
	其中右端再经叶到 \(M\) 的浸入。单射性迫使 \(\pi_L,\pi_N\) 都只有一层；两叶因此单连通。坐标轴在乘积中嵌入，经整体等距映射后便知两叶在 \(M\) 中嵌入，得到所述 \(\Phi\)。
\end{proof}

上述证明先用完备性完成矩形延拓，再从叶的道路类构造一个已在全域定义的局部等距映射，最后才使用覆盖结论。单连通性负责消除覆盖层数。有限个平行正交因子的情形可重复应用该定理；得到的每个因子仍完备且单连通。

\subsection{因子的不可约性与唯一性}
\label{g3:subsec:0503:03}

群 \(G\) 在非零内积空间 \(V\) 上的正交表示称为\term[bukeyuebiaoshi]{不可约表示}[irreducible representation]，如果仅有 \(0,V\) 两个 \(G\) 不变子空间。正交表示总可分解为有限个不可约正交子空间：选择一个维数最小的非零不变子空间，其正交补仍不变，再对正交补重复。这个有限维论证不要求 \(G\) 在正交群中闭。

\begin{theorem}\label{g3:thm:de-rham-decomposition}
	设 \((M,g)\) 为连通、单连通、完备的光滑黎曼流形。则存在等距分解
	\begin{equation}\label{g3:eq:de-rham-product}
		(M,g)\cong
		(\mathbb R^k,g_{\mathrm{Eucl}})
		\times(M_1,g_1)\times\cdots\times(M_s,g_s),
	\end{equation}
	其中每个 \(M_i\) 连通、单连通、完备、非平坦，且其 Levi-Civita 完整群在切空间上的表示不可约。允许 \(k=0\) 或 \(s=0\)，零维因子理解为一点。Euclid 因子的维数等于 \(M\) 上整体平行向量场空间的维数。
\end{theorem}
\begin{proof}
	固定 \(p\)，令 \(G=\operatorname{Hol}_p(M)\)，先取全部固定向量组成的子空间
	\[
		V_0=(\mathrm{T}_pM)^G,
	\]
	再将其正交补分解为不可约不变子空间
	\[
		\mathrm{T}_pM=V_0\oplus V_1\oplus\cdots\oplus V_s.
	\]
	每个 \(V_i\)（\(i\ge1\)）上的表示非平凡，否则其中的向量属于 \(V_0\)。这些子空间确定正交平行分布。逐次使用\cref{g3:thm:global-parallel-product}，得到完备单连通因子的整体乘积
	\[
		M\cong M_0\times M_1\times\cdots\times M_s.
	\]

	先识别 \(M_0\)。由\cref{g3:thm:holonomy-parallel-tensors}，\(V_0\) 的一组标准正交基延拓成整体平行场，限制到 \(M_0\) 后成为平行标准正交标架 \(X_1,\ldots,X_k\)。每个 \(X_i\) 的积分曲线满足
	\(\nabla_{X_i}X_i=0\)，因完备性而对所有时间存在。又由无挠性，
	\([X_i,X_j]=0\)。零 Lie 括号说明 \(X_i\) 的局部流保持 \(X_j\)：对流推前的向量场求时间导数，得到由 \([X_i,X_j]\) 控制的齐次方程，其初值为 \(X_j\)。因此先沿 \(X_i\) 流动后再沿 \(X_j\) 流动，与交换次序所得的曲线满足同一初值问题，局部流遂交换。将任意有限时间区间分段，这一关系逐段给出整体流的交换。记其流为 \(\phi_i^t\)，定义
	\[
		F:\mathbb R^k\longrightarrow M_0,\qquad
		F(t_1,\ldots,t_k)=\phi_1^{t_1}\circ\cdots\circ\phi_k^{t_k}(p).
	\]
	流的交换保证 \(\mathrm dF(\partial_{t_i})=X_i\)，故 \(F\) 是局部等距。源空间 \(\mathbb R^k\) 完备，目标 \(M_0\) 单连通，\cref{g3:lem:complete-local-isometry-covering}使 \(F\) 成为整体等距同构。若 \(k=0\)，该步骤给出一个点。

	乘积联络还给出完整群的精确关系
	\begin{equation}\label{g3:eq:product-holonomy}
		\operatorname{Hol}_{(p_0,\ldots,p_s)}(M)
		=\{I\}\times
		\operatorname{Hol}_{p_1}(M_1)\times\cdots\times
		\operatorname{Hol}_{p_s}(M_s).
	\end{equation}
	一条乘积闭路的平行移动按各因子的投影闭路分别进行，所以左端包含于右端；让闭路仅在一个因子中运动，又可分别实现右端的每一个群元素，故反向包含也成立。因此 \(V_i\) 上的不可约性正是 \(M_i\) 自身完整群的不可约性。若某个 \(M_i\) 平坦，\cref{g3:cor:flat-restricted-holonomy}与单连通性使其整体完整群平凡，和该表示非平凡矛盾。最后，平行场与固定向量的对应说明 Euclid 因子的维数恰为 \(k\)。
\end{proof}

\begin{proposition}\label{g3:prop:de-rham-uniqueness}
	分解\eqref{g3:eq:de-rham-product}中的 Euclid 因子及非平坦不可约因子，在因子重排和各因子的等距同构意义下唯一。固定基点后，它们在切空间中给出的子空间，以及通过该点的相应积分叶，在重排非平坦因子后也唯一。
\end{proposition}
\begin{proof}
	在第一种分解下记切空间为 \(V_0\oplus V_1\oplus\cdots\oplus V_s\)。式\eqref{g3:eq:product-holonomy}说明
	\(G=G_1\times\cdots\times G_s\)：\(G_i\) 只作用于 \(V_i\)，对其余子空间恒等。另一种分解中的 Euclid 子空间必须也是 \(G\) 的全部固定空间，因而仍为 \(V_0\)。

	设 \(W\) 为另一种分解的一个非平凡不可约子空间。正交投影
	\(P_i:W\to V_i\) 与 \(G\) 的作用交换。对于 \(i\ge1\)，非零 \(P_i\) 的核是 \(W\) 的不变子空间，像是 \(V_i\) 的不变子空间，所以它必为线性同构。投影 \(P_0\) 必为零：若非零，其核为零，而 \(V_0\) 上作用平凡，便使 \(W\) 上的作用也平凡，与假设矛盾。

	不能同时有 \(P_i,P_j\) 非零而 \(i\ne j\)。否则对 \(a\in G_i\) 与 \(w\in W\)，有
	\[
		P_j(aw-w)=aP_jw-P_jw=0.
	\]
	\(P_j\) 的单射性说明 \(aw=w\)。再由 \(P_i\) 满射，\(G_i\) 在 \(V_i\) 上的作用就平凡，仍为矛盾。因此 \(W\) 仅有一个非零投影，且 \(W=V_i\)。交换两种分解可知非平凡不可约子空间恰好一一对应。

	通过基点的相应叶由分布唯一决定，故两组因子逐一等距。Euclid 子空间整体唯一，但它内部的标准正交基可以旋转，因此再把 \(\mathbb R^k\) 分成 \(k\) 条直线时不具有这样的唯一性。
\end{proof}

证明中的关键是不同因子的完整群能够独立作用。一般抽象群的不可约正交分解，若有同构的表示重复出现，并不必然具有唯一的不变子空间；整体乘积所给出的式\eqref{g3:eq:product-holonomy}排除了这种混合。另一个直接结论是，式\eqref{g3:eq:de-rham-product}中的非平坦因子维数至少为 \(2\)，因为一维黎曼流形的曲率恒为零。

\subsection{一般商空间中的扭曲现象}
\label{g3:subsec:0503:04}

若 \(M\) 连通且完备，但未假设单连通，可先给万有覆盖
\(\pi:\widetilde M\to M\) 赋予拉回度量。和上述叶覆盖的证明一样，\(M\) 中的全时测地线可唯一提升，故 \(\widetilde M\) 完备。于是 de Rham 定理分解的是 \(\widetilde M\)。每个覆盖变换保持拉回度量，因而是等距变换；唯一性说明它保持整个 Euclid 分布，并可置换彼此等距的非平坦不可约因子。它还可能把一个因子中的平移和另一个因子中的等距变换联结在一起。取商后，局部乘积仍在，整体却未必是因子商的乘积。
% 跨卷依赖：第一卷第5.5节 g1:prop:universal-deck-discontinuous，已核读覆盖变换群在纤维上的自由传递作用。

\begin{example}\label{g3:exm:flat-mobius-holonomy}
	在 Euclid 平面 \(\mathbb R^2\) 上考虑等距变换
	\[
		\tau(s,t)=(s+1,-t),\qquad
		M=\mathbb R^2/\langle\tau\rangle.
	\]
	这个作用自由，因为非恒等偶次幂平移 \(s\)，奇次幂也使 \(s\) 改变奇整数。每个紧集只与有限个自身的平移像相交；在每点附近取足够小的矩形，可以使其不同平移像彼此不交。因此商映射是光滑覆盖，Euclid 度量下降为商度量。该曲面是无边界的\term[kaimobiusdai]{开 Möbius 带}[open Möbius band]。\footnote{August Ferdinand Möbius（1790--1868），德国数学家、天文学家，研究射影几何、重心坐标及单侧曲面。}平坦性由局部等距给出；任意测地初值可提升到 \(\mathbb R^2\) 中的一条全时直线，其投影为商中的全时测地线，故商测地完备，也即度量完备。
	% 商覆盖的拓扑依据为第一卷第5.5节 g1:prop:action-quotient-cover，已核读；完备性也可用本卷 g3:prop:riemannian-covering-distance。

	两条坐标直线分布均在 \(\tau\) 下不变，故下降为正交平行分布。向量场 \(\partial_s\) 也下降为整体平行场，而 \(\partial_t\) 经过一圈后变号。以原点的像为基点，提升生成元闭路后，沿 Euclid 平面的平行移动恒为恒等，终点切空间再由覆盖变换的微分识别，得到
	\[
		\operatorname{Hol}_p(M)
		=\left\{
			I,\begin{pmatrix}1&0\\0&-1\end{pmatrix}
		\right\},
		\qquad
		\operatorname{Hol}_p^0(M)=\{I\}.
	\]
	所有闭路的提升终点均为某个 \(\tau^m(0,0)\)，所以没有其他完整群元素。完整群的固定空间只有水平方向，说明这里的两条平行直线分布并不都由整体平行单位场生成。

	它也不可能等距于两个正维一维黎曼流形的乘积。这样的因子必完备；每个因子中的全时单位速测地线给出一个从 \(\mathbb R\) 出发的局部等距映射，由\cref{g3:lem:complete-local-isometry-covering}成为万有覆盖。其覆盖变换是直线的等距变换，而反射具有不动点，不可能是非恒等的覆盖变换。因此覆盖变换均为平移，直线的单位平行场可以下降，因子的完整群平凡。两个因子的乘积完整群仍平凡，与上式矛盾。这里整体乘积的障碍来自绕闭路后的反射，曲率已经处处为零。
\end{example}

更一般地，若 \(N\) 完备，\(\varphi:N\to N\) 是等距变换，\(a>0\)，则
\[
	(\mathbb R\times N)/\bigl((s,x)\sim(s+a,\varphi(x))\bigr)
\]
继承完备的局部乘积度量。沿圆周方向走一圈后，另一因子按照 \(\varphi\) 重新识别。只有分析这项整体识别，才能判断商是否仍可写成整体黎曼乘积；局部乘积和万有覆盖上的唯一分解本身不作此保证。
% 纲要条目排为编号小节标题，并进入目录与书签；说明段仍保留于纲要注释中。

% !TeX root = ../../Book3.tex
% 纲要标题：### 5.4 齐性与对称空间
% 纲要位置：几何学计划纲要.md，第 2287 行。
% 正文状态：初稿；尚未正式审核。
% 规划任务：保留现有齐性度量、点反射、对称分解与曲率模型，补充最大平坦、秩及紧非紧对偶的构造。完整核验球面、双曲空间与一个经典矩阵商的相应几何；局部 Lie 代数对偶与整体覆盖、离散商分别处理，不据局部对应推断任意整体商相同。

% 纲要内容开始
% 
% > 补充写作范围与外部输入：规划任务：保留现有齐性度量、点反射、对称分解与曲率模型，补充最大平坦、秩及紧非紧对偶的构造。完整核验球面、双曲空间与一个经典矩阵商的相应几何；局部 Lie 代数对偶与整体覆盖、离散商分别处理，不据局部对应推断任意整体商相同。
% 
% * 齐性度量及不变联络
% * 对称空间的点反射
% * Lie 代数分解与曲率
% * 紧型、非紧型和 Euclid 模型
% 
% * 极大交换子空间与最大平坦
% * 对称空间的秩及乘积结构
% * 紧型与非紧型的对偶
% 
% 
% 纲要内容结束

\section{齐性与对称空间}
\label{g3:sec:0504}

平行结构约束了一个流形上不同点的比较方式。若流形还有足够多的等距变换，许多计算可以集中在一个切空间里进行：群作用负责移动基点，稳定子负责变换切向方向。球面和双曲空间是秩一的例子；正定矩阵构成的空间还会出现维数大于一的全测地平坦子空间。为了看见这些差别，本节先把联络和曲率写成 Lie 括号，再追踪交换方向所产生的平坦。

\subsection{齐性度量及不变联络}
\label{g3:subsec:0504:01}

设有限维 Lie 群 \(G\) 光滑地、传递地作用在流形 \(M\) 上，固定 \(o\in M\)，令 \(K\) 为其稳定子。第一卷第17.1节的闭子群商定理给出 \(M\simeq G/K\)，其中 \(K\) 闭，\(o=eK\)，且
\[
	\mathrm T_o(G/K)\simeq\mathfrak g/\mathfrak k,
	\qquad \mathfrak g=\operatorname{Lie}(G),\quad
	\mathfrak k=\operatorname{Lie}(K).
\]
这里用到的具体结论是：商映射为光滑淹没，具有局部截面；稳定子在上述切空间上的作用由 \(\operatorname{Ad}(K)\) 诱导。它不要求 \(K\) 正规，\(G/K\) 通常不是群。

\begin{proposition}\label{g3:prop:homogeneous-metric-isotropy}
	设 \(G\) 为 Lie 群，\(K\) 为闭子群。\(G/K\) 上的 \(G\) 不变黎曼度量，与 \(\mathfrak g/\mathfrak k\) 上的 \(K\) 不变正定内积一一对应。若存在 \(\operatorname{Ad}(K)\) 不变分解
	\[
		\mathfrak g=\mathfrak k\oplus\mathfrak m,
	\]
	便可把这样的内积写在 \(\mathfrak m\) 上；当 \(K\) 紧时，总能选到这种分解和正定内积。
\end{proposition}
\begin{proof}
	不变度量在 \(o\) 的值必须被稳定子保持。反过来，给定这一内积，在 \(gK\) 处用左平移 \(L_g\) 搬运它。换代表 \(g\mapsto gk\) 时，差别正是稳定子 \(k\) 的切作用，所以定义不变。利用商映射的局部截面可见所得度量光滑，且构造与取基点值互逆。

	若 \(K\) 紧，将 \(\mathfrak g\) 上任意正定内积按其归一化不变测度平均，得到 \(\operatorname{Ad}(K)\) 不变内积，再取 \(\mathfrak m=\mathfrak k^\perp\)。这就是第一卷第19.1节的表示平均构造：积分保持正定性，变量替换保证不变性。
\end{proof}

带上述不变度量的空间称为\term[qixing liman kongjian]{齐性黎曼空间}[homogeneous Riemannian space]。在分解 \(\mathfrak g=\mathfrak k\oplus\mathfrak m\) 下，记括号的两个分量为 \([X,Y]_{\mathfrak k}\)、\([X,Y]_{\mathfrak m}\)。下面的公式同时说明怎样构造不变联络，以及度量怎样选出 Levi-Civita 联络。

\begin{proposition}\label{g3:prop:homogeneous-connection-formula}
	设 \(G/K\) 有上述不变分解与 \(\operatorname{Ad}(K)\) 不变正定内积。定义 \(\Lambda(X)Y\in\mathfrak m\) 为
	\begin{equation}\label{g3:eq:homogeneous-levi-civita}
		\begin{aligned}
		2\langle\Lambda(X)Y,Z\rangle
		={}&\langle[X,Y]_{\mathfrak m},Z\rangle
		-\langle[Y,Z]_{\mathfrak m},X\rangle
		+\langle[Z,X]_{\mathfrak m},Y\rangle.
		\end{aligned}
	\end{equation}
	它给出不变度量的 Levi-Civita 联络。其基点曲率为
	\begin{equation}\label{g3:eq:homogeneous-curvature}
		R_o(X,Y)
		=[\Lambda(X),\Lambda(Y)]-\Lambda([X,Y]_{\mathfrak m})
		-\operatorname{ad}([X,Y]_{\mathfrak k})|_{\mathfrak m}.
	\end{equation}
\end{proposition}
\begin{proof}
	把 \(G\) 上的左 Maurer--Cartan 形式按分解写成 \(\omega+\vartheta\)，其中 \(\omega\) 取值于 \(\mathfrak k\)，\(\vartheta\) 取值于 \(\mathfrak m\)。在商的局部截面上，\(\vartheta\) 把切向量表示为 \(\mathfrak m\) 中的列分量；改变截面时，\(\omega\) 按联络的规则变化，\(\vartheta\) 按稳定子表示变化。式\eqref{g3:eq:homogeneous-levi-civita}由不变内积与括号构成，故 \(\Lambda\) 也等变。因此
	\[
		\beta=\operatorname{ad}(\omega)|_{\mathfrak m}+\Lambda(\vartheta)
	\]
	是商切丛上良定义的不变联络矩阵。

	交换式\eqref{g3:eq:homogeneous-levi-civita}中的 \(Y,Z\)，两式相加，得到 \(\Lambda(X)\) 反对称。\(\operatorname{ad}(\mathfrak k)\) 也反对称，故该联络保持度量。再交换 \(X,Y\)，得
	\[
		\Lambda(X)Y-\Lambda(Y)X=[X,Y]_{\mathfrak m}.
	\]
	Maurer--Cartan 方程的 \(\mathfrak m\) 分量说明，联络的挠率在基点正是左边再减 \([X,Y]_{\mathfrak m}\)，所以为零。由 Levi-Civita 联络的唯一性，构造就是所求联络。

	最后计算 \(\mathrm d\beta+\beta\wedge\beta\)。在基点取水平输入 \(X,Y\in\mathfrak m\)，Maurer--Cartan 方程给出
	\[
		\mathrm d\omega(X,Y)=-[X,Y]_{\mathfrak k},\qquad
		\mathrm d\vartheta(X,Y)=-[X,Y]_{\mathfrak m}.
	\]
	于是外微分贡献式\eqref{g3:eq:homogeneous-curvature}的后两项，\(\beta\wedge\beta\) 贡献第一项。它与第一章采用的曲率矩阵约定一致。
\end{proof}

\subsection{对称空间的点反射}
\label{g3:subsec:0504:02}

球面上，把一点固定并同时反转所有切向方向，是一个全局等距变换。一般齐性空间只保证能移动点；能够在每一点实现这种反转，是更强的要求。

\begin{definition}\label{g3:def:riemannian-symmetric-space}
	设 \(M\) 为连通黎曼流形。若对每个 \(p\in M\)，存在全局等距变换 \(s_p\)，满足
	\[
		s_p(p)=p,\qquad (\mathrm ds_p)_p=-\mathrm{Id}.
	\]
	则称 \(M\) 为\term[liman duicheng kongjian]{黎曼对称空间}[Riemannian symmetric space]。此等距变换称为 \(p\) 处的\term[dian fanshe]{点反射}[geodesic symmetry]。
\end{definition}

一个连通黎曼流形上的等距变换，由一点的像及该点的微分唯一确定。事实上，它保持联络和测地线，在法邻域内满足 \(f(\exp_pv)=\exp_{f(p)}(\mathrm df_pv)\)；两个映射若初值一致，便在一个邻域相同。把“一阶数据相同”的点集看成既开又闭的集合，就能沿连通流形延伸这一相同性。因此 \(s_p\) 唯一，且 \(s_p^2=\mathrm{Id}\)。

\begin{proposition}\label{g3:prop:symmetric-complete-parallel-curvature}
	每个连通黎曼对称空间都测地完备，其等距变换群在点集上传递，并且 \(\nabla\mathrm{Rm}=0\)。
\end{proposition}
\begin{proof}
	若一条测地线 \(\gamma\) 只定义到有限正时间 \(b\)，取已定义的 \(a\) 满足 \(b/2<a<b\)。公式
	\[
		\widetilde\gamma(t)=s_{\gamma(a)}(\gamma(2a-t))
	\]
	把它经过 \(a\) 的部分反射到 \(a\) 之后，并延伸到 \(2a>b\)。在 \(a\) 处的位置和速度相同，测地线唯一性保证两段光滑接合。这与最大区间的有限右端点矛盾；左端点同理。

	由 Hopf--Rinow 定理，任意两点 \(p,q\) 可由一条极小测地线连接，其中点的反射把 \(p\) 送到 \(q\)，故等距群传递。最后，等距变换保持 \(\nabla\mathrm{Rm}\)。它是五阶协变张量，在 \(p\) 处被 \(\mathrm ds_p=-\mathrm{Id}\) 作用后变号，故只能为零。
\end{proof}

对称性也可由群直接构造。设 \(G\) 连通，\(\sigma:G\to G\) 是满足 \(\sigma^2=\mathrm{Id}\) 的光滑群自同构，并令闭子群 \(K\) 满足
\[
	(G^\sigma)_0\subseteq K\subseteq G^\sigma,
	\qquad G^\sigma=\{g:\sigma(g)=g\}.
\]
这样的数据称为\term[duicheng dui]{对称对}[symmetric pair]。微分 \(\theta=\mathrm d\sigma_e\) 的特征空间分解为
\begin{equation}\label{g3:eq:symmetric-lie-algebra-splitting}
	\mathfrak g=\mathfrak k\oplus\mathfrak p,
	\qquad \theta|_{\mathfrak k}=1,\quad \theta|_{\mathfrak p}=-1.
\end{equation}
若在 \(\mathfrak p\) 上给定 \(\operatorname{Ad}(K)\) 不变正定内积，则 \(s_o(gK)=\sigma(g)K\) 良定义，在 \(o\) 固定点的微分为负恒等。由 \(s_oL_g=L_{\sigma(g)}s_o\) 和度量不变性，它是等距；其共轭 \(L_gs_oL_g^{-1}\) 给出其他点的反射。

\ProseParagraphBreak

一般黎曼对称空间也有这种 Lie 齐性描述：由成对点反射生成的连通变换群给出有效的对称对；单连通情形可在相应覆盖群中选取连通不动子群。这里使用这一标准结构结果，详见 \cite[第II章第1--2节、第IV章第1节]{Loos1969SymmetricI}。其作用是把一般对称空间纳入上述构造，以下联络、曲率和具体模型的计算都将展开。改变作用群的有限中心核，并不改变有效作用，也不应被误认为改变了空间的度量。

\subsection{Lie 代数分解与曲率}
\label{g3:subsec:0504:03}

由于 \(\theta\) 保持括号，两个特征空间满足
\begin{equation}\label{g3:eq:symmetric-bracket-relations}
	[\mathfrak k,\mathfrak k]\subseteq\mathfrak k,
	\qquad[\mathfrak k,\mathfrak p]\subseteq\mathfrak p,
	\qquad[\mathfrak p,\mathfrak p]\subseteq\mathfrak k.
\end{equation}
例如最后一式来自 \(\theta[X,Y]=[-X,-Y]=[X,Y]\)。因此取 \(\mathfrak m=\mathfrak p\) 时，式\eqref{g3:eq:homogeneous-levi-civita}的三项全部为零，\(\Lambda=0\)。这不是说曲率为零；\([\mathfrak p,\mathfrak p]\) 仍能通过稳定子的作用旋转切空间。

\begin{proposition}\label{g3:prop:symmetric-curvature-bracket}
	设 \(G/K\) 为上述带不变黎曼度量的对称对，\(o=eK\)。对 \(X,Y,Z\in\mathfrak p\)，有
	\begin{equation}\label{g3:eq:symmetric-curvature-bracket}
		R_o(X,Y)Z=-[[X,Y],Z].
	\end{equation}
	从 \(o\) 出发、初速度为 \(X\) 的测地线为
	\begin{equation}\label{g3:eq:symmetric-geodesic}
		\gamma_X(t)=\exp_G(tX)K.
	\end{equation}
\end{proposition}
\begin{proof}
	曲率式直接由 \(\Lambda=0\) 和 \eqref{g3:eq:homogeneous-curvature} 得到。沿群中的曲线 \(\exp(tX)\)，左 Maurer--Cartan 形式等于 \(X\,\mathrm dt\)，其 \(\mathfrak k\) 分量为零。所以由左平移得到的 \(\mathfrak p\) 标架沿投影曲线平行；曲线速度在此标架中是常向量 \(X\)，协变加速度为零。
\end{proof}

为判断曲率符号，使用伴随表示的迹型
\[
	\mathcal B(X,Y)=\operatorname{tr}(\operatorname{ad}X\circ\operatorname{ad}Y),
\]
即\term[Killing xing]{Killing 型}[Killing form]。\footnote{Wilhelm Karl Joseph Killing（1847--1923），德国数学家，研究非 Euclid 几何及简单 Lie 代数的结构与分类。}
迹的循环性给出 \(\mathcal B([X,Y],Z)=\mathcal B(X,[Y,Z])\)。在紧半单情形，\(-\mathcal B\) 正定；在非紧型的 Cartan 分解中，\(\mathcal B\) 在 \(\mathfrak p\) 上正定，在 \(\mathfrak k\) 上负定，且两者正交。等价地，\(\mathcal B_\theta(X,Y)=-\mathcal B(X,\theta Y)\) 在整个 \(\mathfrak g\) 上正定。后者可作为半单 Lie 代数上\term[Cartan duihe]{Cartan 对合}[Cartan involution]的定义。

\ProseParagraphBreak

分别在 \(\mathfrak p\) 上选取 \(-\mathcal B\) 和 \(\mathcal B\) 为度量，由 \eqref{g3:eq:symmetric-curvature-bracket} 得到
\begin{equation}\label{g3:eq:symmetric-sectional-sign}
	\begin{aligned}
	\langle R(X,Y)Y,X\rangle
	&=-\mathcal B([X,Y],[X,Y])\ge0&&\text{（紧型）},\\
	\langle R(X,Y)Y,X\rangle
	&=\mathcal B([X,Y],[X,Y])\le0&&\text{（非紧型）}.
	\end{aligned}
\end{equation}
两式取零都当且仅当 \([X,Y]=0\)。乘以正常数改变度量时，截面曲率按该常数的倒数变化。对于多个不可约因子的乘积，应分别给各因子选尺度，再用乘积曲率；各因子的尺度需要分别约定。

\subsection{紧型、非紧型和 Euclid 模型}
\label{g3:subsec:0504:04}

\(\mathbb R^m\) 上的 \(s_p(x)=2p-x\) 给出欧氏对称空间。其平移方向全部交换，曲率为零。紧型与非紧型则是去掉欧氏因子后的两种半单情形。一般单连通黎曼对称空间分成欧氏、紧型和非紧型因子，这是上述结构定理与本章 de Rham 分解相结合的结果；“紧型”在这里不包括平坦环面，“非紧型”也不只是底空间非紧的意思。

\begin{example}\label{g3:ex:sphere-hyperbolic-symmetric-models}
	设 \(n\ge2\)。半径 \(a>0\) 的球面为 \(\operatorname{SO}(n+1)/\operatorname{SO}(n)\)：旋转群传递作用，基点 \(ae_0\) 的稳定子恰为作用在后 \(n\) 个坐标上的 \(\operatorname{SO}(n)\)。相应的 \(\mathfrak p\) 由矩阵
	\[
		X_v=\begin{pmatrix}0&-v^{\mathsf T}\\v&0\end{pmatrix},\qquad v\in\mathbb R^n,
	\]
	组成，基点度量是 \(\langle X_v,X_w\rangle=a^2v\cdot w\)。计算两次交换子可得
	\[
		-[[X_v,X_w],X_z]
		=X_{\,v(w\cdot z)-w(v\cdot z)}.
	\]
	因此 \(K=1/a^2\)。点反射是线性映射在球面上的限制：
	\[
		s_p(x)=\frac{2\langle x,p\rangle}{a^2}p-x.
	\]
	它固定 \(p\)，在 \(p^\perp=\mathrm T_pS^n(a)\) 上为负恒等。

	双曲模型取 Minkowski 内积 \(\langle(t,u),(s,v)\rangle_-=-ts+u\cdot v\)，以及双曲面的上片
	\[
		\mathbb H^n_{-1/a^2}
		=\{(t,u):t^2-|u|^2=a^2,\ t>0\}.
	\]
	空间旋转与平面内的双曲旋转可把 \((a,0)\) 送到任意点，稳定子仍为 \(\operatorname{SO}(n)\)，故它是 \(\operatorname{SO}_0(1,n)/\operatorname{SO}(n)\)。此时
	\[
		X_v=\begin{pmatrix}0&v^{\mathsf T}\\v&0\end{pmatrix},\qquad
		-[[X_v,X_w],X_z]=X_{-v(w\cdot z)+w(v\cdot z)}.
	\]
	诱导度量仍在基点给出 \(a^2v\cdot w\)，所以 \(K=-1/a^2\)。点反射为
	\[
		s_p(x)=-\frac{2\langle x,p\rangle_-}{a^2}p-x.
	\]
	在基点它是 \((t,u)\mapsto(t,-u)\)，保持上片；一般点的情形由共轭得到。两模型因而都是全局对称空间，完备性也由 \cref{g3:prop:symmetric-complete-parallel-curvature} 得到。

	这里的半径 \(a\) 是几何尺度。若改用上述半单迹型的标准度量，由 \(\mathcal B_{\mathfrak{so}(n+1)}(X,Y)=(n-1)\operatorname{tr}(XY)\) 及相应实型的同一公式，得到基点内积 \(2(n-1)v\cdot w\)。因此 Killing 型归一化对应 \(a^2=2(n-1)\)，其曲率是 \(\pm1/(2(n-1))\)，并非自动等于 \(\pm1\)。
\end{example}

\begin{example}\label{g3:ex:sl-symmetric-positive-matrices}
	令 \(n\ge2\)，并设
	\[
		\mathcal P_n^1=\{P=P^{\mathsf T}>0:\det P=1\}.
	\]
	\(G=\operatorname{SL}(n,\mathbb R)\) 通过 \(P\mapsto gPg^{\mathsf T}\) 作用。任意 \(P\) 的正定平方根行列式为一，故作用传递；\(I\) 的稳定子为 \(\operatorname{SO}(n)\)。于是
	\begin{equation}\label{g3:eq:sl-spd-identification}
		\operatorname{SL}(n,\mathbb R)/\operatorname{SO}(n)
		\longrightarrow\mathcal P_n^1,
		\qquad gK\longmapsto gg^{\mathsf T}
	\end{equation}
	是微分同胚。其逆由 \(P\mapsto P^{1/2}K\) 给出。

	在正定矩阵模型中取度量
	\begin{equation}\label{g3:eq:spd-trace-metric}
		g_P(U,V)=\frac14\operatorname{tr}(P^{-1}UP^{-1}V),
		\qquad \operatorname{tr}(P^{-1}U)=\operatorname{tr}(P^{-1}V)=0.
	\end{equation}
	共轭变换和迹的循环性保证它不变。对合 \(\sigma(g)=(g^{\mathsf T})^{-1}\) 给出
	\[
		\mathfrak k=\mathfrak{so}(n),\qquad
		\mathfrak p=\{X=X^{\mathsf T}:\operatorname{tr}X=0\}.
	\]
	商映射在单位元的微分为 \(X\mapsto2X\)，故式\eqref{g3:eq:spd-trace-metric}在 \(\mathfrak p\) 上恰对应 \(\operatorname{tr}(XY)\)。系数 \(1/4\) 正是为保持这一个识别。

	还可直接验证联络：若 \(V\) 为矩阵值切向量场，则
	\begin{equation}\label{g3:eq:spd-connection}
		(\nabla_UV)_P=(\mathrm dV)_P[U]
		-\frac12\bigl(UP^{-1}V+VP^{-1}U\bigr).
	\end{equation}
	由 \(\mathrm d(P^{-1})[U]=-P^{-1}UP^{-1}\) 对度量求导，两个逆矩阵的导数分别被上式的两个修正项分担，因而它保持度量；修正项对 \(U,V\) 对称，所以无挠。对切向条件求导又得到
	\[
		\operatorname{tr}(P^{-1}\mathrm dV[U])
		=\operatorname{tr}(P^{-1}UP^{-1}V),
	\]
	故 \(\nabla_UV\) 仍切于 \(\det P=1\)，证明它就是 Levi-Civita 联络。代入可核验所有初值的测地线：
	\begin{equation}\label{g3:eq:spd-geodesic}
		P(t)=P_0^{1/2}\exp\bigl(tP_0^{-1/2}UP_0^{-1/2}\bigr)P_0^{1/2}.
	\end{equation}
	它对所有实数 \(t\) 有定义。矩阵对数把 \(\mathcal P_n^1\) 与无迹对称矩阵空间微分同胚，因此本模型单连通且完备。平方根、指数及对数的光滑性可在正交对角化后由光滑函数演算得到；重复特征值不影响这些矩阵函数。

	点反射是 \(s_Q(P)=QP^{-1}Q\)。在 \(Q=I\) 处，它由矩阵求逆给出，微分为 \(-\mathrm{Id}\)，且直接保持式\eqref{g3:eq:spd-trace-metric}；一般 \(Q\) 由不变作用得到。曲率在 \(\mathfrak p\) 中为
	\begin{equation}\label{g3:eq:sl-trace-sectional}
		K(X\wedge Y)
		=-\frac{\lVert[X,Y]\rVert_{\mathrm F}^2}
		{\operatorname{tr}(X^2)\operatorname{tr}(Y^2)-\operatorname{tr}(XY)^2},
	\end{equation}
	其中 \(\lVert A\rVert_{\mathrm F}^2=\operatorname{tr}(A^{\mathsf T}A)\)。若用正定矩阵的实际切向量 \(U,V,W\) 表示，则在 \(I\) 处同一个公式写为 \(R(U,V)W=-\tfrac14[[U,V],W]\)。这两个系数的差来自 \(U=2X\)。

	当 \(n=2\) 时，\(\mathfrak p\) 是二维空间，取正交单位基
	\[
		X=\frac1{\sqrt2}\begin{pmatrix}1&0\\0&-1\end{pmatrix},\qquad
		Y=\frac1{\sqrt2}\begin{pmatrix}0&1\\1&0\end{pmatrix},
	\]
	便得 \(\lVert[X,Y]\rVert_{\mathrm F}^2=2\)，故曲率恒为 \(-2\)。更具体地，映射
	\[
		(x,y)\longmapsto\frac1y\begin{pmatrix}x^2+y^2&x\\x&1\end{pmatrix},\qquad y>0,
	\]
	把度量拉回为 \((\mathrm dx^2+\mathrm dy^2)/(2y^2)\)。对任意 \(n\)，在矩阵代数上展开 \(\operatorname{ad}X=L_X-R_X\) 可得
	\(\mathcal B(X,Y)=2n\operatorname{tr}(XY)\)（无迹时）；恒等矩阵方向的伴随作用为零，因此取 \(\mathfrak{sl}(n,\mathbb R)\) 的迹仍是此式。Killing 型度量是当前度量的 \(2n\) 倍，所有截面曲率相应再除以 \(2n\)。
\end{example}

\subsection{极大交换子空间与最大平坦}
\label{g3:subsec:0504:05}

在球面上，两个独立切向方向不交换；在 \(\mathcal P_n^1\) 中，所有无迹对角矩阵却彼此交换。由曲率式可见，后一组方向在内部没有曲率。关键还在于这些方向能够积分为全测地子空间。

\begin{definition}\label{g3:def:flat-and-maximal-abelian}
	黎曼流形中的\term[pingtan ziliuxing]{平坦子流形}[flat]指带诱导度量的连通、完备、全测地平坦浸入子流形。若它不真包含于更大的此类子流形，则称为\term[zuidapingtan]{最大平坦}[maximal flat]。对称分解中的交换子空间是满足 \([\mathfrak a,\mathfrak a]=0\) 的线性子空间 \(\mathfrak a\subseteq\mathfrak p\)；不能再扩大者称为极大交换子空间。
\end{definition}

\begin{proposition}\label{g3:prop:symmetric-flats-abelian}
	在紧型或非紧型对称空间的标准度量中，经过基点的平坦子流形，其切空间是 \(\mathfrak p\) 的交换子空间。反过来，交换子空间 \(\mathfrak a\) 的像
	\[
		F_{\mathfrak a}=\{\exp_G(H)K:H\in\mathfrak a\}
	\]
	带有自然的完备全测地平坦浸入结构；极大交换子空间对应最大平坦。在非紧型中它是 Euclid 空间；在紧型中最大平坦是平坦环面。
\end{proposition}
\begin{proof}
	若 \(F\) 全测地且平坦，则它的截面曲率就是环境在其切平面上的截面曲率。式\eqref{g3:eq:symmetric-sectional-sign}因而给出 \([X,Y]=0\) 对 \(X,Y\in\mathrm T_oF\) 成立。

	反过来，\(\mathfrak a\) 交换时，\(\exp(H+X)=\exp(H)\exp(X)\)。因此映射 \(H\mapsto\exp(H)K\) 在每点的微分由左平移识别为 \(\mathfrak a\) 的包含，拉回度量恒定。其直线映为式\eqref{g3:eq:symmetric-geodesic}的左平移，所以像全测地且平坦。令 \(\Gamma=\{H\in\mathfrak a:\exp H\in K\}\)，局部单射性说明 \(\Gamma\) 离散；两个参数有相同像当且仅当其差属于 \(\Gamma\)。故 \(\mathfrak a/\Gamma\) 给出所述完备浸入结构。任何更大的全测地平坦在基点有更大的交换切空间，反向构造亦成立。

	非紧型中，\(\operatorname{ad}H\) 对正定内积 \(\mathcal B_\theta\) 自伴。若 \(\exp H\in K\)，则 \(\exp(\operatorname{ad}H)\) 同时为正定自伴和正交算子，故等于恒等；实特征值进而给出 \(\operatorname{ad}H=0\)。半单性使中心为零，所以 \(H=0\)，即 \(\Gamma=0\)。

	紧型中取 \(A=\exp\mathfrak a\) 在紧群 \(G\) 内的闭包 \(T\)。它是紧连通交换群，且被 \(\sigma\) 保持。其 Lie 代数的 \(-1\) 特征空间仍与 \(\mathfrak a\) 交换；若 \(\mathfrak a\) 极大，该空间便等于 \(\mathfrak a\)。因此 \(A\) 在紧轨道 \(T\cdot o\) 上的轨道维数等于 \(T\cdot o\) 的维数，每条 \(A\) 轨道都是开的。后者连通，故 \(A\cdot o=T\cdot o\)。于是 \(\mathfrak a/\Gamma\) 紧，\(\Gamma\) 为满格，所得平坦是环面。
\end{proof}

\begin{proposition}\label{g3:prop:maximal-abelian-conjugacy}
	在紧型或非紧型的有效对称对中，\(\mathfrak p\) 的极大交换子空间都由 \(\operatorname{Ad}(K)\) 共轭，因而维数相同。每个 \(Y\in\mathfrak p\) 都可被共轭到任一预先选定的极大交换子空间 \(\mathfrak a\) 中。
\end{proposition}
\begin{proof}
	先选 \(H\in\mathfrak a\)，使 \(Z_{\mathfrak p}(H)=\mathfrak a\)。这样的 \(H\) 存在：交换算子族 \(\operatorname{ad}(\mathfrak a)\) 在非紧型中自伴，在紧型中斜自伴，复化后都可同时对角化。其共同零空间与 \(\mathfrak p\) 的交是 \(Z_{\mathfrak p}(\mathfrak a)=\mathfrak a\)，最后的等号来自极大性。避开有限个非零特征线性函数的核，就得到所需 \(H\)。

	对给定 \(Y\in\mathfrak p\)，在紧的稳定子作用群上取函数 \(k\mapsto\mathcal B(\operatorname{Ad}(k)Y,H)\) 的极值。记极值处的向量为 \(Y_0\)。对任意 \(Z\in\mathfrak k\) 求导，有
	\[
		0=\mathcal B([Z,Y_0],H)=\mathcal B(Z,[Y_0,H]).
	\]
	\([Y_0,H]\in\mathfrak k\)，而 \(\mathcal B\) 在 \(\mathfrak k\) 上负定，故 \([Y_0,H]=0\)，即 \(Y_0\in\mathfrak a\)。这里若所用覆盖群中的 \(K\) 不紧，只须在其紧的有效伴随像上取极值。

	再给定另一个极大交换子空间 \(\mathfrak a'\)，取其相应正则元素 \(Y\)，使 \(Z_{\mathfrak p}(Y)=\mathfrak a'\)。把 \(Y\) 共轭到 \(\mathfrak a\) 后，有 \(\mathfrak a\subseteq\operatorname{Ad}(k)\mathfrak a'\)。交换两者再作一次，得到维数相等，于是这一包含为等号。
\end{proof}

\subsection{对称空间的秩及乘积结构}
\label{g3:subsec:0504:06}

最大平坦的公共维数称为对称空间的\term[duicheng kongjian de zhi]{秩}[rank]，记为 \(\operatorname{rank}M\)。\symidx{\(\operatorname{rank}M\)}在纯紧型或纯非紧型中，刚才的命题使它也等于 \(\dim\mathfrak a\)。经过任何点、包含任何给定切向方向，都能找到一个最大平坦；这来自先将该方向扩充为极大交换子空间，再用群作用移动基点。

\ProseParagraphBreak

球面与双曲空间的秩均为 \(1\)：其非零截面曲率排除了二维平坦，而每条完整测地线已是一维平坦。\(\mathcal P_n^1\) 的秩为 \(n-1\)，因为对角无迹矩阵构成 \((n-1)\) 维交换子空间，而任意一族交换实对称矩阵可同时正交对角化。若对角化矩阵行列式为负，只须反转一个特征向量的符号，就可令它属于 \(\operatorname{SO}(n)\)。对应最大平坦为
\[
	\{\operatorname{diag}(e^{2h_1},\ldots,e^{2h_n}):\textstyle\sum_i h_i=0\},
\]
其度量是超平面 \(\sum_i h_i=0\) 上的 \(\sum_i\mathrm dh_i^2\)。它的维数远小于 \(\dim\mathcal P_n^1=n(n+1)/2-1\)。

\begin{proposition}\label{g3:prop:symmetric-rank-product}
	若单连通黎曼对称空间按 de Rham 分解写成
	\[
		M=\mathbb R^d\times M_1\times\cdots\times M_s,
	\]
	其中 \(M_i\) 为非平坦不可约因子，则
	\[
		\operatorname{rank}M=d+\sum_i\operatorname{rank}M_i.
	\]
\end{proposition}
\begin{proof}
	每个因子的最大平坦与 \(\mathbb R^d\) 的乘积给出右端维数的平坦。反过来，设 \(F\) 是经过基点的全测地平坦，\(V=\mathrm T_oF\)。对 \(X,Y\in V\)，全测地性和平坦性给出 \(R(X,Y)X=0\)。投影到第 \(i\) 个因子，得到 \(R_i(X_i,Y_i)X_i=0\)；再与 \(Y_i\) 配对，利用该因子的定号公式，便知 \([X_i,Y_i]=0\)。所以 \(V\) 在每个因子上的投影都是交换子空间，维数不超过该因子的秩；欧氏投影维数不超过 \(d\)。\(V\) 嵌入这些投影的直和，故有反向维数不等式。
\end{proof}

这里使用的整体分解是 \cref{g3:thm:de-rham-decomposition}。平坦因子使秩增加，但秩大于一并不意味着空间可约：\(\operatorname{SL}(n,\mathbb R)/\operatorname{SO}(n)\) 在 \(n\ge3\) 时正是不可约而有高维最大平坦的模型。这里也能直接核验不可约性。若无迹对称矩阵的非零子空间对 \(\operatorname{SO}(n)\) 共轭不变，先对角化其中一个非零矩阵，再对平面旋转求导，就能得到某个 \(E_{ij}+E_{ji}\)；换坐标得到全部非对角对称矩阵，再对旋转求导得到所有对角差，因而该子空间就是整个 \(\mathfrak p\)。稳定子表示不可约，而连通稳定子必须保持 de Rham 的每个非平坦因子，故本模型不可约。

\subsection{紧型与非紧型的对偶}
\label{g3:subsec:0504:07}

把一个对称分解复化，在复 Lie 代数中令
\begin{equation}\label{g3:eq:symmetric-dual-lie-algebra}
	\mathfrak g^d=\mathfrak k\oplus\mathrm i\mathfrak p.
\end{equation}
由括号关系可知它仍是一个实 Lie 代数：\([\mathrm iX,\mathrm iY]=-[X,Y]\)，而 \([Z,\mathrm iX]=\mathrm i[Z,X]\)。保留稳定子方向，把 \(X\mapsto\mathrm iX\) 作为切空间的等距识别，这就是\term[jin feijin duious]{紧非紧对偶}[compact--noncompact duality]。

\begin{proposition}\label{g3:prop:symmetric-duality-curvature}
	设 \((\mathfrak g,\mathfrak k)\) 为紧型或非紧型半单对称对。式\eqref{g3:eq:symmetric-dual-lie-algebra}把它变成相反类型的对称对；以相应 Killing 型度量及切空间等距 \(X\mapsto\mathrm iX\) 比较时，
	\[
		R^d(\mathrm iX,\mathrm iY)\mathrm iZ=-\mathrm iR(X,Y)Z.
	\]
	因此对应切平面的截面曲率互为相反数，两空间的秩相同。
\end{proposition}
\begin{proof}
	对偶实型的 Killing 型是复化迹型的限制，故在 \(\mathrm i\mathfrak p\) 上有 \(\mathcal B^d(\mathrm iX,\mathrm iY)=-\mathcal B(X,Y)\)，在 \(\mathfrak k\) 上不变。它恰把正、负两种切向限制交换。曲率公式给出
	\[
		R^d(\mathrm iX,\mathrm iY)\mathrm iZ
		=-[-[X,Y],\mathrm iZ]=\mathrm i[[X,Y],Z].
	\]
	度量识别保持分母，故截面曲率变号。又 \([X,Y]=0\) 当且仅当 \([\mathrm iX,\mathrm iY]=0\)，极大交换子空间对应，秩不变。
\end{proof}

球面与双曲空间的矩阵块之间的符号变化，正是这一对偶。\(\mathfrak{sl}(n,\mathbb R)=\mathfrak{so}(n)\oplus\operatorname{Sym}_0(n,\mathbb R)\) 的紧对偶是
\[
	\mathfrak{su}(n)=\mathfrak{so}(n)\oplus\mathrm i\operatorname{Sym}_0(n,\mathbb R),
\]
所以 \(\operatorname{SU}(n)/\operatorname{SO}(n)\) 与 \(\operatorname{SL}(n,\mathbb R)/\operatorname{SO}(n)\) 对偶。在 \(\operatorname{SU}(n)\) 中，复共轭的固定子群为 \(\operatorname{SO}(n)\)，从而这是实际的紧对称模型，其秩也为 \(n-1\)。

\ProseParagraphBreak

另一个重要来源是紧单群 \(L\) 的双不变度量。它作为对称空间写为 \((L\times L)/\operatorname{diag}L\)，对合交换两因子，点反射在单位元为 \(x\mapsto x^{-1}\)。若直接在 \(\mathfrak l\) 上采用双不变内积，则 \(\nabla_XY=\tfrac12[X,Y]\)，从定义计算得
\[
	R(X,Y)Z=-\tfrac14[[X,Y],Z].
\]
这里 \(\mathfrak p=\{(X,-X)\}\) 在商映射 \((a,b)\mapsto ab^{-1}\) 下对应 \(2X\)，所以这个 \(1/4\) 与式\eqref{g3:eq:symmetric-curvature-bracket}一致。其非紧对偶是复群商 \(L_{\mathbb C}/L\)。分类时不能漏掉这一族，也不能把 Lie 群的秩与任意商空间的秩混为一谈。

\ProseParagraphBreak

对偶首先是 Lie 代数对和局部度量的对应。取各自的单连通对称模型，才能得到无额外选择的整体对应；任意离散商还带有覆盖群与周期数据。例如 \(S^n\) 与 \(\mathbb{RP}^n\) 局部等距，却不是同一个整体空间；双曲空间的离散商通常只局部对称。两边的基本群和商作用并不由乘上 \(\mathrm i\) 的局部公式决定。\secref{g3:sec:0506}的分类表将始终区分单连通模型与这些整体商。

% 纲要条目排为编号小节标题，并进入目录与书签；说明段仍保留于纲要注释中。

% !TeX root = ../../Book3.tex
% 纲要标题：### 5.5* 特殊完整群的几何内容
% 纲要位置：几何学计划纲要.md，第 2301 行。
% 正文状态：初稿；尚未正式审核。

% 纲要内容开始
% 
% * Kähler、超 Kähler 与平行形式
% * G₂、Spin(7) 的代数模型
% * Berger 完整群分类
% 
% ---
% 
% 纲要内容结束

\OptionalSection{特殊完整群的几何内容}{g3:sec:0505}

一个正交变换群越小，它能够共同保持的张量通常越多。例如，保持一个正交复结构的正交变换组成酉群；同时保持三个满足四元数关系的复结构，群还会进一步缩小。\cref{g3:thm:holonomy-parallel-tensors}使这种切空间中的线性代数成为流形上的几何：被整体完整群固定的张量，恰好延拓为整体平行张量。本节先具体识别这些张量，再说明不可约情形的分类范围。

\subsection{Kähler、超 Kähler 与平行形式}
\label{g3:subsec:0505:01}

设 \(M\) 为实 \(2m\) 维光滑流形。一个丛自同态 \(J:TM\to TM\) 若满足 \(J^2=-I\)，便称为\term[jifujiegou]{几乎复结构}[almost complex structure]；它把每个切空间变成复向量空间，复数单位 \(\mathrm i\) 的作用就是 \(J\)。若黎曼度量还满足 \(g(JX,JY)=g(X,Y)\)，则 \((g,J)\) 为几乎 Hermite 结构\footnote{Charles Hermite（1822--1901），法国数学家，研究椭圆函数与代数方程，并证明 \(e\) 的超越性。}，其基本二形式为
\[
	\omega(X,Y)=g(JX,Y).
\]
由 \(J^{-1}=-J\) 与正交性可得 \(g(JX,Y)=-g(X,JY)\)，所以 \(\omega\) 确为交替形式。相容的实正交基可以写成 \(e_1,Je_1,\ldots,e_m,Je_m\)；此时 \(\omega=\sum_a e^a\wedge(Je_a)^\flat\)，并且 \(\omega^m/m!\) 给出复方向对应的体积形式。

\begin{definition}\label{g3:def:parallel-kahler-structure}
	设 \((M,g)\) 为实 \(2m\) 维光滑黎曼流形。若正交几乎复结构 \(J\) 对 Levi-Civita 联络满足 \(\nabla J=0\)，则称 \((M,g,J)\) 为 \term[kahlerjiegou]{Kähler 结构}[Kähler structure]。本节采用这一联络表述。\footnote{Erich Kähler（1906--2000），德国数学家，研究复几何、微分形式与外微分方程。}
\end{definition}

这一条件同时控制长度、角度及复方向。它与常见的复流形表述相联系，首先是因为 \(\nabla\omega=0\) 蕴含 \(\mathrm d\omega=0\)：无挠联络下，外微分是协变导数的交替化。其次，复结构的可积性障碍
\[
	N_J(X,Y)=[JX,JY]-J[JX,Y]-J[X,JY]-[X,Y]
\]
也为零。确实，将每个括号换成 \(\nabla_XY-\nabla_YX\)，再用 \(\nabla_X(JY)=J\nabla_XY\)，全部项两两抵消。这个计算说明平行性怎样消除可积性障碍；以下关于完整群的论证直接使用 \(\nabla J=0\)。

\begin{proposition}\label{g3:prop:unitary-holonomy-parallel-complex}
	设 \((M,g)\) 为连通、无边界的实 \(2m\) 维光滑黎曼流形，\(p\in M\)。存在 Kähler 结构 \(J\)，当且仅当某个正交识别 \(T_pM\cong\mathbb C^m\) 下有 \(\operatorname{Hol}_p(g)\subseteq\operatorname U(m)\)。指定 \(J_p\) 后，相应平行结构唯一。
\end{proposition}
\begin{proof}
	在实内积空间 \(V\) 上，正交变换 \(A\) 保持 \(J\) 的条件是 \(AJA^{-1}=J\)，也就是 \(AJ=JA\)。把 \(J\) 当作乘以 \(\mathrm i\)，这正说 \(A\) 是复线性的；同时保持实内积便等价于保持对应 Hermite 内积。因此 \(J\) 在正交群中的稳定子为 \(\operatorname U(m)\)。若 \(J\) 平行，完整群固定 \(J_p\)，于是得到包含关系。反之，完整群固定标准的 \(J_p\)，由\cref{g3:thm:holonomy-parallel-tensors}得到唯一平行延拓。平行运输保持 \(g\)，且与张量乘积、收缩相容，故 \(J_p^2=-I\) 和正交性在所有点成立。
\end{proof}

\ProseParagraphBreak

在 \(\mathbb C^m\) 上再给定非零复体积形式 \(\Theta_0=\mathrm dz^1\wedge\cdots\wedge\mathrm dz^m\)，酉变换对它的作用由复行列式决定。因此，保持 \(g,J,\Theta_0\) 的群是 \(\operatorname{SU}(m)\)。这里的复 \(m\) 形式指复值实交替形式，并要求在每个槽放入 \(JX\) 等于将原值乘以 \(\mathrm i\)。这种类型条件与非零性都由平行运输保持。

\begin{corollary}\label{g3:cor:special-unitary-parallel-volume}
	设 \((M,g,J)\) 为连通的 Kähler 流形，复维数为 \(m\)。整体完整群包含在 \(\operatorname{SU}(m)\) 中，当且仅当存在处处非零的平行复 \((m,0)\) 形式。对标准体积形式的一次定长选择可以沿运输保持；其常数相位不影响完整群的包含关系。
\end{corollary}
\begin{proof}
	对上一段的稳定子计算应用\cref{g3:thm:holonomy-parallel-tensors}，分别延拓复形式的实部和虚部。反向则由平行形式在每条闭路运输下保持不变得到。
\end{proof}

\ProseParagraphBreak

四元数向量空间 \(\mathbb H^r\) 还提供一组互相关联的复结构。约定在右四元数模上取 \(I(v)=v\mathrm i\)、\(J(v)=v\mathrm j\)，并取 \(K=IJ\)；由于算子复合次序，\(K(v)=-v\mathrm k\)。这样始终有
\[
	I^2=J^2=K^2=-I_{TM},\qquad IJ=K=-JI.
\]
采用左四元数作用也能得到同一组算子关系；真正需要固定的是这里写出的复合式。

\begin{definition}\label{g3:def:hyperkahler-structure}
	设 \((M,g)\) 为实 \(4r\) 维光滑黎曼流形，\(r\ge1\)，平行性相对于 Levi-Civita 联络。三个正交平行丛自同态 \(I,J,K\) 若满足 \(I^2=J^2=K^2=-I_{TM}\)、\(IJ=K=-JI\)，就称 \((g,I,J,K)\) 为\term[chaokahlerjiegou]{超 Kähler 结构}[hyperkähler structure]。
\end{definition}

\begin{proposition}\label{g3:prop:hyperkahler-holonomy}
	设 \((M,g)\) 为连通、无边界的实 \(4r\) 维光滑黎曼流形，\(r\ge1\)。存在超 Kähler 结构，当且仅当整体完整群在适当正交识别下包含于 \(\operatorname{Sp}(r)\subset\operatorname{SO}(4r)\)。每个超 Kähler 结构给出三个平行闭二形式 \(\omega_I,\omega_J,\omega_K\)。相对于复结构 \(I\)，形式 \(\Psi=\omega_J+\mathrm i\omega_K\) 是处处非退化的平行复 \((2,0)\) 形式。
\end{proposition}
\begin{proof}
	与 \(I,J,K\) 同时交换的实线性映射，正是四元数线性映射；其中保持实内积的群定义为紧辛群 \(\operatorname{Sp}(r)\)。这给出稳定子的识别，再用\cref{g3:thm:holonomy-parallel-tensors}得到前两个条件的等价。算子关系与正交性都是张量恒等式，沿运输保持。三个二形式平行，因此闭。

	验证 \(\Psi\) 的类型只需代入关系 \(JI=-K\)、\(KI=J\)：
	\[
		\Psi(IX,Y)=-\omega_K(X,Y)+\mathrm i\omega_J(X,Y)
		=\mathrm i\Psi(X,Y).
	\]
	另一槽由交替性得到相同结果。若 \(\Psi(X,Y)=0\) 对所有实向量 \(Y\) 成立，取实部可得 \(g(JX,Y)=0\)，因而 \(X=0\)。将实切空间以 \(I\) 看成复空间，这便是复二形式的非退化性。在四元数正交基中，它也可写成 \(r\) 对复坐标的标准辛形式。于是 \(\Psi^r\) 是非零平行复体积形式，特别有 \(\operatorname{Sp}(r)\subset\operatorname{SU}(2r)\)。
\end{proof}

对任意常数 \((a,b,c)\in S^2\)，算子 \(aI+bJ+cK\) 仍然平方为 \(-I_{TM}\)，并且正交、平行；超 Kähler 结构因而给出一个球面族的 Kähler 结构。另一方面，群
\[
	\operatorname{Sp}(r)\operatorname{Sp}(1)
	=\bigl(\operatorname{Sp}(r)\times\operatorname{Sp}(1)\bigr)/\{(1,1),(-1,-1)\}
\]
通过 \(v\mapsto Avq^{-1}\) 作用在 \(\mathbb H^r\) 上。它保持三维空间 \(\operatorname{span}\{I,J,K\}\)，却允许旋转其中的基。实维数 \(4r\ge8\) 时，相应的完整群包含关系给出\term[siyuanshukahlerjiegou]{四元数 Kähler 结构}[quaternionic Kähler structure]；所得平行对象是一整个三维子丛。保持一个子空间与固定它的每个向量，在这里产生了不同的几何。

\subsection{\texorpdfstring{\(G_2\)、Spin(7) 的代数模型}{G₂、Spin(7) 的代数模型}}
\label{g3:subsec:0505:02}

有些小群保持的基本张量是三形式或四形式。先在定向 Euclid 空间 \(\mathbb R^7\) 中取标准正交余标架 \(e^1,\ldots,e^7\)，并简记 \(e^{ijk}=e^i\wedge e^j\wedge e^k\)。定义
\begin{align}
	\varphi_0&=e^{123}+e^{145}+e^{167}+e^{246}-e^{257}-e^{347}-e^{356},\label{g3:eq:standard-g2-three-form}\\
	\psi_0=*\varphi_0&=e^{4567}+e^{2367}+e^{2345}+e^{1357}-e^{1346}-e^{1256}-e^{1247}.\label{g3:eq:standard-g2-four-form}
\end{align}
星号取标准度量和方向 \(e^{1234567}\)。例如 \(*e^{123}=e^{4567}\)，而 \(*e^{246}=e^{1357}\)；其余项按同样的排列符号规则确定。

\begin{proposition}\label{g3:prop:g2-form-recovers-metric}
	在定向 Euclid 空间 \(\mathbb R^7\) 中，记 \(e^1,\ldots,e^7\) 为标准正交余标架，\(e^{1234567}\) 为正体积形式，\(\varphi_0\) 为式\eqref{g3:eq:standard-g2-three-form}的三形式。则它与 Euclid 内积满足
	\begin{equation}\label{g3:eq:g2-metric-recovery}
		(\iota_u\varphi_0)\wedge(\iota_v\varphi_0)\wedge\varphi_0
		=6\langle u,v\rangle e^{1234567}
	\end{equation}
	对所有 \(u,v\in\mathbb R^7\) 成立。因此 \(\varphi_0\) 的每个实线性稳定变换都保持 Euclid 度量及定向。
\end{proposition}
\begin{proof}
	两边对 \(u,v\) 都是对称双线性的，故只需在标准基上计算。由\eqref{g3:eq:standard-g2-three-form}收缩得到
	\[
	\begin{array}{c|l}
		a&\iota_{e_a}\varphi_0\\\hline
		1&e^{23}+e^{45}+e^{67}\\
		2&-e^{13}+e^{46}-e^{57}\\
		3&e^{12}-e^{47}-e^{56}\\
		4&-e^{15}-e^{26}+e^{37}\\
		5&e^{14}+e^{27}+e^{36}\\
		6&-e^{17}+e^{24}-e^{35}\\
		7&e^{16}-e^{25}-e^{34}
	\end{array}
	\]
	每行自乘产生三个不同四元指标项，各出现两次；再乘 \(\varphi_0\) 后，三个非零项都为 \(2e^{1234567}\)。两行不同时，把表中两个二元指标与\eqref{g3:eq:standard-g2-three-form}的三元指标合并，每次或者重复指标，或者缺少所需的互补三元组，所得七形式为零。因此系数矩阵为 \(6I_7\)，得到\eqref{g3:eq:g2-metric-recovery}。

	若 \(A^*\varphi_0=\varphi_0\)，将恒等式应用于 \(Au,Av\) 并拉回，可得
	\[
		\langle Au,Av\rangle\det A=\langle u,v\rangle.
	\]
	正定性先给出 \(\det A>0\)。对矩阵等式 \((\det A)A^{\mathsf T}A=I\) 取行列式，得到 \((\det A)^9=1\)，所以 \(\det A=1\)、\(A^{\mathsf T}A=I\)。
\end{proof}

这个稳定子是紧例外 Lie 群 \(G_2\)，维数为14。八维的对应模型是在 \(\mathbb R e_0\oplus\mathbb R^7\) 上取\term[cayley sixingshi]{Cayley 四形式}[Cayley four-form]\footnote{Arthur Cayley（1821--1895），英国数学家，研究矩阵代数、不变量与高维几何。}
\begin{equation}\label{g3:eq:standard-spin7-cayley-form}
	\Phi_0=e^0\wedge\varphi_0+\psi_0.
\end{equation}
其稳定子是 \(\operatorname{Spin}(7)\) 的八维实旋量表示像，是 \(\operatorname{SO}(8)\) 中维数21的紧群。这两个稳定子的群识别属于例外群的代数输入；模型、维数及识别可见 Joyce 的正式综述第2.2--2.3节\cite[第2.2--2.3节，第112--113页]{Joyce2006ExceptionalHolonomy}。式\eqref{g3:eq:standard-spin7-cayley-form}按 \(e^0,e^1,\ldots,e^7\) 排列，与该文八维模型只差指标起点。

\ProseParagraphBreak

在七维流形上，一个处处可由线性同构化成 \(\varphi_0\) 的三形式称为\term[zhengsanxingshi]{正三形式}[positive three-form]。上一命题说明它确定度量和方向：在每点选这样的标架，两种选择之间的变换保持标准内积，故所得度量与选择无关。若此度量的 Levi-Civita 联络还满足 \(\nabla\varphi=0\)，便得到无挠 \(G_2\) 结构。类似地，点点等价于 \(\Phi_0\) 的四形式称为 \(\operatorname{Spin}(7)\) 型四形式；要求它对诱导度量平行，得到无挠 \(\operatorname{Spin}(7)\) 结构。

\begin{proposition}\label{g3:prop:exceptional-parallel-forms}
	设 \((M,g)\) 为连通光滑黎曼流形。以式\eqref{g3:eq:standard-g2-three-form}和式\eqref{g3:eq:standard-spin7-cayley-form}的形式 \(\varphi_0,\Phi_0\) 定义标准稳定子群 \(G_2\subset\operatorname{SO}(7)\)、\(\operatorname{Spin}(7)\subset\operatorname{SO}(8)\)；相容形式指在每点的某正交标架中等于相应标准形式，\(*\) 是所诱导定向的 Hodge 星算子。在七维情形，存在与 \(g\) 相容的平行正三形式，当且仅当整体完整群包含在标准 \(G_2\) 中；在八维情形，存在相容的平行 \(\operatorname{Spin}(7)\) 型四形式，当且仅当整体完整群包含在此标准 \(\operatorname{Spin}(7)\) 中。平行的 \(\varphi\)、\(*\varphi\) 和 \(\Phi\) 均为闭形式。
\end{proposition}
\begin{proof}
	对稳定子中的模型形式应用\cref{g3:thm:holonomy-parallel-tensors}。运输是线性等距，故标准型及相容度量沿路径保持。反向由每条环路都固定平行形式得到。Levi-Civita 联络保持度量和方向，因而与 Hodge 星算子交换；外微分又是无挠协变导数的交替化，所以这些形式都闭。
\end{proof}

一个可直接核验的例子来自三复维 Kähler 流形 \((N,h,J)\) 及定长平行复体积形式 \(\Theta\)，将其长度归一化为标准酉余标架中 \(\mathrm dz^1\wedge\mathrm dz^2\wedge\mathrm dz^3\) 的长度。在 \(\mathbb R\times N\) 上取乘积度量，并令
\[
	\varphi=\mathrm dt\wedge\omega+\operatorname{Re}\Theta.
\]
在 \(e^1=\mathrm dt\)、\(z^1=x^2+\mathrm ix^3\)、\(z^2=x^4+\mathrm ix^5\)、\(z^3=x^6+\mathrm ix^7\) 对应的酉余标架中展开，正好得到\eqref{g3:eq:standard-g2-three-form}。三个因子均平行，所以 \(\varphi\) 平行。但 \(\partial_t\) 也是平行向量场，完整群实际包含于固定这一方向的 \(\operatorname{SU}(3)\)，而不会等于 \(G_2\)。这把“包含于某群”和“恰好等于某群”的区别落实到一个平行方向上。

\subsection{Berger 完整群分类}
\label{g3:subsec:0505:03}

de Rham 分解先把单连通完备流形分为 Euclid 因子与不可约因子；局部对称的不可约因子由下一节讨论的对称空间分类处理。余下的非局部对称不可约情形，具有一个有限的完整群名单。以下引用 \term[bergerfenlei]{Berger 分类}[Berger classification] 的黎曼版本。\footnote{Marcel Berger（1927--2016），法国数学家，研究黎曼完整群、曲率与整体几何。}其陈述可核对 Joyce 第2.1节定理2.3\cite[第2.1节，定理2.3，第111--112页]{Joyce2006ExceptionalHolonomy}；一般分类的证明作为本节外部输入。

\termidx[Berger fenlei]{Berger 完整群分类}
\foreignidx{Berger holonomy classification}
\begin{theorem}[Berger 分类]\label{g3:thm:berger-classification}
	设 \((M,g)\) 为连通、单连通、无边界的光滑黎曼流形，维数 \(n\ge2\)，其 Levi-Civita 完整群在 \(T_pM\) 上的实表示不可约，并且 \(\nabla R\) 不恒为零。则在适当正交基下，\(\operatorname{Hol}_p(g)\) 是下表中的一个群，采用表中标准的实正交表示。这里不要求完备性；当从整体 de Rham 因子进入本定理时，完备性已由该分解的假设提供。
\end{theorem}

\begin{table}[htbp]
	\centering
	\caption{非局部对称不可约黎曼完整群及其平行结构}
	\label{g3:tab:berger-holonomy}
	\begin{tabularx}{\textwidth}{@{}lllX@{}}
		\toprule
		群 & 实维数 & 参数范围 & 由群识别的几何结构 \\
		\midrule
		\(\operatorname{SO}(n)\) & \(n\) & \(n\ge2\) & 度量及定向 \\
		\(\operatorname U(m)\) & \(2m\) & \(m\ge2\) & 平行正交复结构 \\
		\(\operatorname{SU}(m)\) & \(2m\) & \(m\ge2\) & 平行复结构及复体积形式 \\
		\(\operatorname{Sp}(r)\) & \(4r\) & \(r\ge2\) & 平行四元数三元组 \\
		\(\operatorname{Sp}(r)\operatorname{Sp}(1)\) & \(4r\) & \(r\ge2\) & 平行四元数结构子丛 \\
		\(G_2\) & \(7\) & & 平行正三形式 \\
		\(\operatorname{Spin}(7)\) & \(8\) & & 平行 Cayley 四形式 \\
		\bottomrule
	\end{tabularx}
\end{table}

表中省去了低维重复：\(\operatorname U(1)=\operatorname{SO}(2)\)，\(\operatorname{Sp}(1)=\operatorname{SU}(2)\)，而 \(\operatorname{Sp}(1)\operatorname{Sp}(1)\) 在四维表示中的像为 \(\operatorname{SO}(4)\)。特殊酉群在 \(m=1\) 时为平凡群，产生的平坦二维表示可约，因而不属于这里的不可约非局部对称范围。

\ProseParagraphBreak

几项假设各有具体作用。单连通性使整体与受限完整群一致；在一般流形上，分类先适用于万有覆盖上的相应不可约因子，基本群还可能增加离散分支。不可约性排除了直接乘积，例如两个曲面的乘积通常具有 \(\operatorname{SO}(2)\times\operatorname{SO}(2)\) 完整群。\(\nabla R\ne0\) 的条件把局部对称情形交给另一套分类；对称空间可以具有名单以外的不可约各向同性表示。

\ProseParagraphBreak

本节已将平行张量与群包含关系对应起来。分类定理进一步限制哪些群可以恰好出现为非局部对称不可约黎曼完整群；构造达到等号的度量还有独立的存在问题。例如平坦七维环面带有平行 \(\varphi_0\)，但其完整群是平凡群。关于 \(\operatorname{SU}(m)\)、\(\operatorname{Sp}(r)\)、\(G_2\)、\(\operatorname{Spin}(7)\) 完整群的 Ricci 平坦性，以及例外完整群的实际存在，可进一步查阅同一综述第2.1--2.3节及第3--4节；这些曲率限制与存在结果在此按该文外引。
% 纲要条目排为编号小节标题，并进入目录与书签；说明段仍保留于纲要注释中。

% !TeX root = ../../Book3.tex
% 纲要标题：### 5.6*　限制根与对称空间分类
% 纲要位置：几何学计划纲要.md，第 2309 行。

% 纲要内容开始
% 
% > 写作范围与外部输入：规划任务：在第5.4节的对称分解与对偶模型之后建立限制根语言，完整计算一个经典矩阵商的秩、限制根和重数。给出不可约单连通紧型与非紧型模型的分类表及准确适用条件；一般 Cartan 分类作为外部输入，不以第一卷第19章的有限根算例冒充完整分类证明。所用根系与李群覆盖结果按实际版本就近恢复。
% 
% * 限制根、根空间与重数
% * 特殊线性群商的限制根
% * 不可约经典模型与低秩同构
% * 例外对称空间与 Cartan 分类
% * 局部等距、单连通覆盖与整体商
% * 分类表与紧非紧对偶
% 
% ---
% 
% 纲要内容结束
% 正文状态：初稿；尚未正式审核。
% 规划任务：在第5.4节的对称分解与对偶模型之后建立限制根语言，完整计算一个经典矩阵商的秩、限制根和重数。给出不可约单连通紧型与非紧型模型的分类表及准确适用条件；一般 Cartan 分类作为外部输入，不以第一卷第19章的有限根算例冒充完整分类证明。所用根系与李群覆盖结果按实际版本就近恢复。

\OptionalSection{限制根与对称空间分类}{g3:sec:0506}

最大平坦给出一组彼此交换的方向，但同样的秩仍可能对应不同的对称空间。进一步观察这些方向怎样作用在其他切向和稳定子方向上，便得到限制根及其重数。它们同时记录平坦的维数、横向方向的数量和曲率算子的特征值。下面先建立这一语言，完整计算特殊线性群商，再把已算的模型放入一般分类；分类的穷尽性使用正式教材中的 Cartan 分类，而不从少数矩阵算例推断。

\subsection{限制根、根空间与重数}
\label{g3:subsec:0506:01}

先在非紧型有效半单对称对 \(\mathfrak g=\mathfrak k\oplus\mathfrak p\) 中工作。令 \(\theta\) 为 Cartan 对合，\(\mathcal B\) 为 Killing 型，\(\mathfrak a\subseteq\mathfrak p\) 为极大交换子空间。在正定内积 \(\mathcal B_\theta(X,Y)=-\mathcal B(X,\theta Y)\) 下，迹型的不变性给出
\begin{equation}\label{g3:eq:cartan-adjoint-identity}
	(\operatorname{ad}X)^*=-\operatorname{ad}(\theta X).
\end{equation}
特别地，每个 \(\operatorname{ad}H\)、\(H\in\mathfrak a\) 都自伴，且这些算子彼此交换，故存在共同的实特征空间分解。

\begin{definition}\label{g3:def:restricted-root}
	设 \(\mathfrak g=\mathfrak k\oplus\mathfrak p\) 为非紧型有效半单对称 Lie 代数，\(\theta\) 为 Cartan 对合，\(\mathcal B\) 为 Killing 型，\(\mathfrak a\subset\mathfrak p\) 为极大交换子空间。对 \(\alpha\in\mathfrak a^*\)，令
	\[
		\mathfrak g_\alpha
		=\{X\in\mathfrak g:[H,X]=\alpha(H)X\text{ 对所有 }H\in\mathfrak a\}.
	\]
	使 \(\mathfrak g_\alpha\ne0\) 的非零线性函数称为\term[xianzhi gen]{限制根}[restricted root]，其集合记为 \(\Sigma\)；\(\mathfrak g_\alpha\) 称为限制根空间，\(m_\alpha=\dim_{\mathbb R}\mathfrak g_\alpha\) 称为\term[gen de chongshu]{根的重数}[root multiplicity]。
\end{definition}
\symidx{\(\Sigma,\mathfrak g_\alpha,m_\alpha\)：限制根、实根空间与重数}

这里的根空间在实 Lie 代数中取维数，与复半单 Lie 代数的复一维根空间不能混同。称为“限制”，是因为把 \(\mathfrak a\) 扩充到复化 Lie 代数的 Cartan 子代数后，非零限制根正是其普通根在 \(\mathfrak a\) 上的非零限制；不同普通根可以限制成同一个函数。以下构造只需共同特征空间，不以这个扩充解释作为证明前提。

\begin{proposition}\label{g3:prop:restricted-root-decomposition}
	设 \(\mathfrak g=\mathfrak k\oplus\mathfrak p\) 为非紧型有效半单对称 Lie 代数，\(\theta\) 为 Cartan 对合，\(\mathcal B\) 为 Killing 型，\(\mathfrak a\subset\mathfrak p\) 为极大交换子空间。记 \(\Sigma\)、\(\mathfrak g_\alpha\)、\(m_\alpha\) 分别为限制根集、相应根空间和重数，并以 \(G/K\) 表示该对称 Lie 代数的局部齐性模型。令 \(\mathfrak m=Z_{\mathfrak k}(\mathfrak a)\)。则
	\begin{equation}\label{g3:eq:restricted-root-decomposition}
		\mathfrak g=\mathfrak m\oplus\mathfrak a
		\oplus\bigoplus_{\alpha\in\Sigma}\mathfrak g_\alpha,
		\qquad \mathfrak g_0=\mathfrak m\oplus\mathfrak a.
	\end{equation}
	此外，\(\theta\mathfrak g_\alpha=\mathfrak g_{-\alpha}\)，且
	\[
		[\mathfrak g_\alpha,\mathfrak g_\beta]\subseteq\mathfrak g_{\alpha+\beta},
	\]
	其中不存在的特征空间理解为零。选取 \(H_0\in\mathfrak a\) 使所有根在其上非零，并令 \(\Sigma^+=\{\alpha:\alpha(H_0)>0\}\)，则
	\begin{equation}\label{g3:eq:restricted-root-dimension}
		\dim(G/K)=\dim\mathfrak a+\sum_{\alpha\in\Sigma^+}m_\alpha.
	\end{equation}
\end{proposition}
\begin{proof}
	共同自伴算子的谱分解给出 \(\mathfrak g=\bigoplus_\alpha\mathfrak g_\alpha\)。零空间是 \(Z_{\mathfrak g}(\mathfrak a)\)；它被 \(\theta\) 保持，其 \(\mathfrak p\) 分量恰为 \(\mathfrak a\)，否则可扩大交换子空间，\(\mathfrak k\) 分量则为 \(\mathfrak m\)。若 \(X\in\mathfrak g_\alpha\)，则
	\[
		[H,\theta X]=-\theta[H,X]=-\alpha(H)\theta X,
	\]
	故 \(\theta\) 交换正负根空间，特别有 \(m_\alpha=m_{-\alpha}\)。Jacobi 恒等式给出
	\[
		[H,[X,Y]]=[[H,X],Y]+[X,[H,Y]]
		=(\alpha(H)+\beta(H))[X,Y],
	\]
	证明括号关系。最后，每对 \(\mathfrak g_\alpha\oplus\mathfrak g_{-\alpha}\) 的 \(-1\) 特征空间为
	\[
		\mathfrak p_\alpha=\{X-\theta X:X\in\mathfrak g_\alpha\},
	\]
	其维数为 \(m_\alpha\)，与 \(\mathfrak a\) 一起直和为 \(\mathfrak p\)，于是得到维数公式。
\end{proof}

\begin{proposition}\label{g3:prop:restricted-root-system}
	设 \(\mathfrak g=\mathfrak k\oplus\mathfrak p\) 为非紧型有效半单对称 Lie 代数，\(\theta\) 为 Cartan 对合，\(\mathcal B\) 为 Killing 型，\(\mathfrak a\subset\mathfrak p\) 为极大交换子空间。记 \(\Sigma\) 为限制根集、\(m_\alpha\) 为根重数。以 \(\mathcal B|_{\mathfrak a}\) 识别 \(\mathfrak a\) 与其对偶，\(\Sigma\) 是张成 \(\mathfrak a^*\) 的有限根系，允许非约化。具体地，它保持于反射
	\[
		s_\alpha(\beta)=\beta-
		\frac{2\langle\beta,\alpha\rangle}{\langle\alpha,\alpha\rangle}\alpha,
	\]
	并满足 \(2\langle\beta,\alpha\rangle/\langle\alpha,\alpha\rangle\in\mathbb Z\)。若 \(t\alpha\in\Sigma\)、\(t>0\)，则 \(t\in\{1/2,1,2\}\)。反射保持根重数。
\end{proposition}
\begin{proof}
	若 \(H\in\mathfrak a\) 被所有根消去，则它与式\eqref{g3:eq:restricted-root-decomposition}中的每个空间交换，因而在半单代数的中心中，只能为零；故根张成对偶空间。

	记 \(\alpha^\sharp\) 为 \(\alpha\) 对应的向量。取 \(0\ne E\in\mathfrak g_\alpha\)，令 \(F=-\theta E\)。\([E,F]\) 属于 \(\mathfrak a\)，且由不变性
	\[
		\mathcal B([E,F],H)=\alpha(H)\mathcal B_\theta(E,E).
	\]
	缩放 \(E\) 使 \(\mathcal B_\theta(E,E)=2/\lVert\alpha\rVert^2\)，便得
	\[
		H_\alpha=[E,F]=\frac{2\alpha^\sharp}{\lVert\alpha\rVert^2},\qquad
		[H_\alpha,E]=2E,\quad[H_\alpha,F]=-2F.
	\]
	在伴随表示中，\(\operatorname{ad}E\) 与 \(\operatorname{ad}F\) 互为伴随，\(\operatorname{ad}H_\alpha\) 自伴。第一卷第19.4节的有限升降链论证在这里适用；恢复其要点如下。若 \(v\) 是最高权向量，\(Hv=\lambda v\)、\(Ev=0\)，则交换关系递推给出
	\[
		E F^jv=j(\lambda-j+1)F^{j-1}v.
	\]
	有限维性使链在某个 \(F^{N+1}v=0\)、\(F^Nv\ne0\) 处终止，取 \(j=N+1\) 得 \(\lambda=N\)。正交补保持三个算子，再重复便得到所有特征值为整数。因此 \(\operatorname{ad}H_\alpha\) 在 \(\mathfrak g_\beta\) 上的特征值 \(\beta(H_\alpha)\) 是整数。

	\(E-F\in\mathfrak k\)。由三维括号关系，\(\operatorname{Ad}(\exp(\tfrac\pi2(E-F)))\) 把 \(H_\alpha\) 送到 \(-H_\alpha\)，并固定 \(\ker\alpha\subseteq\mathfrak a\)，故在 \(\mathfrak a\) 上实现反射。它将 \(\mathfrak g_\beta\) 同构地送到反射后的根空间，证明根集和重数均被保持。若 \(t\alpha\) 也是根，对两种次序应用整性得到 \(2t\in\mathbb Z\)、\(2/t\in\mathbb Z\)，于是 \(t=1/2,1,2\)。
\end{proof}

若根集中没有与 \(\alpha\) 同时出现的 \(2\alpha\)，称为约化；否则称为非约化。典型非约化根系为
\[
	BC_r=\{\pm e_i,\ \pm2e_i,\ \pm e_i\pm e_j\ (i<j)\}.
\]
同一直线上的短根与倍根须分别记录重数。例如后面的 FII 型有 \(m_\alpha=8\)、\(m_{2\alpha}=7\)，所以秩一却有维数 \(1+8+7=16\)。

\ProseParagraphBreak

这些根也直接进入曲率。对 \(H\in\mathfrak a\)、\(V=X-\theta X\in\mathfrak p_\alpha\)，有 \((\operatorname{ad}H)^2V=\alpha(H)^2V\)，从而
\begin{equation}\label{g3:eq:restricted-root-curvature}
	R(V,H)H=-\alpha(H)^2V.
\end{equation}
紧对偶把这个符号反转，而根重数与秩不变。它说明限制根记录的是沿平坦方向看到的横向曲率特征值。

\subsection{特殊线性群商的限制根}
\label{g3:subsec:0506:02}

回到 \(\operatorname{SL}(n,\mathbb R)/\operatorname{SO}(n)\)，\(n\ge2\)，取
\[
	\mathfrak a=\{H=\operatorname{diag}(h_1,\ldots,h_n):\textstyle\sum_i h_i=0\}.
\]
前节已用实对称矩阵的同时对角化证明它极大，故秩为 \(n-1\)。令 \(E_{ij}\) 为矩阵单位，直接相乘得
\[
	[H,E_{ij}]=(h_i-h_j)E_{ij}.
\]
因此全部限制根、根空间和重数为
\begin{equation}\label{g3:eq:sl-restricted-roots}
	\Sigma=\{\alpha_{ij}:i\ne j\},\quad
	\alpha_{ij}(H)=h_i-h_j,\quad
	\mathfrak g_{\alpha_{ij}}=\mathbb R E_{ij},\quad m_{\alpha_{ij}}=1.
\end{equation}
零特征空间是无迹对角矩阵，故 \(\mathfrak m=0\)。这就是 \(A_{n-1}\) 根系。在当前的迹度量 \(\operatorname{tr}(XY)\) 下，\(\alpha_{ij}\) 的对偶向量为 \(E_{ii}-E_{jj}\)，平方长度为 \(2\)；若改用 \(\mathcal B=2n\operatorname{tr}(XY)\)，根的平方长度为 \(1/n\)。

\ProseParagraphBreak

选 \(h_1>h_2>\cdots>h_n\)，正根为 \(\alpha_{ij}\)、\(i<j\)，单根可取 \(\alpha_{i,i+1}\)。所有正根都满足
\[
	\alpha_{ij}=\alpha_{i,i+1}+\cdots+\alpha_{j-1,j}.
\]
故正根数为 \(n(n-1)/2\)，维数公式逐项给出
\[
	(n-1)+\frac{n(n-1)}2=\frac{n(n+1)}2-1
	=\dim\operatorname{SL}(n,\mathbb R)-\dim\operatorname{SO}(n).
\]
与其配对的切向根空间和稳定子方向分别由
\[
	S_{ij}=\frac{E_{ij}+E_{ji}}{\sqrt2},\qquad
	A_{ij}=\frac{E_{ij}-E_{ji}}{\sqrt2}
\]
张成，满足 \([H,S_{ij}]=(h_i-h_j)A_{ij}\)、\([H,A_{ij}]=(h_i-h_j)S_{ij}\)。于是式\eqref{g3:eq:restricted-root-curvature}在这里成为一个逐个二维块即可核验的矩阵等式。

\begin{proposition}\label{g3:prop:sl-restricted-weyl-group}
	设 \(n\ge2\)，在 \(\operatorname{SL}(n,\mathbb R)/\operatorname{SO}(n)\) 中取 \(\mathfrak a=\{\operatorname{diag}(h_1,\ldots,h_n):\sum_i h_i=0\}\)。则\term[xianzhi Weyl qun]{限制 Weyl 群}[restricted Weyl group]
	\[
		W=N_{\operatorname{SO}(n)}(\mathfrak a)/Z_{\operatorname{SO}(n)}(\mathfrak a)
	\]
	为置换群 \(S_n\)，其作用是置换对角元。根超平面为 \(h_i=h_j\)，每个室由对角元的一种严格次序给出。
\end{proposition}
\begin{proof}
	若一个正交矩阵正规化全部对角无迹矩阵，它必须把某个具有互异特征值的对角矩阵的特征直线集合映到坐标直线集合，所以它是带符号的置换矩阵。反过来，这样的矩阵都正规化 \(\mathfrak a\)。每种置换都可通过调整一个符号，使代表的行列式为一；逐点固定 \(\mathfrak a\) 的恰是行列式为一的对角符号矩阵。取商后只剩置换，得到 \(S_n\)。交换第 \(i,j\) 个坐标正是根 \(\alpha_{ij}\) 的反射。
\end{proof}

\(\operatorname{SU}(n)/\operatorname{SO}(n)\) 的最大交换切空间为 \(\mathrm i\mathfrak a\)。对偶保留上述根的数目与重数，把横向曲率的符号反转。两边的最大平坦分别是环面与 \(\mathbb R^{n-1}\)；相同的 \(A_{n-1}\) 数据没有抹掉这一整体差别。

\subsection{不可约经典模型与低秩同构}
\label{g3:subsec:0506:03}

下面各表都在正维、不可约、单连通的黎曼对称模型范围内使用；度量在每个不可约因子上允许乘一个正常数。第一张表的两列模型分别为 \(U/K\) 和 \(G/K\)，\(U\) 表示紧群，\(G\) 表示相应非紧实型。表中的 \(K\) 是实际矩阵嵌入的子群，\(\operatorname S(\operatorname U(p)\times\operatorname U(q))\) 指两块行列式乘积为一的子群。对正交群商使用表列的连通稳定子，而非其可能更大的非连通扩张。

\ProseParagraphBreak

记 \(\operatorname{Sp}(n)\) 为保持标准 四元数 内积的紧群；\(\operatorname{Sp}(2n,\mathbb R)\) 和 \(\operatorname{Sp}(2n,\mathbb C)\) 则分别是作用在 \(\mathbb R^{2n}\)、\(\mathbb C^{2n}\) 上的辛矩阵群。\(\operatorname{SU}^*(2n)\) 是 \(\operatorname{SL}(2n,\mathbb C)\) 中与标准 四元数 结构交换的实型，\(\operatorname{SO}^*(2n)\) 是相应的正交实型；\(\operatorname{Sp}(p,q)\) 保持签名 \((p,q)\) 的 四元数 Hermitian 型。下标 \(0\) 表示恒等连通分支。这样的维数记号避免把 \(\operatorname{Sp}(n)\) 与实辛群的参数混用。

\begingroup
\small
\setlength{\tabcolsep}{3pt}
\begin{longtable}{@{}p{.10\textwidth}p{.14\textwidth}p{.16\textwidth}p{.23\textwidth}p{.19\textwidth}p{.08\textwidth}@{}}
	\caption{七个经典对称空间族：紧模型 \(U/K\) 与非紧模型 \(G/K\)}\label{g3:tab:classical-symmetric-spaces}\\
	\toprule
	类型、范围 & \(U\) & \(G\) & \(K\) & 实维数 & 秩\\
	\midrule
	\endfirsthead
	\multicolumn{6}{l}{\thetable\ （续）}\\
	\toprule
	类型、范围 & \(U\) & \(G\) & \(K\) & 实维数 & 秩\\
	\midrule
	\endhead
	\bottomrule
	\endfoot
	AI\newline \(n\ge2\) & \(\operatorname{SU}(n)\) & \(\operatorname{SL}(n,\mathbb R)\) & \(\operatorname{SO}(n)\) & \(\tfrac12(n-1)(n+2)\) & \(n-1\)\\[3pt]
	AII\newline \(n\ge2\) & \(\operatorname{SU}(2n)\) & \(\operatorname{SU}^*(2n)\) & \(\operatorname{Sp}(n)\) & \((n-1)(2n+1)\) & \(n-1\)\\[3pt]
	AIII\newline \(1\le p\le q\) & \(\operatorname{SU}(p+q)\) & \(\operatorname{SU}(p,q)\) & \(\operatorname S(\operatorname U(p)\times\operatorname U(q))\) & \(2pq\) & \(p\)\\[3pt]
	BDI\newline 见下文 & \(\operatorname{SO}(p+q)\) & \(\operatorname{SO}_0(p,q)\) & \(\operatorname{SO}(p)\times\operatorname{SO}(q)\) & \(pq\) & \(p\)\\[3pt]
	CI\newline \(n\ge1\) & \(\operatorname{Sp}(n)\) & \(\operatorname{Sp}(2n,\mathbb R)\) & \(\operatorname U(n)\) & \(n(n+1)\) & \(n\)\\[3pt]
	CII\newline \(1\le p\le q\) & \(\operatorname{Sp}(p+q)\) & \(\operatorname{Sp}(p,q)\) & \(\operatorname{Sp}(p)\times\operatorname{Sp}(q)\) & \(4pq\) & \(p\)\\[3pt]
	DIII\newline \(n\ge2\) & \(\operatorname{SO}(2n)\) & \(\operatorname{SO}^*(2n)\) & \(\operatorname U(n)\) & \(n(n-1)\) & \(\lfloor n/2\rfloor\)\\
\end{longtable}
\endgroup

BDI 的参数为 \(1\le p\le q\)、\(p+q\ge3\)，并排除 \((p,q)=(2,2)\)。被排除的商是 \(S^2\times S^2\) 及其非紧对偶的局部类型，属于可约情形；\((1,1)\) 是一维平坦情形，已经由 \(p+q\ge3\) 排除。第一张表来自 \cite[第VII章第2节，第100--120页，尤其第119页表4]{Loos1969SymmetricII}，把原书分列的奇、偶正交族统一为 BDI。

\ProseParagraphBreak

这些范围保留了便于理解的低维重合，所以表不是不重复的计数表。把尺度相应调整后，有
\[
	\begin{aligned}
	\operatorname{SU}(2)/\operatorname{SO}(2)
	&\simeq\operatorname{Sp}(1)/\operatorname U(1)\simeq S^2,\\
	\operatorname{SU}(4)/\operatorname{Sp}(2)&\simeq S^5,\qquad
	\operatorname{Sp}(2)/(\operatorname{Sp}(1)\times\operatorname{Sp}(1))\simeq S^4,\\
	\operatorname{SO}(6)/\operatorname U(3)&\simeq\mathbb{CP}^3,\qquad
	\operatorname{SO}(4)/\operatorname U(2)\simeq S^2.
	\end{aligned}
\]
更一般地，AIII 中 \(p=1\) 给复射影空间，CII 中 \(p=1\) 给 四元数 射影空间，BDI 中 \(p=1\) 给球面。正交与辛族的低维交叉还包括 \(\operatorname{Sp}(2)/\operatorname U(2)\) 与 \(\operatorname{SO}(5)/(\operatorname{SO}(2)\times\operatorname{SO}(3))\)。这些重合由 \(B_2=C_2\)、\(D_3=A_3\) 等 Lie 代数同构解释，原书第120页列出相应低维对照。\(D_2=A_1\oplus A_1\) 则解释上述可约例外。它们说的是正确尺度下的单连通模型，不能直接把任意同名矩阵群的有限商也视为相同。

\ProseParagraphBreak

还有一整族不能并入这七行：紧单群 \(L\) 本身配双不变度量，及其非紧对偶 \(L_{\mathbb C}/L\)，通常分别称为 II 型与 IV 型。这里 \(L\) 与复化群 \(L_{\mathbb C}\) 都取单连通形式。全部可能的简单类型列于下表；\(B_2\) 与 \(C_2\) 的重复保留并明确注明。普通复根在这种非紧商中成为实二维限制根空间，所以每个限制根的重数为 \(2\)，而非 \(1\)。这一点也可由把复根空间忘掉复标量、视为实空间直接看出。

\begingroup
\small
\setlength{\tabcolsep}{3pt}
\begin{longtable}{@{}p{.13\textwidth}p{.23\textwidth}p{.28\textwidth}p{.16\textwidth}p{.08\textwidth}@{}}
	\caption{紧简单群型与复简单群型：\(L\) 对偶于 \(L_{\mathbb C}/L\)}\label{g3:tab:group-type-symmetric-spaces}\\
	\toprule
	类型、范围 & \(L\) & \(L_{\mathbb C}\) & 实维数 & 秩\\
	\midrule
	\endfirsthead
	\multicolumn{5}{l}{\thetable\ （续）}\\
	\toprule
	类型、范围 & \(L\) & \(L_{\mathbb C}\) & 实维数 & 秩\\
	\midrule
	\endhead
	\bottomrule
	\endfoot
	\(A_\ell,\ \ell\ge1\) & \(\operatorname{SU}(\ell+1)\) & \(\operatorname{SL}(\ell+1,\mathbb C)\) & \(\ell(\ell+2)\) & \(\ell\)\\[3pt]
	\(B_\ell,\ \ell\ge2\) & \(\operatorname{Spin}(2\ell+1)\) & \(\operatorname{Spin}(2\ell+1,\mathbb C)\) & \(\ell(2\ell+1)\) & \(\ell\)\\[3pt]
	\(C_\ell,\ \ell\ge2\) & \(\operatorname{Sp}(\ell)\) & \(\operatorname{Sp}(2\ell,\mathbb C)\) & \(\ell(2\ell+1)\) & \(\ell\)\\[3pt]
	\(D_\ell,\ \ell\ge4\) & \(\operatorname{Spin}(2\ell)\) & \(\operatorname{Spin}(2\ell,\mathbb C)\) & \(\ell(2\ell-1)\) & \(\ell\)\\[3pt]
	\(G_2\) & \(G_2\) & \(G_2(\mathbb C)\) & \(14\) & \(2\)\\
	\(F_4\) & \(F_4\) & \(F_4(\mathbb C)\) & \(52\) & \(4\)\\
	\(E_6\) & \(E_6\) & \(E_6(\mathbb C)\) & \(78\) & \(6\)\\
	\(E_7\) & \(E_7\) & \(E_7(\mathbb C)\) & \(133\) & \(7\)\\
	\(E_8\) & \(E_8\) & \(E_8(\mathbb C)\) & \(248\) & \(8\)\\
\end{longtable}
\endgroup

此表采用简单 Lie 代数的分类及其单连通紧形式，参见 \cite[第V章第4节、第VII章第1节]{Loos1969SymmetricII}。\(\operatorname{Spin}(3)\simeq\operatorname{Sp}(1)\simeq\operatorname{SU}(2)\)、\(\operatorname{Spin}(5)\simeq\operatorname{Sp}(2)\)、\(\operatorname{Spin}(6)\simeq\operatorname{SU}(4)\) 解释未另列的低阶类型；\(\operatorname{Spin}(4)\) 是两个 \(\operatorname{SU}(2)\) 的乘积，故不列入简单群表。紧群 \(\operatorname{SU}(2)\) 的适当双不变度量又是圆球面 \(S^3\)，与第一张表的秩一球面行重合。

\subsection{例外对称空间与 Cartan 分类}
\label{g3:subsec:0506:04}

最后还有十二个例外的非群型对称对。为把整体含义说清楚，下表中 \(U\) 均取单连通紧群，\(\mathfrak k\) 指表列对合的固定 Lie 代数，\(K_U=(U^\sigma)_0\)，紧模型为 \(U/K_U\)。非紧列给出对偶实型的标准记号，模型取相应非紧对称空间的单连通形式。各向同性群由这个对称对确定；表中的直和是 Lie 代数直和，不能把它擅自换成无中心商的群直积。例如 \(\mathfrak{so}(10)\oplus\mathfrak u(1)\) 并没有声称实际子群是 \(\operatorname{Spin}(10)\times\operatorname U(1)\)。

\ProseParagraphBreak

记号 \(E_{6(6)}\)、\(E_{6(-26)}\) 等括号内的数，表示相应实 Lie 代数 Killing 型的签名差 \(\dim\mathfrak p-\dim\mathfrak k\)，不是对称空间的维数或秩。表中同时列限制根系，方便与维数、秩核对。

\begingroup
\small
\setlength{\tabcolsep}{3pt}
\begin{longtable}{@{}p{.07\textwidth}p{.07\textwidth}p{.15\textwidth}p{.29\textwidth}p{.08\textwidth}p{.06\textwidth}p{.09\textwidth}@{}}
	\caption{十二个例外对称对及其紧非紧对偶}\label{g3:tab:exceptional-symmetric-spaces}\\
	\toprule
	类型 & \(U\) & 非紧实型 & \(\mathfrak k\) & 维数 & 秩 & \(\Sigma\)\\
	\midrule
	\endfirsthead
	\multicolumn{7}{l}{\thetable\ （续）}\\
	\toprule
	类型 & \(U\) & 非紧实型 & \(\mathfrak k\) & 维数 & 秩 & \(\Sigma\)\\
	\midrule
	\endhead
	\bottomrule
	\endfoot
	EI & \(E_6\) & \(E_{6(6)}\) & \(\mathfrak{sp}(4)\) & \(42\) & \(6\) & \(E_6\)\\
	EII & \(E_6\) & \(E_{6(2)}\) & \(\mathfrak{su}(6)\oplus\mathfrak{sp}(1)\) & \(40\) & \(4\) & \(F_4\)\\
	EIII & \(E_6\) & \(E_{6(-14)}\) & \(\mathfrak{so}(10)\oplus\mathfrak u(1)\) & \(32\) & \(2\) & \(BC_2\)\\
	EIV & \(E_6\) & \(E_{6(-26)}\) & \(\mathfrak f_4\) & \(26\) & \(2\) & \(A_2\)\\
	EV & \(E_7\) & \(E_{7(7)}\) & \(\mathfrak{su}(8)\) & \(70\) & \(7\) & \(E_7\)\\
	EVI & \(E_7\) & \(E_{7(-5)}\) & \(\mathfrak{so}(12)\oplus\mathfrak{sp}(1)\) & \(64\) & \(4\) & \(F_4\)\\
	EVII & \(E_7\) & \(E_{7(-25)}\) & \(\mathfrak e_6\oplus\mathfrak u(1)\) & \(54\) & \(3\) & \(C_3\)\\
	EVIII & \(E_8\) & \(E_{8(8)}\) & \(\mathfrak{so}(16)\) & \(128\) & \(8\) & \(E_8\)\\
	EIX & \(E_8\) & \(E_{8(-24)}\) & \(\mathfrak e_7\oplus\mathfrak{sp}(1)\) & \(112\) & \(4\) & \(F_4\)\\
	FI & \(F_4\) & \(F_{4(4)}\) & \(\mathfrak{sp}(3)\oplus\mathfrak{sp}(1)\) & \(28\) & \(4\) & \(F_4\)\\
	FII & \(F_4\) & \(F_{4(-20)}\) & \(\mathfrak{so}(9)\) & \(16\) & \(1\) & \(BC_1\)\\
	G & \(G_2\) & \(G_{2(2)}\) & \(\mathfrak{so}(4)\) & \(8\) & \(2\) & \(G_2\)\\
\end{longtable}
\endgroup

例外固定子代数的依据为 \cite[第VII章第3节，第124页表5、第130--132页及表6]{Loos1969SymmetricII}；限制根、秩与维数的汇总见同书第146页表8。FII 的紧模型是 \(F_4/\operatorname{Spin}(9)\)，也称八元数射影平面；G 型的紧模型是 \(G_2/\operatorname{SO}(4)\)。这里仅用这两个名称识别表中空间，不以八元数坐标构造作为分类证明。

\termidx[Cartan duicheng kongjian fenlei]{Cartan 对称空间分类}
\foreignidx{Cartan classification of symmetric spaces}
\begin{theorem}[Cartan 分类]\label{g3:thm:cartan-symmetric-classification}
	每个正维、不可约、单连通、非平坦的黎曼对称空间，按正常数倍的度量等价，属于\ref{g3:tab:classical-symmetric-spaces}、\ref{g3:tab:group-type-symmetric-spaces}或\ref{g3:tab:exceptional-symmetric-spaces}中的紧模型或非紧模型。反过来，各表在所注明参数范围内均给出这样的模型，允许已说明的低维重复。一般单连通黎曼对称空间是 Euclid 空间与有限个这些不可约因子的黎曼乘积，各非平坦因子的尺度独立选择。
\end{theorem}

此定理的穷尽性，以及例外 Lie 代数和对合的分类，采用 \cite[第V章第4节、第VII章第1--3节]{Loos1969SymmetricII} 的一般分类。第一卷第19章提供了根空间、反射和部分低秩根图，本节又完整算出 AI 型；这些计算核验了表的一部分，没有承担所有根图和所有对合的穷尽证明。乘积部分还使用本章已证的 de Rham 分解与对称空间的紧非紧结构。

\subsection{局部等距、单连通覆盖与整体商}
\label{g3:subsec:0506:05}

表中的 Lie 代数对首先确定局部模型。若使用紧侧的单连通群 \(U\) 与连通固定子群 \(K_U\)，商 \(U/K_U\) 单连通：商映射具有局部截面，可提升路径；一条商中的环路提升后的端点差在 \(K_U\) 中，利用 \(K_U\) 连通可把它补成 \(U\) 中的环路，再由 \(U\) 单连通缩去。这是第一卷覆盖与齐性空间结论在此处的具体用法。经典表中即使写的是 \(\operatorname{SO}\) 而非 \(\operatorname{Spin}\)，所列商也已经是相应的单连通空间：\(K\) 中的平面转动映到 \(\pi_1(\operatorname{SO}(m))\) 的生成元，因此提升到 \(\operatorname{Spin}(m)\) 后不再产生额外覆盖。被排除的低维情形已经单独处理。

\ProseParagraphBreak

一般连通完备局部对称空间的万有覆盖是单连通对称模型的乘积；这是对称空间的局部到整体结构结果，见 \cite[第II章第3节、第IV章第1节]{Loos1969SymmetricI}。原空间还需由一个自由、固有不连续的离散等距群 \(\Gamma\) 取商。表格没有记录 \(\Gamma\)，因此没有分类所有基本群、闭测地线长度或商空间。特别地，把一个 Lie 群替换为有限中心商，同时把稳定子取为相应像，可能根本不改变齐性空间；而扩大稳定子的离散分支，则可能真正改变空间的覆盖类型。

\begin{proposition}\label{g3:prop:symmetric-quotient-reflection-criterion}
	设 \(\widetilde M\) 为单连通黎曼对称空间，\(\Gamma\) 自由、固有不连续地作用为等距变换。以 \(s_x\) 表示 \(\widetilde M\) 在 \(x\) 处的点反射，\(\pi\) 为商投影。令 \(M=\Gamma\backslash\widetilde M\)，\(x\in\widetilde M\)。在 \(\pi(x)\) 处存在全局点反射，当且仅当
	\[
		s_x\Gamma s_x^{-1}=\Gamma.
	\]
	所以商空间全局对称，当且仅当这个条件对每个 \(x\) 都成立。
\end{proposition}
\begin{proof}
	若正规化条件成立，则 \(s_x\) 下降为商上的等距变换，固定 \(\pi(x)\) 并在其切空间上为负恒等，故是点反射。反过来，商上的点反射可提升为固定 \(x\) 的等距变换；其微分也为负恒等，所以由一点一阶数据的唯一性，它正是 \(s_x\)。覆盖空间的提升等距变换共轭保持覆盖变换群，因此得到正规化条件。
\end{proof}

\(S^n\) 的对径变换与所有点反射交换，故 \(\mathbb{RP}^n\) 全局对称，它与球面的差别是覆盖数据。在 \(\mathbb H^2\) 上，取沿一条轴的非平凡双曲平移 \(\gamma\)，则 \(\langle\gamma\rangle\backslash\mathbb H^2\) 完备且局部对称；若 \(x\) 不在这条轴上，\(s_x\) 把轴送到另一条测地线，因此不能共轭保持同一个循环群。这个商在该点没有全局点反射。于是“完备局部对称”与“全局对称”即使在二维也必须分开。

\subsection{分类表与紧非紧对偶}
\label{g3:subsec:0506:06}

阅读分类表时，可以从三个可直接计算的量检查所选行。首先，\(\dim M=\dim\mathfrak g-\dim\mathfrak k\)，在对偶两边保持不变。例如 EIII 的维数为 \(78-(45+1)=32\)，EVI 为 \(133-(66+3)=64\)。其次，秩来自 \(\mathfrak p\) 的极大交换子空间，或限制根系所在空间的维数，不是紧群 \(U\) 的秩。最后，根重数必须与式\eqref{g3:eq:restricted-root-dimension}相容。

\ProseParagraphBreak

例如 EIV 的限制根系为 \(A_2\)，三个正根的重数都为 \(8\)，所以维数为 \(2+3\cdot8=26\)。FII 的 \(BC_1\) 则给出 \(1+8+7=16\)。同是秩一，实球面、复射影空间、四元数 射影空间和 FII 的根重数不同，横向曲率的特征值及其出现次数也不同。这些数值与表中的一般分类相符，核验并不把根重数的实质来源替换成维数凑数。

\ProseParagraphBreak

对偶变换保持秩、维数、根系与重数，改变的是 \([\mathfrak p,\mathfrak p]\) 的符号及由此产生的曲率符号。若要比较数值，还须使用同一尺度：本节 \(\operatorname{SL}(n,\mathbb R)/\operatorname{SO}(n)\) 的迹度量与 Killing 型度量相差 \(2n\)，而球面和双曲模型中的半径另有几何归一化。分类表确定不可约模型及其尺度自由度；一个具体写出的度量，仍应通过联络、括号或模型坐标把这个常数核实。


\begin{chaptersummary}
	完整群把沿闭路的平行运输汇集为基点纤维上的变换群。受限完整群是其单位元连通分支；一般度量向量丛的完整群可以不是正交群中的闭子群，因而构造时须保留浸入 Lie 群的自身拓扑。Ambrose--Singer 定理把它的 Lie 代数写成所有搬运曲率算子的线性张成：局部曲率沿不同路径进入同一个基点切空间，才共同决定连续的返回方向。

	\ProseParagraphBreak

	整体完整群固定的张量唯一延拓为整体平行张量。不变正交子空间给出平行分布，而无挠性与度量相容性使互补分布形成局部黎曼乘积。完备性保证沿叶和道路矩形继续构造，单连通性消除最后的覆盖歧义，于是得到整体 de Rham 分解。非平坦不可约因子由度量决定，平坦方向合并为一个 Euclid 因子；一般流形还要保留覆盖变换对这些因子的作用。

	\ProseParagraphBreak

	对称空间的点反射把切空间与稳定子方向分成 \(\mathfrak p\) 和 \(\mathfrak k\)，曲率由双重 Lie 括号表达，平坦方向由 \(\mathfrak p\) 内的交换子空间描述。最大平坦的维数是秩，紧非紧对偶改变 \([\mathfrak p,\mathfrak p]\) 的符号。特殊线性群商把这些概念化为对称矩阵、对角化和矩阵单位的计算；限制根的重数则进一步记录各根方向的维数。特殊完整群与对称空间的分类处理不同的条件，两者都须区分模型的局部结构、单连通实现以及离散商。
\end{chaptersummary}

\BookExercises

\begin{exercise}\label{g3:ex:surface-holonomy}
	设 \((M,g)\) 为连通无边界的二维光滑黎曼流形，\(K\) 为其截面曲率。
	\begin{enumerate}
		\item 证明：若 \(K\not\equiv0\)，则在任意基点选正交标架后，受限完整群是 \(\mathrm{SO}(2)\)。若 \(K\equiv0\)，受限完整群是什么？
		\item 在 \(K\not\equiv0\) 的条件下，证明整体完整群在 \(M\) 可定向时是 \(\mathrm{SO}(2)\)，在 \(M\) 不可定向时是 \(\mathrm O(2)\)。
	\end{enumerate}
\end{exercise}
\begin{solution}
	1. 二维中，\(\mathfrak{so}(\mathrm T_pM)\) 是一维实向量空间。若某点 \(q\) 的 \(K(q)\ne0\)，则对 \(q\) 处的正交基 \(e_1,e_2\)，算子 \(R_q(e_1,e_2)\ne0\)。取路径 \(\gamma:p\to q\)，其共轭 \(P_\gamma^{-1}R_q(e_1,e_2)P_\gamma\) 仍非零。由\cref{g3:thm:ambrose-singer}，完整代数是整个 \(\mathfrak{so}(2)\)。受限完整群是对应的连通浸入子群，所以等于 \(\mathrm{SO}(2)\)。若 \(K\equiv0\)，二维曲率张量为零，由\cref{g3:cor:flat-restricted-holonomy}得受限完整群为 \(\{I\}\)。

	2. 若 \(M\) 已定向，沿路径选取与该定向一致的局部正交标架。每一小段运输矩阵连续地从单位矩阵出发，行列式为 \(+1\)；相邻正定向标架的过渡也在 \(\mathrm{SO}(2)\) 中。因此环路运输都保定向，\(\operatorname{Hol}_p\subseteq\mathrm{SO}(2)\)。结合第一问的反向包含，得到相等。

	若 \(M\) 不可定向，必有一个环路的运输反定向。否则，在 \(\mathrm T_pM\) 选定一个定向，沿任意路径把它送到各点；两条路径所得定向的差由一个环路运输给出，因而假设保证定义无关路径。在小坐标球内，用一条固定引路加上随终点光滑变化的径向路径，这些定向连续变化，于是得到整个 \(M\) 的定向，矛盾。因此整体完整群包含一个行列式为 \(-1\) 的元素 \(h\)。它已包含 \(\mathrm{SO}(2)\)，故也包含另一个陪集 \(h\mathrm{SO}(2)\)，两者之并是 \(\mathrm O(2)\)。
\end{solution}

\begin{exercise}\label{g3:ex:smooth-flat-curvature-jet}
	在 \(\mathbb R^2\) 的平凡秩二正交丛上，令
	\[
		f(x)=\begin{cases}e^{-1/x^2},&x>0,\\0,&x\leq0,\end{cases}
		\qquad J=\begin{pmatrix}0&-1\\1&0\end{pmatrix},
		\qquad \nabla^E=\mathrm d+Jf(x)\,\mathrm dy.
	\]
	证明该联络光滑，且曲率在原点的全部阶坐标偏导都为零，但其整体与受限完整群都等于 \(\mathrm{SO}(2)\)。进一步证明：把联络限制到以原点为中心的任意开圆盘，所得完整群仍为 \(\mathrm{SO}(2)\)。这个例子对“只用基点曲率的全部阶导数计算光滑联络的完整代数”说明了什么？
\end{exercise}
\begin{solution}
	对 \(x>0\)，\(f\) 的任意阶导数都是 \(x^{-1}\) 的一个多项式乘以 \(e^{-1/x^2}\)。对任意 \(N\geq0\)，令 \(u=1/x\) 即得
	\[
		\lim_{x\downarrow0}x^{-N}e^{-1/x^2}
		=\lim_{u\to\infty}u^Ne^{-u^2}=0.
	\]
	逐阶求导并比较左右极限，得到 \(f\in C^\infty(\mathbb R)\) 且 \(f^{(j)}(0)=0\) 对所有 \(j\) 成立。因此联络光滑。联络矩阵只有 \(\mathrm dy\) 分量，故
	\[
		R^E=Jf'(x)\,\mathrm dx\wedge\mathrm dy.
	\]
	它的全部阶坐标偏导在原点都为零。若还在底流形上使用通常的平直联络来定义曲率的迭代协变导数，结论也不变：丛上的交换子项为 \([Jf(x),Jf'(x)]=0\)，以后各阶同样为零，故这些协变导数在原点仍全为零。

	另一方面，\(x>0\) 时 \(f'(x)=2x^{-3}e^{-1/x^2}>0\)，所以任何以原点为中心的开圆盘中都有非零曲率。所有曲率值及其平行共轭均为 \(J\) 的实倍数；至少有一个非零倍数，故由\cref{g3:thm:ambrose-singer}完整代数为 \(\mathbb RJ=\mathfrak{so}(2)\)。平面和开圆盘都单连通，整体完整群等于受限完整群，也就都等于 \(\mathrm{SO}(2)\)。因此在光滑范畴内，即使基点的曲率全部阶导数已知，也不能一般地把它们的 Lie 代数张成代替沿实际路径搬运的曲率；本例中前者为零，后者是一维。
\end{solution}

\begin{exercise}\label{g3:ex:constant-connection-full-so3}
	在 \(\mathbb R^2\times\mathbb R^3\) 上取标准纤维内积，令
	\[
		A=\begin{pmatrix}0&-1&0\\1&0&0\\0&0&0\end{pmatrix},
		\qquad B=\begin{pmatrix}0&0&0\\0&0&-1\\0&1&0\end{pmatrix},
		\qquad\nabla^E=\mathrm d+A\,\mathrm dx+B\,\mathrm dy.
	\]
	计算曲率，并证明完整群是 \(\mathrm{SO}(3)\)。解答中应具体给出三个线性无关的搬回曲率算子，说明为什么每一点的曲率算子只张成一维，仍可能产生三维的完整代数。
\end{exercise}
\begin{solution}
	\(A,B\) 均斜对称，所以联络保持内积。系数为常矩阵，曲率来自交换子项：
	\[
		R^E=C\,\mathrm dx\wedge\mathrm dy,
		\qquad C=[A,B]
		=\begin{pmatrix}0&0&1\\0&0&0\\-1&0&0\end{pmatrix}.
	\]
	直接相乘可核对 \([A,C]=-B\)、\([B,C]=A\)。因此
	\[
		\exp(sA)C\exp(-sA)=C\cos s-B\sin s,
		\qquad
		\exp(tB)C\exp(-tB)=C\cos t+A\sin t.
	\]
	例如第一式的左端对 \(s\) 求导等于与 \(A\) 的交换子，右端满足同一线性方程和同一初值，故两边相等；第二式同理。

	从原点沿 \(x\) 轴到 \((s,0)\) 的运输为 \(\exp(-sA)\)，所以在那里取 \(R^E(\partial_x,\partial_y)=C\) 并搬回，得到 \(C\cos s-B\sin s\)。沿 \(y\) 轴到 \((0,t)\) 则得到 \(C\cos t+A\sin t\)。取常值路径、\(s=\pi/2\) 和 \(t=\pi/2\)，分别得到 \(C,-B,A\)。这三个矩阵线性无关并张成 \(\mathfrak{so}(3)\)，所以由\cref{g3:thm:ambrose-singer}完整代数为 \(\mathfrak{so}(3)\)。

	底流形单连通，整体完整群连通。联络矩阵斜对称且迹为零，在全局正交标架中每条运输从 \(I\) 连续出发，故所有运输都在 \(\mathrm{SO}(3)\) 内。对应 \(\mathfrak{so}(3)\) 的连通浸入子群是整个 \(\mathrm{SO}(3)\)，于是完整群确为 \(\mathrm{SO}(3)\)。每一点的二形式 \(R^E\) 只有矩阵方向 \(C\)，但运输会把这个方向共轭成其他方向；完整代数正是同时记录这些比较后的方向。
\end{solution}

\begin{exercise}\label{g3:ex:parallel-symmetric-tensors-de-rham}
	设连通、单连通、完备黎曼流形具有 de Rham 分解
	\[
		M=\mathbb R^k\times M_1\times\cdots\times M_s,
		\qquad
		g=g_0\oplus g_1\oplus\cdots\oplus g_s,
	\]
	其中 \(M_i\) 非平坦且完整群不可约。证明每个平行对称协变二阶张量都唯一写成
	\[
		h=h_0+\sum_{i=1}^s c_i g_i,
	\]
	这里各项均通过因子投影拉回，\(h_0\) 是 \(\mathbb R^k\) 上的常系数对称双线性形式，\(c_i\in\mathbb R\)。求这些张量构成的向量空间的维数，并说明何时 \(h\) 是黎曼度量；在这种情形证明它与 \(g\) 有同一个 Levi-Civita 联络，且仍完备。
\end{exercise}
\begin{solution}
	用 \(g\) 将 \(h\) 写成 \(h(X,Y)=g(AX,Y)\)。对称性等价于 \(A\) 自伴，由 \(\nabla g=0\) 知 \(\nabla h=0\) 等价于 \(\nabla A=0\)。在基点切空间分解
	\(V=V_0\oplus V_1\oplus\cdots\oplus V_s\) 上，\(A\) 与完整群的每个元素交换。由\cref{g3:thm:de-rham-decomposition}证明中的\eqref{g3:eq:product-holonomy}，第 \(i\) 个完整群 \(G_i\) 可独立作用于 \(V_i\)，并对其余因子恒等。

	固定 \(i\ge1\)。对 \(v\in V_i^\perp\)，其 \(V_i\) 分量 \(P_iAv\) 被 \(G_i\) 固定，因为
	\[
		b(P_iAv)=P_iA(bv)=P_iAv\qquad(b\in G_i).
	\]
	\(V_i\) 上的不可约非平凡表示没有非零固定向量，故 \(P_iAv=0\)。再用 \(A\) 自伴可得反方向的混合块也为零。因此 \(A\) 按上述正交分解成为块对角自同态。每个自伴块 \(A|_{V_i}\) 的特征空间为 \(G_i\) 不变子空间，不可约性迫使它只有一个特征值，即 \(A|_{V_i}=c_iI\)。而 \(V_0\) 上群作用平凡，自伴块任意。

	平行张量由一个点的值唯一决定，所以这些块给出题中的整体公式；反过来，公式中的每一项都平行。由此空间维数为
	\[
		\frac{k(k+1)}2+s.
	\]
	\(h\) 正定当且仅当 \(h_0\) 正定且所有 \(c_i>0\)，零维 Euclid 块不附加条件。此时原联络无挠且满足 \(\nabla h=0\)，由 Levi-Civita 联络的唯一性，它也是 \(h\) 的联络。因此两度量有相同的仿射参数测地线；\(g\) 测地完备使 \(h\) 测地完备，再由 Hopf--Rinow 定理，\(h\) 完备。也可取各有限维块的最小、最大特征值，得到常数 \(0<c\le C\) 使 \(cg\le h\le Cg\)，直接比较两种距离。
\end{solution}

\begin{exercise}\label{g3:ex:incomplete-disk-parallel-product}
	设 \(D=\{(x,y)\in\mathbb R^2:x^2+y^2<1\}\) 取 Euclid 度量。证明 \(D\) 单连通，且两条坐标直线分布为正交平行分布，但 \(D\) 不是两个正维黎曼流形的整体黎曼乘积。具体取 \(1/\sqrt2<a<1\)，考察从原点出发的道路
	\[
		\alpha(s)=(as,0),\qquad \beta(t)=(0,at),
		\qquad 0\le s,t\le1,
	\]
	说明平行矩形的整体延拓在何处失败。
\end{exercise}
\begin{solution}
	收缩 \((u,z)\mapsto(1-u)z\) 保持在凸集 \(D\) 中，故 \(D\) 单连通。Euclid 联络在直角坐标下的系数为零，两条常值直线分布显然正交且平行。

	若 \(\alpha,\beta\) 存在平行矩形 \(Q(s,t)=(q_1(s,t),q_2(s,t))\)，适应分布的条件给出
	\[
		\partial_tq_1=0,\qquad \partial_sq_2=0.
	\]
	结合两条给定边，必有 \(Q(s,t)=(as,at)\)。但
	\(|Q(1,1)|=\sqrt2a>1\)，不属于 \(D\)，所以完整矩形不存在。沿参数对角线 \(s=t\) 填充时，在 \(s=t=1/(\sqrt2a)\) 已到达被删去的圆周。

	若 \(D\) 是尊重这两个分布的整体乘积，则分别在两因子中行走 \(\alpha,\beta\)，其配对必给出整个平行矩形，与上面矛盾。实际上任一非平凡黎曼乘积在二维都给出两条正交平行直线分布；Euclid 平行移动保持常向量，所以这种平行分布必须由一对固定正交方向给出。绕原点旋转坐标后，仍可用完全相同的两条道路得到矛盾，故不存在任何非平凡的整体黎曼乘积分解。

	最后，径向单位速测地线在有限时间到达圆周而不能继续留在 \(D\)，故 \(D\) 不完备。在\cref{g3:lem:parallel-rectangle-extension}的证明中，包含这些道路的闭有界集不再保证紧，靠近圆周时可用乘积盒的统一大小趋于零，这正是本例缺失的条件。
\end{solution}

\begin{exercise}\label{g3:ex:irrational-flat-torus-parallel-leaf}
	固定无理数 \(\alpha\)，令
	\[
		\Lambda=\mathbb Z(1,0)+\mathbb Z(\alpha,1),
		\qquad T=\mathbb R^2/\Lambda,
	\]
	赋予由 Euclid 度量下降的平坦度量。证明 \(T\) 紧且完备，整体完整群平凡，水平与竖直分布都是平行分布。证明每条竖直积分叶在 \(T\) 中稠密，因而这两个指定分布不能成为整体黎曼乘积的两个因子方向。解释它与完整群不变子空间对应定理并不矛盾。
\end{exercise}
\begin{solution}
	格点平移自由且适当不连续，故商是光滑黎曼曲面，投影为局部等距覆盖。闭平行四边形
	\[
		\{u(1,0)+v(\alpha,1):0\le u,v\le1\}
	\]
	紧且映满 \(T\)，所以 \(T\) 紧，因而完备。常场 \(\partial_x,\partial_y\) 在所有格点平移下保持不变，下降为整体平行标准正交标架；完整群固定这组基，故平凡。它们各自张成的分布也平行。

	用商坐标
	\[
		[u,v]=[x-\alpha y,y]\in\mathbb R^2/\mathbb Z^2
	\]
	识别底层环面。经过原点的竖直线变为
	\[
		t\longmapsto[-\alpha t,t].
	\]
	固定目标第二坐标 \(b\)，令 \(t=m+b\)、\(m\in\mathbb Z\)，第二坐标便恰为 \([b]\)，第一坐标为 \([-\alpha m-\alpha b]\)。无理旋转的整数轨道在圆周上稠密：给定 \(\varepsilon>0\)，将足够多个 \([j\alpha]\) 放入长度小于 \(\varepsilon\) 的有限分箱中，抽屉原理给出非零整数 \(q\) 和整数 \(l\)，使 \(0<|q\alpha-l|<\varepsilon\)。换 \(q\) 的符号后可令 \(\delta=q\alpha-l>0\)。点 \(0,\delta,\ldots,\lfloor1/\delta\rfloor\delta\) 与圆周任一点的距离都小于 \(\varepsilon\)，且这些点属于同一个整数旋转轨道。这证明所需稠密性。因此竖直线可任意逼近每个 \([u,b]\)，其像在整个环面稠密；平移给出其他竖直叶同样稠密。

	这条叶并非整个环面。例如在固定第二坐标 \([b]\) 的圆周上，它只取可数多个第一坐标值，而该圆周不可数。若两个指定分布来自整体乘积，竖直叶应是 \(\{x\}\times N\) 的像，因而是闭嵌入子流形，不可能是一条真且稠密的叶。

	完整群不变子空间对应的结论确实给出了整体平行分布，也给出了局部乘积；整体 de Rham 定理还要求单连通。这里每条竖直叶以自身的叶拓扑是一条完备直线，其到 \(T\) 的浸入不为嵌入。这同时说明不能在全局乘积证明开始时就把积分叶视为闭子流形。
\end{solution}

\begin{exercise}\label{g3:ex:parallel-one-form-ricci}
	设 \((M,g)\) 为连通、无边界的光滑黎曼流形，\(\alpha\) 为平行实一形式，\(V=\alpha^\sharp\)。证明 \(V\) 平行、\(\mathrm d\alpha=0\)，且 \(R(X,Y)V=0\)、\(\operatorname{Ric}(V,W)=0\) 对所有切向量成立。由此证明：若某点 Ricci 二次型正定，则 \(\alpha=0\)。若 \(M\) 完备且单连通而 \(\alpha\ne0\)，说明 \(M\) 有一个等距的实直线乘积因子。最后证明在紧 \(M\) 上，一个恰当的平行一形式必为零。
\end{exercise}
\begin{solution}
	音乐同构与度量联络相容，故 \(\nabla V=0\)。无挠性给出
	\[
		\mathrm d\alpha(X,Y)=(\nabla_X\alpha)(Y)-(\nabla_Y\alpha)(X)=0.
	\]
	把平行场代入曲率定义得到 \(R(X,Y)V=0\)。由\cref{g3:prop:curvature-symmetries}的两对交换律及末两槽斜对称性，
	\[
		g(R(X,V)Y,Z)=g(R(Z,Y)V,X)=0,
	\]
	于是 \(R(X,V)=0\)。取正交基收缩即得 \(\operatorname{Ric}(V,W)=0\)。在 Ricci 正定点，\(\operatorname{Ric}(V,V)=0\) 迫使 \(V=0\)；平行场长度恒定，因此 \(V\) 处处为零。

	若 \(V\ne0\)，其单位化生成一个平行直线丛，正交补也平行。由\cref{g3:thm:global-parallel-product}，完备单连通的 \(M\) 分裂出一个完备单连通的一维因子。沿其单位平行场取弧长坐标，该因子就是 \(\mathbb R\)。若 \(M\) 紧且 \(\alpha=\mathrm df\)，实函数 \(f\) 在最大值点满足 \(\mathrm df=0\)。\(\alpha\) 的长度恒定，故 \(\alpha=0\)。
\end{solution}

\begin{exercise}\label{g3:ex:spd-distance-formula}
	在 \(\mathcal P_n^1\) 上使用式\eqref{g3:eq:spd-trace-metric}的度量。证明任意 \(P,Q\in\mathcal P_n^1\) 之间存在唯一的极小测地线段（固定仿射参数），并求出距离
	\[
		d(P,Q)=\frac12\left\lVert\log(P^{-1/2}QP^{-1/2})\right\rVert_{\mathrm F}.
	\]
	若 \(P,Q\) 可交换，把答案写成它们在共同正交特征基中的特征值表达式。
\end{exercise}
\begin{solution}
	合同变换 \(S\mapsto P^{-1/2}SP^{-1/2}\) 是等距，故只需连接 \(I\) 与 \(Q_0=P^{-1/2}QP^{-1/2}\)。从 \(I\) 出发的测地线全为 \(t\mapsto\exp(tU)\)，其中 \(U\) 为无迹实对称矩阵。正定矩阵有唯一的实对称对数，因此在 \([0,1]\) 上到达 \(Q_0\) 的唯一测地线是 \(\exp(t\log Q_0)\)。本模型完备，Hopf--Rinow 定理保证存在极小段，故它必是这一段。还原到原基点，得到
	\[
		\gamma(t)=P^{1/2}\exp\bigl(t\log Q_0\bigr)P^{1/2}.
	\]
	它常速，初始速度在 \(I\) 处的范数为 \(\tfrac12\lVert\log Q_0\rVert_{\mathrm F}\)，因此长度及距离正是题设表达式。若 \(P,Q\) 交换，选取共同正交特征基，分别记其特征值为 \(p_i,q_i>0\)，则
	\[
		d(P,Q)^2=\frac14\sum_{i=1}^n\left(\log\frac{q_i}{p_i}\right)^2.
	\]
	此时测地线一直位于共同对角最大平坦中，其第 \(i\) 个对角元为 \(p_i^{1-t}q_i^t\)。
\end{solution}

\begin{exercise}\label{g3:ex:sl-sectional-sharp-range}
	在 \(\operatorname{SL}(n,\mathbb R)/\operatorname{SO}(n)\) 上采用 \(\mathfrak p=\operatorname{Sym}_0(n,\mathbb R)\) 中的迹内积。证明所有截面曲率均在 \([-2,0]\) 中；当 \(n\ge3\) 时，这个区间的每个值都出现。说明 \(n=2\) 时的情形，并换算到 Killing 型度量。
\end{exercise}
\begin{solution}
	对所求平面取迹内积下的正交单位基 \(X,Y\)，再正交对角化 \(X=\operatorname{diag}(h_i)\)。此时 \(\sum_i h_i^2=1\)，且
	\[
		\lVert[X,Y]\rVert_{\mathrm F}^2
		=\sum_{i,j}(h_i-h_j)^2y_{ij}^2.
	\]
	由 \((h_i-h_j)^2\le2(h_i^2+h_j^2)\le2\) 及 \(\sum_{i,j}y_{ij}^2=1\)，得到 \(\lVert[X,Y]\rVert_{\mathrm F}^2\le2\)。式\eqref{g3:eq:sl-trace-sectional}遂给出 \(-2\le K\le0\)。

	当 \(n\ge3\)，在左上角三阶块中取
	\[
		X=\frac1{\sqrt2}\operatorname{diag}(1,-1,0),\qquad
		Y_0=\frac{E_{12}+E_{21}}{\sqrt2},\qquad
		Y_1=\frac1{\sqrt6}\operatorname{diag}(1,1,-2),
	\]
	其余分量取零。三个矩阵迹为零且两两正交单位，\([X,Y_1]=0\)、\(\lVert[X,Y_0]\rVert_{\mathrm F}^2=2\)。因此 \(Y_t=\cos t\,Y_0+\sin t\,Y_1\) 与 \(X\) 正交单位，并有
	\[
		K(X\wedge Y_t)=-2\cos^2t.
	\]
	令 \(0\le t\le\pi/2\) 即得到整个区间。当 \(n=2\) 时，\(\mathfrak p\) 只有二维，正文给出的基已计算出唯一切平面的曲率 \(-2\)，因此没有零曲率平面。Killing 型度量为迹度量的 \(2n\) 倍，故 \(n\ge3\) 时的区间变成 \([-1/n,0]\)，\(n=2\) 时为常数 \(-1/2\)。
\end{solution}

\begin{optionalexercise}\label{g3:ex:varying-hyperkahler-sphere}
	设连通流形 \(M\) 带超 Kähler 结构 \((g,I,J,K)\)，并设光滑实函数 \(a,b,c\) 满足 \(a^2+b^2+c^2=1\)。令 \(L=aI+bJ+cK\)。证明 \(L\) 是正交几乎复结构，且它平行当且仅当 \(a,b,c\) 都为常数。再在 Euclid \(\mathbb R^4\) 上使用
	\[
		\omega_I=\mathrm dx^{12}+\mathrm dx^{34},\quad
		\omega_J=\mathrm dx^{13}-\mathrm dx^{24},\quad
		\omega_K=\mathrm dx^{14}+\mathrm dx^{23},
	\]
	其中 \(\mathrm dx^{ij}=\mathrm dx^i\wedge\mathrm dx^j\)，取 \(L=\cos x^1 I+\sin x^1 J\)。计算其基本二形式的外微分，判断该结构是否 Kähler。
\end{optionalexercise}
\begin{solution}
	由四元数关系，交叉项相消，故 \(L^2=-(a^2+b^2+c^2)I_{TM}=-I_{TM}\)。三个算子均斜自伴，故 \(L^*=-L\)，从而 \(L^*L=I_{TM}\)。又有
	\[
		\nabla_XL=X(a)I+X(b)J+X(c)K.
	\]
	若实数 \(u,v,w\) 满足 \(uI+vJ+wK=0\)，平方后可得 \(u^2+v^2+w^2=0\)，所以这三个算子线性无关。因此 \(\nabla L=0\) 当且仅当 \(\mathrm da=\mathrm db=\mathrm dc=0\)，在连通 \(M\) 上等价于三者为常数。

	给定的三个常系数二形式通过 \(\omega_A(X,Y)=g(AX,Y)\) 确定满足四元数关系的 \(I,J,K\)。变量组合的基本形式为
	\[
		\omega_L=\cos x^1(\mathrm dx^{12}+\mathrm dx^{34})
		+\sin x^1(\mathrm dx^{13}-\mathrm dx^{24}),
	\]
	因此
	\[
		\mathrm d\omega_L=-\sin x^1\,\mathrm dx^{134}-\cos x^1\,\mathrm dx^{124}.
	\]
	两个系数不同时为零，故它处处不闭，不能成为 Kähler 结构。每点的代数关系与沿流形的平行性在这个例子中分别可检验。
\end{solution}

\begin{optionalexercise}\label{g3:ex:flat-g2-quotient}
	在平坦环面 \(T^7=\mathbb R^7/\mathbb Z^7\) 上取\eqref{g3:eq:standard-g2-three-form}中的常系数三形式。考虑
	\[
		\sigma[x^1,\ldots,x^7]
		=[x^1+\tfrac12,x^2,x^3,-x^4,-x^5,-x^6,-x^7].
	\]
	证明 \(\sigma\) 是无不动点的二阶等距变换，并保持 \(\varphi_0\)。说明商 \(M=T^7/\langle\sigma\rangle\) 是紧完备光滑平坦流形，带平行正三形式。计算它的整体和受限完整群，以及平行一形式空间的维数。
\end{optionalexercise}
\begin{solution}
	变换保持整格，且平方仅把第一坐标加1，所以在环面上平方为恒等。若有不动点，则第一坐标要求 \(1/2\in\mathbb Z\)，矛盾。其微分为
	\[
		A=\operatorname{diag}(1,1,1,-1,-1,-1,-1),
	\]
	故是等距。三形式中每个三元组含零个或两个被反号的指标，因此七项都保持。自由有限群作用给出光滑覆盖，度量和三形式由局部等距图下降；商紧且无边界，故完备，局部平坦性与平行性也随之下降。

	把 \(\sigma\) 提升到 \(\mathbb R^7\)，商的覆盖变换群由整平移和此仿射映射生成。Euclid 平行运输保持常向量，沿商中一条闭路提升后，端点通过某个覆盖变换与起点相对应；回到同一纤维时出现其线性部分的逆。因此所有返回变换属于 \(\{I,A\}\)。整平移给出 \(I\)，连接一点与其 \(\sigma\) 提升像的路径给出 \(A^{-1}=A\)，故
	\[
		\operatorname{Hol}_p(M)=\{I,A\},\qquad
		\operatorname{Hol}_p^0(M)=\{I\}.
	\]
	平行一形式对应于 \(A^*\) 的不变量，恰为 \(\mathrm dx^1,\mathrm dx^2,\mathrm dx^3\) 的线性张成，维数为3。这个非平凡整体完整群包含于 \(G_2\)，但远小于 \(G_2\)。
\end{solution}

\begin{optionalexercise}\label{g3:ex:complex-hyperbolic-restricted-roots}
	设 \(m\ge1\)，取对称空间 \(\operatorname{SU}(1,m)/\operatorname S(\operatorname U(1)\times\operatorname U(m))\)。其切向矩阵为
	\[
		X_z=\begin{pmatrix}0&z^*\\z&0\end{pmatrix},\qquad z\in\mathbb C^m,
	\]
	采用内积 \(\langle X_z,X_w\rangle=\operatorname{Re}(z^*w)\)。证明秩为一。取 \(H=X_{e_1}\)，计算全部限制根与重数，并由此求出曲率算子 \(V\mapsto R(V,H)H\) 在 \(H^\perp\) 上的特征值与重数。特别说明 \(m=1\) 与 \(m\ge2\) 的区别。
\end{optionalexercise}
\begin{solution}
	直接相乘可得
	\[
		[X_z,X_w]=\begin{pmatrix}z^*w-w^*z&0\\0&zw^*-wz^*\end{pmatrix}.
	\]
	若 \(z\ne0\)，先用酉坐标令 \(z=te_1\)、\(t>0\)。右下块为零迫使 \(w_j=0\) 对所有 \(j>1\) 成立，并迫使 \(w_1\) 为实数；因此 \(w\) 是 \(z\) 的实数倍。反向蕴含直接成立，故交换切向子空间至多一维，\(\mathfrak a=\mathbb RH\) 极大，秩为一。

	为算特征空间，先把实 Lie 代数复化为 \(\mathfrak{sl}(m+1,\mathbb C)\)。这一步不改变实特征值对应的重数：\(H\) 在实代数中，故其复化特征空间是原实特征空间的复化。具体地，令 \(J=\operatorname{diag}(-1,1,\ldots,1)\)，反线性对合 \(Z\mapsto-J\overline Z^{\mathsf T}J\) 的不动点是 \(\mathfrak{su}(1,m)\)，每个无迹复矩阵都由这个对合分成 \(Z_1+\mathrm iZ_2\)，\(Z_i\in\mathfrak{su}(1,m)\)。

	\(H\) 在 \(\mathbb C^{m+1}\) 上的特征值为 \(1,-1,0,\ldots,0\)，其中零的重数为 \(m-1\)。伴随作用的特征值是两个这样的数之差。于是 \(2\) 与 \(-2\) 的重数各为 \(1\)；\(1\) 与 \(-1\) 的重数各为 \(2(m-1)\)。令 \(\alpha(H)=1\)，则 \(m\ge2\) 时
	\[
		\Sigma=\{\pm\alpha,\pm2\alpha\},\qquad
		m_\alpha=2(m-1),\quad m_{2\alpha}=1.
	\]
	维数核验为 \(1+2(m-1)+1=2m\)。当 \(m=1\) 时，\(\pm\alpha\) 根空间消失，只剩 \(\pm2\alpha\)，所以根系是约化的 \(A_1\)，不是仍含零重数根的 \(BC_1\)。

	题设 \(H\) 为单位向量。由式\eqref{g3:eq:restricted-root-curvature}，\(H^\perp\) 上的曲率算子有特征值 \(-1\)，重数 \(2m-2\)，以及特征值 \(-4\)，重数 \(1\)。当 \(m=1\)，只剩常曲率 \(-4\) 的二维模型；当 \(m\ge2\)，在两个正交特征空间之间取单位线性组合可得到 \([-4,-1]\) 中每个径向截面曲率。稳定子在单位切向量上等价于酉群的传递作用，故这也是所有截面曲率的范围。这个归一化由题设内积决定，改变内积的常数倍后须相应缩放曲率。
\end{solution}

\begin{optionalexercise}\label{g3:ex:same-rank-dimension-distinct-symmetric}
	比较两个非紧型不可约对称空间
	\[
		M_1=\operatorname{SL}(3,\mathbb C)/\operatorname{SU}(3),\qquad
		M_2=G_{2(2)}/K,
	\]
	其中 \(K\) 的 Lie 代数为 \(\mathfrak{so}(4)\)，模型均取单连通形式。由分类表及限制根计算，证明二者维数与秩相同，却不能通过常数缩放变成等距空间。要求用沿最大平坦的曲率算子核说明区别，而不仅比较名称。
\end{optionalexercise}
\begin{solution}
	\(M_1\) 的限制根系为 \(A_2\)，每个根重数为 \(2\)；它有三个正根，所以秩为 \(2\)，维数为 \(2+3\cdot2=8\)。\(M_2\) 的限制根系为 \(G_2\)，所有根重数为 \(1\)，共有六个正根，故秩也是 \(2\)，维数为 \(2+6=8\)。这些数据分别来自群型表及例外表的 G 行。

	固定最大平坦切空间 \(\mathfrak a\)。对非零 \(H\in\mathfrak a\)，记 \(J_H(V)=R(V,H)H\)。由根空间分解，
	\[
		\ker J_H=\mathfrak a\oplus
		\bigoplus_{\substack{\alpha\in\Sigma^+\\\alpha(H)=0}}\mathfrak p_\alpha,
	\qquad
		\dim\ker J_H=2+\sum_{\substack{\alpha\in\Sigma^+\\\alpha(H)=0}}m_\alpha.
	\]
	一般方向的核维数为 \(2\)。在 \(A_2\) 平面中，有三条不同的根核直线；非零 \(H\) 落在其中一条时，核维数增为 \(4\)。在 \(G_2\) 平面中，则有六条根核直线，非零点上的核维数增为 \(3\)。因为两者都约化，每条直线只对应一对正负根，故这里没有漏计同一直线上的倍根。

	等距变换保持最大平坦及曲率算子核；常数缩放度量也不改变 \(R(V,H)H=0\) 这一条件。因此它们必须保持上述奇异直线与核维数。两模型的直线数及核维数跳跃均不同，故任何常数缩放后也不等距。这说明维数与秩不足以识别不可约对称空间。
\end{solution}
